% concatenated from main.tex, 00_abstract.tex, 01_intro.tex, 02_two-face.tex, 03_general.tex, 04_appendix.tex
\documentclass[11pt]{article}
\usepackage{amsmath,amsfonts,amssymb,amsthm}
\usepackage{fullpage}
\usepackage[ruled]{algorithm2e}
\renewcommand{\algorithmcfname}{ALGORITHM}
\usepackage{subfigure} 
\usepackage{epstopdf}
\usepackage{enumerate}
\usepackage{enumitem}
\usepackage{graphicx}
%\usepackage{cite}
\usepackage{natbib}
\usepackage{color}
\usepackage[sans]{dsfont}
\usepackage{mathtools}
\usepackage{hyperref}
\usepackage{cleveref}
\usepackage{tcolorbox}
\usepackage{multirow}
\usepackage{xspace}
\tcbuselibrary{skins,breakable}
\tcbset{enhanced jigsaw}
%\hypersetup{colorlinks,linkcolor={blue},citecolor={blue},urlcolor={red}}
%\usepackage[style=alphabetic,natbib=true,maxbibnames=99]{biblatex}
%\addbibresource{refs.bib}

\newcommand{\george}[1]{{\color{green} George: #1}}
\newcommand{\zihan}[1]{{\color{blue} Zihan: #1}}
\newcommand{\tianyi}[1]{{\color{purple} Tianyi: #1}}

\newtheorem{theorem}{Theorem}
\newtheorem{observation}[theorem]{Observation}
\newtheorem{lemma}[theorem]{Lemma}
\newtheorem{claim}[theorem]{Claim}
\newtheorem{corollary}[theorem]{Corollary}
\newtheorem{proposition}[theorem]{Proposition}
\newtheorem{definition}[theorem]{Definition}
\newtheorem{assumption}[theorem]{Assumption}
\newtheorem{remark}[theorem]{Remark}

\newcommand{\poly}{\mathsf{poly}}
\newcommand{\aset}{{\mathcal{A}}}
\newcommand{\bset}{{\mathcal{B}}}
\newcommand{\gset}{{\mathcal{G}}}
\newcommand{\hset}{{\mathcal{H}}}
\newcommand{\dset}{{\mathcal{D}}}
\newcommand{\sset}{{\mathcal{S}}}
\newcommand{\dist}{\textnormal{\textsf{dist}}}
\newcommand{\eps}{\varepsilon}

\newcommand{\set}[1]{\left\{ #1 \right\}}

\newenvironment{properties}[2][0]
{\renewcommand{\theenumi}{#2\arabic{enumi}}
	\begin{enumerate} \setcounter{enumi}{#1}}{\end{enumerate}\renewcommand{\theenumi}{\arabic{enumi}}}


\setlength{\parskip}{2mm} \setlength{\parindent}{0mm}

\begin{document}

\begin{titlepage}
	
\title{Paths and Intersections: Exact Emulators for Planar Graphs}
\date{}
%\iffalse{STOC submission is double blind}
\author{George Z. Li\thanks{Carnegie Mellon University, PA, USA. Email: {\tt gzli929@gmail.com}. Part of the work was done at the 2023 DIMACS REU program, supported by NSF grants CCF-1836666 and CNS-2150186.}  
\and 
Zihan Tan\thanks{University of Minnesota Twin Cities, MN, USA. Email: {\tt zihantan1993@gmail.com}.}
\and 
Tianyi Zhang \thanks{ETH Zürich, \href{}{tianyi.zhang@inf.ethz.ch. Supported by the starting grant ``A New Paradigm for Flow and Cut Algorithms'' (no. TMSGI2\_218022) of the Swiss National Science Foundation. Work done at Tel Aviv University, supported by European Research Council (ERC) under the European Union’s Horizon 2020 research and innovation programme (grant agreement No 803118 UncertainENV).}} } 
%\fi

\maketitle
	
\thispagestyle{empty}

\begin{abstract}
    We study vertex sparsification for preserving distances in planar graphs. Given an edge-weighted planar graph with $k$ terminals, the goal is to construct an emulator, which is a smaller edge-weighted planar graph that contains the terminals and exactly preserves the pairwise distances between them. We construct exact planar emulators of size $O(f^2k^2)$ in the setting where terminals lie on $f$ faces in the planar embedding of the input graph. Our result generalizes and interpolates between the previous results of Chang and Ophelders and Goranci, Henzinger, and Peng which is an $O(k^2)$ bound in the setting where all terminals lie on a single face (i.e., $f=1$), and the result of Krauthgamer, Nguyen, and Zondiner, which is an $O(k^4)$ bound for the general case (i.e., $f=k$).

    Our construction follows a recent new way of analyzing graph structures, by viewing graphs as paths and their intersections, which we believe is of independent interest.
\end{abstract}

\end{titlepage}

\tableofcontents

\newpage

\section{Introduction}

Graph compression is an approach that reduces the size of large graphs while preserving key properties, such as flow/cut/distance information. By shrinking the graph before performing computations, it saves computational resources. This approach has played a crucial role in designing faster and more efficient graph algorithms.









We study a specific type of graph compression, called \emph{vertex sparsifiers for preserving distances}, on planar graphs. The input is an edge-weighted planar graph $G$ and a subset $T\subseteq V(G)$ of its vertices called \emph{terminals}. The goal is to compute a smaller planar graph $H$ that contains all terminals in $T$ and preserves their pairwise distances in $G$.
Such a graph $H$ is called an \emph{emulator} of $G$. Formally, $H$ is a \emph{quality-$q$} emulator of $G$ with respect to $T$, for some real number $q\ge 1$, iff
\[
\forall t,t'\in T, \quad \dist_G(t,t')\le \dist_H(t,t')\le q\cdot \dist_G(t,t').
\]
The research focus in vertex sparsification lies in the trade-off between the quality and the emulator size, measured by $|V(H)|$, the number of vertices in the emulator. Ideally, we would like to preserve all distances accurately with small-sized emulators.

If there were no constraints on the structure of the emulator, then for any graph $G$ (not necessarily planar) with $k$ terminals, we could simply set $H$ as the clique on the terminals and give each edge $(t,t')$ length $\dist_G(t,t')$, trivially achieving optimal quality $q=1$ and optimal size $|V(H)|=k$.
Therefore, for distance-preserving vertex sparsification to be interesting and useful, the emulator needs to inherit some structural properties from the input graph $G$. A natural requirement is that $H$ be a \emph{minor} of $G$. Such emulators are called \emph{distance-preserving minors} and have been studied extensively.
In the setting where $q=1$, it was shown by Krauthgamer, Nguyen, and Zondiner \cite{krauthgamer2014preserving} via a simple analysis that every graph $G$ admits a quality-$1$ distance-preserving minor with size $O(k^4)$, and they also proved a lower bound of $\Omega(k^2)$, leaving a gap to be closed. Another interesting setting is when $|V(H)|=k$, and is called the \emph{Steiner Point Removal} problem (first studied by Gupta \cite{gupta2001steiner}), requiring that the emulator only contain terminals.
For general graphs, the best achievable quality is shown to be $O(\log |T|)$ \cite{kamma2015cutting,cheung2018steiner,filtser2018steiner} and $\tilde \Omega(\sqrt{\log |T|})$ \cite{chan2006tight,chen2024lower}.
For minor-free graphs, after an exciting line of work
\cite{basu2008steiner,filtser2020scattering,hershkowitz20211,chang2023covering,chang2023shortcut}, Chang, Conroy, Le, Milenkovic, Solomon and Than recently showed that quality $O(1)$ can be achieved.


Another regime that recently received much attention is planar emulators for planar graphs. The input is a planar graph $G$, and goal is to construct an emulator $H$ that is also a planar graph but not necessarily a minor of $G$.
A recent work by Chang, Krauthgamer, and Tan \cite{chang2022almost} showed that quality $(1+\eps)$ can be achieved by emulators of near-linear size $O(k\cdot\poly(\log k/\eps))$. For exact (quality-$1$) planar emulators, the best upper and lower bounds remain $O(k^4)$ and $\Omega(k^2)$ \cite{krauthgamer2014preserving}.
Recently, Chang and Ophelders~\cite{ChangO20} and Goranci, Henzinger, and Peng~\cite{goranci2020improved} studied the case where all terminals lie on a single face in the planar embedding of $G$, and they showed that such graphs have exact planar emulators of size $O(k^2)$, which is also shown to be tight.

Therefore, the natural next question is:
%\vspace{-1pt}
\[\emph{Does every planar graph with k terminals have an exact planar emulator of size $O(k^2)$?}
\]
%\vspace{-1pt}
















\iffalse
In graph sparsification problems, the goal is to transform an input graph $G$ into a smaller graph $G^\prime$ while preserving, possibly approximately, some properties of the original graph such as cut or distance values. These problems have many natural algorithmic applications. For instance, one can compute the smaller graph $G^\prime$ as a preprocessing step, which can reduce computation for subsequent algorithms. 

Our work considers a subclass of these problems called vertex sparsification, where the goal is to minimize the size of the vertex set of $G^\prime$. In this setting, we are given a terminal set $T$ of $k$ nodes and we wish to output a new graph containing $T$ and a small number of other vertices which preserves some property among the terminals. Specifically, our work considers the case where we wish to preserve the pairwise distances between the terminals exactly. We focus on the case where $G$ is known to be planar, and we restrict $G^\prime$ to also be planar (which excludes the trivial solution of outputting the complete graph on $T$). This restriction has natural applications where there are faster algorithms for shortest path problems on planar graphs. 

Let us briefly overview the known bounds for planar emulators. For general planar graphs, it is known that an emulator of size $O(k^4)$ suffices to preserve all pairwise distances exactly. It is also known that $\Omega(k^2)$ nodes is required, even in the special case where all terminal nodes lie on the same face \cite{ChangO20}. Complementing this result, \citet{ChangO20} showed that when all terminals lie on the same face, $O(k^2)$ nodes actually suffices. Our work generalizes this result: we show that even the terminals lie on two faces, we can construct a planar emulator of size $O(k^2)$. 
\fi







\subsection{Our Result}

In this paper, we make progress towards this question. Our result is summarized in the following theorem.

\begin{theorem}
\label{thm: main}
Let $G$ be an edge-weighted planar graph. Let $T\subseteq V(G)$ be a set of $k$ vertices, such that there exist $f$ faces in the planar embedding of $G$ that contain all vertices of $T$. Then there exists another edge-weighted planar graph $H$ with $T\subseteq V(H)$ and $|V(H)|=O(f^2k^2)$, such that for all  $t,t'\in T$, $\dist_H(t,t')=\dist_G(t,t')$.
\end{theorem}

When $f$ takes its minimum value $1$, our result recovers the $O(k^2)$ bound given by \cite{ChangO20,goranci2020improved}.
When $f$ takes its maximum value $k$, our result recovers the $O(k^4)$ bound given by \cite{krauthgamer2014preserving}.
Therefore, our result can be viewed as interpolating and generalizing these previous results.

Unlike the previous work \cite{ChangO20} which showed a universal construction and gave clean formula for directly computing edge weights from terminal distances, or the previous work \cite{goranci2020improved,krauthgamer2014preserving} which constructed the emulators by performing a series of operations on the input graph, we first construct a skeleton graph that captures the shortest path structures in the input graph, and then \emph{use a linear program} to find the edge weights that preserve the shortest-path distances between terminals.
For the first part, our construction of the skeleton graph relies on a new understanding recently adopted in \cite{chen2025path}: instead of viewing a graph as consisting of vertices and edges, we view it as being formed by paths and their intersections.
For the second part, the crux is establishing the feasibility of the edge-weight linear program. We take the dual of this LP, reducing the feasibility to proving the flow-equivalence properties of certain path-morphing operations, which is then done by Wye-Delta transformations\footnote{The Wye-Delta transformation is a celebrated technique for analyzing electrical networks. This technique has also been used in \cite{goranci2020improved} for simplifying half-grid graphs with terminals on the diagonal. We employ this technique in a different way, to show non-existence of certain flows.}. We believe that our approaches in both parts are of independent interest and will prove useful for more distance-based graph problems.




\paragraph{Related Work.}
Chang, Gawrychowski, Mozes, and Weimann \cite{chang2018near} also studied emulators for unweighted planar graphs, and proved that every $n$-vertex undirected unweighted planar graph admits an exact emulator (not necessarily planar) of size $\tilde O(\min\{k^2,\sqrt{kn}\})$, which is near-optimal.

%In addition to vertex sparsification for preserving distances, recently there has also been much work on vertex sparsification for preserving cut and flow structures of a graph \cite{hagerup1998characterizing,moitra2009approximation,leighton2010extensions,charikar2010vertex,chuzhoy2012vertex,kratsch2012representative,krauthgamer2013mimicking,andoni2014towards,khan2014mimicking,krauthgamer2017refined,karpov2017exponential,chen2020fast,goranci2021expander,krauthgamer2023exact,chen20241+,chen2024cut,das2024nearly}, leading to to many interesting algorithmic implications.

\paragraph{Organization.} 
We provide a technical overview \Cref{sec: warmup}, featuring our skeleton graph construction, using the $2$-face case ($f=2$) as an example, to provide intuitions for the general $f$-face case. The detailed construction for the $f$-face case is then provided in \Cref{sec: main}. In \Cref{sec: edge weight}, we show how to compute the edge weights for the skeleton graph to preserves the shortest-path distances, where in \Cref{subsec: weight overview} contains a high-level overview of the approach in this part.


%\paragraph{Assumptions.} 
%Here, we state some standard assumptions which can be made without loss of generality. Let the nodes in $O$ be labeled $u_0$ through $u_{|O|-1}$ starting at some node and going clockwise and similarly for nodes in $I$ with the labels $v_0$ through $v_{|I|-1}$.  
%Without loss of generality, we may assume that the graph contains edges between consecutive nodes on the inner and outer faces. Furthermore, we may assume these edges are the shortest paths between them by setting their weight to be the same as the shortest path. Observe that this does not change the planarity of the graph or the pairwise distances. Finally, we will assume that all nodes on the inner and outer faces are terminal nodes; this can be done by adding edges between consecutive terminal nodes on each face.
\section{Technical Overview: the 2-Face Case}
\label{sec: warmup}

Our approach consists of two steps: constructing the emulator graph and setting its edge weight.
In this section, we provide an overview of the emulator construction step, illustrating the ideas in the special case where all terminals lie on $2$ faces. An overview of the edge weight setting step can be found in  \Cref{subsec: weight overview}.

The construction for the $O(k^4)$-size emulators \cite{krauthgamer2014preserving} works as follows.
\begin{itemize}
    \item For every pair of terminals $t,t'$, take its shortest path $\pi_{t,t'}$ in $G$ (assuming it is unique, which can be achieved by standard techniques), so we have in total $\binom{k}{2}$ terminal shortest paths.
    \item For every pair of terminal shortest paths $\pi,\pi'$, their intersection must be either empty or a subpath of both $\pi$ and $\pi'$. Call the endpoints of this subpath \emph{special vertices}, so in total we have $2\cdot\binom{\binom{k}{2}}{2}=O(k^4)$ special vertices.
    \item Let $H$ be the union of all terminal shortest paths $\set{\pi_{t,t'}}_{t,t'\in T}$ (so $H\subseteq G$, and clearly $H$ is an exact emulator). The only vertices in $H$ with degree $\ne 2$ are special vertices (which include all terminals). We replace each induced path between a pair of special vertices with an edge connecting them whose length is the total length of the induced path, and the resulting graph $H$ is a minor of $G$ with $O(k^4)$ vertices. 
\end{itemize}

Their approach is simple and elegant, but also conveys an important message: the key information to be stored in an emulator is the \emph{intersections between terminal shortest paths} in $G$.

But how can we reduce the number of such intersections, which can be $O(k^4)$? The previous work \cite{ChangO20}
showed us a way: let different pairs of shortest paths share their intersections. Specifically, they constructed a compact quarter-grid structure that make  single-terminal-source shortest paths overlap significantly, and thereby ensuring that the intersections between them and other shortest paths are reused heavily. Alternatively, their strategy can be also interpreted as: let shortest paths ``support'' each other, in that some shortest paths simply come from the union of two other shortest paths. See \Cref{fig:quartergrid} for an illustration.

\begin{figure}[h]
	\centering
	\includegraphics[scale=0.1]{figs/quartergrid.jpg}
	\caption{An illustration of the quarter-grid construction in \cite{ChangO20}. Terminals (purple) lie on the boundary. Terminal shortest paths are L-shaped. The $1$-$3$ path (red) and the $2$-$5$ path (pink) intersect at $p$. The $1$-$3$ path and the $2$-$4$ path also intersect at $p$, so $p$ as an intersection is shared/reused. Moreover, the $2$-$3$ paths is the $2$-$p$-$3$ L-shaped path, ``supported'' by the pink and the red paths.\label{fig:quartergrid}}
\end{figure}

Our approach also follows this general strategy of decreasing the number of intersections: \emph{extract important shortest paths and let them support other shortest paths}.
But our way of carrying out this strategy is fundamentally different from the previous work. We now discuss them in more detail, using the $2$-face instances as an example.

\subsection*{Critical paths}

Assume the input graph $G$ is a $2$-face instance. That is, all terminals lie on the boundaries of two faces, that we assume without loss of generality to be the outer face and the inner face.

Denote by $t_1,\ldots,t_{I}$ the terminals lying on the outer face and denote by $t'_1,\ldots,t'_{J}$ the terminals lying on the inner face, both in the clockwise order. For each pair $i,j$, we denote by $P_{i,j}$ the shortest path in $G$ connecting $t_i$ to $t'_j$. Consider now a terminal $t_i$ and all shortest paths $P_{i,1},\ldots,P_{i,J}$. 
%As the intersection between every pair of shortest paths is either empty or a single vertex, and we know that all these paths contain vertex $t_i$, they may not share any other vertex
We use the following observation, whose proof is straightforward and is deferred to \Cref{apd: Proof of obs: split}.

\begin{observation}
\label{obs: split}
For each $i$, there is a unique $j$, such that the inner face is contained in the region enclosed by path $P_{i,j}$, path $P_{i,j+1}$, and the segment of the inner boundary  from $t'_j$ clockwise to $t'_{j+1}$. 
\end{observation}

In other words, there exists a unique $j$, such that when we ``morph'' from $P_{i,j}$ to $P_{i,j+1}$, we suddenly ``switch to the other side''.
In this case, we say that terminal $t_i$ \emph{splits} at the terminal pair $(t'_j,t'_{j+1})$.
Intuitively, the paths $P_{i,j}$ and $P_{i,j+1}$ govern the behavior of all other shortest paths from $t_i$, so we call them \emph{critical paths} from $t_i$
(see \Cref{fig:split}).
Critical paths are the important shortest paths we extract from the input graph $G$.

\begin{figure}[h]
	\centering
	\includegraphics[scale=0.1]{figs/split}
	\caption{An illustration of $t_4$ splitting at $(t'_4,t'_5)$. Paths $P_{4,4}$ and $P_{4,5}$ are critical paths from $t_4$.\label{fig:split}}
\end{figure}

Furthermore, we can show that, as we move from terminal to terminal clockwise on the outer face, their ``split location'' also moves clockwise on the inner face. The proof is deferred to \Cref{apd: Proof of obs: split move}.

\begin{observation}
\label{obs: split move}
Let $t_{i_1},t_{i_2},t_{i_3}$ be terminals on the outer face clockwise in this order. If $t_{i_1},t_{i_2},t_{i_3}$ split at $(t'_{j_1},t'_{j_1+1}),(t'_{j_2},t'_{j_2+1}),(t'_{j_3},t'_{j_3+1})$, respectively, then $t'_{j_1},t'_{j_2},t'_{j_3}$ lie on the inner face clockwise in this order.
\end{observation}

\underline{\emph{Simplifying Assumptions:}} In order to better illustrate the ideas in the construction,
we assume that  $k$ is even and $I=J=k/2$. That is, $k/2$ terminals lie on the outer face, and $k/2$ terminals lie on the inner face. Furthermore, we assume that for each $i\in [k/2]$, terminal $t_i$ splits at $(t'_i,t'_{i+1})$, and so paths $P_{i,i}$ and $P_{i,i+1}$ are critical paths from terminal $t_i$.

\subsection*{Emulator structure for preserving inter-face distances\footnote{We remark that, for the purpose of technical overview, we only discuss the emulator structure needed for preserving distances between a terminal on the inner face and a terminal on the outer face, which is the core part of $2$-face emulators. A complete description of the emulator structure is provided in \Cref{apd: 2-face emulator}.}}

Now that we have found the important shortest paths (i.e., the critical paths), it is time for them to form a structure and support other shortest paths. Specifically, the plan is to construct an emulator consisting of only critical paths, and then design the route of other shortest paths accordingly. 

The first question is: \emph{in our emulator $H$, how should these critical paths intersect each other?}

Consider a pair $P_{1,1},P_{2,2}$ of critical paths, where $P_{1,1}$ connects $t_1$ to $t'_1$ and $P_{2,2}$ connects $t_2$ to $t'_2$. Note that, if we let $P_{1,1}$ and $P_{2,2}$ intersect in $H$, then we can extract a $t_1$-to-$t'_2$ path and a $t_2$-to-$t'_1$ path by non-repetitively using the edges of $E(P_{1,1})\sqcup E(P_{2,2})$ (if an edge belongs to both paths then we are allowed to use it twice). Since in graph $H$, $P_{1,1}$ is meant to be the shortest path connecting $t_1$ and $t'_1$ and $P_{2,2}$ is meant to be the shortest path connecting $t_2$ and $t'_2$, by allowing $P_{1,1}$ and $P_{2,2}$ to intersect in $H$, we are essentially enforcing that
\[\dist_H(t_1,t'_2)+\dist_H(t_2,t'_1)\le \dist_H(t_1,t'_1)+\dist_H(t_2,t'_2).\]
However, this is not necessarily true for the corresponding distances in $G$. Therefore, since we want $H$ to be an emulator for $G$, to be on the safe side, we cannot allow paths $P_{1,1}$ and $P_{2,2}$ to intersect.

%Intuitively, for each terminal $t_i$, its critical paths govern the shapes of all other $t_i$-$t'_j$ shortest paths, or in other words, all other $t_i$-$t'_j$ shortest paths can be viewed as intermediate states of morphing from one critical path to the other. In light of this, we need to keep all critical paths in our emulator, and their crossing/non-crossing relations.

By similar reasons, for each $1\le i\le k/2$, we cannot allow path $P_{i,i}$ and path $P_{i+1,i+1}$ to intersect. Therefore, the paths $P_{1,1},P_{2,2},\ldots,P_{k/2,k/2}$ must be ``parallel'' to each other. Similarly, the paths $P_{1,2},P_{2,3},\ldots,P_{k/2,1}$ must be parallel, too.
This leads us to the structure of $H$:
it is simply obtained by (i) adding, for each $i$, two curves from $t_i$ to $t'_{i}$ and $t'_{i+1}$ respectively, that share only the endpoint $t_i$ and are in the same shape as their corresponding critical paths in $G$; (ii) ensuring that curves corresponding to parallel paths do not intersect; (iii) replacing each intersection between two curves with a new vertex. See \Cref{fig: H3} for an illustration.
By reorganizing the terminal locations, we obtain a grid-shape structure shown on the right of \Cref{fig: H3}.
It is easy to verify that $|V(H)|=O(k^2)$.


\begin{figure}[h]
\centering
\subfigure%[Graphs $H_1,H_2,H_3$.]
{\scalebox{0.1}{\includegraphics{figs/H3_1.jpg}}}
\hspace{0.5cm}
\subfigure%[Graph $H$.]
{\scalebox{0.105}{\includegraphics{figs/H3_2.jpg}}}
\caption{An illustration of the structure of $H$. Left: the picture after adding all curves. The curves corresponding to the critical paths of $t_4$ are shown in pink. Right: the grid-shape structure after reorganizing the location of terminals. Each intersection between a horizontal line and a vertical line represents a new vertex. The down edge from $t_7$ connects to the second vertex on the first row (both marked $a$), and similarly for $b,c,d,e,f$. \label{fig: H3}}
\end{figure}

The next question is: \emph{how should we design the route of other shortest paths on this structure?}

Consider the $t_1$-$t'_3$ shortest path. Since its shape is supposed to be governed by critical paths $P_{1,1}$ and $P_{1,2}$, it may not escape the area enclosed by $P_{1,1}, P_{1,2}$ and the inner face boundary. 
Moreover, its intersections with $P_{1,2}$ and $P_{2,3}$ should be subpaths.
So there are two natural ($1$-bend) choices: 
\begin{enumerate}
\item $t_1$ first goes horizontally to $e$ and then vertically down to $t'_3$; and 
\item $t_1$ first goes vertically down to path $P_{2,3}$ and then follows $P_{2,3}$ horizontally to $t'_3$. 
\end{enumerate}
They are in fact symmetric, so let's take option (2).

Surprisingly, this single choice we made for the $t_1$-$t'_3$ shortest path essentially determines the routes for all other shortest paths. To see why it is true, note that
\begin{itemize}
    \item by similar reasons, the shortest paths between $t_1$-$t'_3$, $t_2$-$t'_4$, $t_3$-$t'_5,\ldots,$ $t_{k/2}$-$t'_2$ should be parallel (mutually vertex-disjoint) as well, and the only way to achieve this property in the structure of $H$ is to route, for each $i$, the $t_i$-$t'_{i+2}$ shortest path as: starting at $t_i$, going one step vertically down to the critical path $P_{i+1,i+2}$, and then following $P_{i+1,i+2}$ horizontally to $t'_{i+2}$; 
    \item when we next consider the $t_1$-$t'_4$ shortest path, since its intersections with $t_1$-$t'_2$, $t_1$-$t'_3$ should both be its subpaths, the only natural ($1$-bend) option is to first go $2$ steps vertically down to the critical path $P_{3,4}$ and then follow $P_{3,4}$ horizontally to $t'_4$ (see \Cref{fig:direction}); this in turn forces the similar routes for all $t_2$-$t'_5$, $t_3$-$t'_6,\ldots, t_{k/2}$-$t'_3$ shortest paths: all first $2$ steps down and then all the way to the right; and
    \item finally, via similar arguments, we can argue that all the $t_i$-$t'_j$ shortest paths should take the ``first vertically down, then horizontally right'' $1$-bend route.
\end{itemize}

\begin{figure}[h]
	\centering
	\includegraphics[scale=0.1]{figs/direction.jpg}
	\caption{If the $t_1$-$t'_3$ shortest path follows the ``down-right'' route (shown in green), while the $t_1$-$t'_4$ shortest path follows the ``right-down'' route (shown in black), then this will create a pair of vertices with multiple shortest paths between them, which is undesired.\label{fig:direction}}
\end{figure}

To summarize, by taking option (2) for the $t_1$-$t'_3$ shortest path, we essentially prioritize the vertical direction over the horizontal direction in the structure of $H$ shown in \Cref{fig: H3}, thereby guiding all $t_i$-$t'_j$ shortest paths to go first vertically and then horizontally. Equivalently, we are prioritizing the vertical critical paths over the horizontal ones. 
Looking back, if we had taken option (1) for routing the $t_1$-$t'_3$ shortest path, then we would have eventually prioritized the horizontal critical paths over the vertical ones. In either way, we have to break the ties between the two directions to establish a correct shortest path structure in $H$.
To formalize this, we introduce some new notions.

%As a last step, we need to set the edge weights in $H_3$ to preserve the distances between all interface pairs $t_i,t'_j$. This is done in the next lemma, which, combined with the properties of $H_1$ and $H_2$ and our construction of $H$, immediately implies that $H$ preserves the distances between all pairs of terminals in $G$.

\paragraph{Primary/secondary paths, and canonical paths.}
From a terminal $t_i$, there are two critical paths connecting it to $t'_i$ and $t'_{i+1}$, respectively. We call $P_{i,i}$ the \emph{primary path from $t_i$} and $P_{i,i+1}$ the \emph{secondary path from $t_i$}. Similarly, from each terminal $t'_j$ there are two critical paths connecting it to $t_{j-1}$ and $t_{j}$. We call $P_{j-1,j}$ the \emph{primary path from $t'_j$} and $P_{j,j}$ the \emph{secondary path from $t'_j$}. %In \Cref{fig: H3}, the primary paths with respect to $t_i$ and the secondary paths with respect to $t_j$ are the paths which go vertically, and vice versa for horizontal paths. Given these, w

Between a pair $t_i, t'_j$ of terminals, we define their \emph{canonical path} $\gamma_{i,j}$ as follows. Let $p$ denote the node where the primary paths from $t_i$ and $t'_j$ intersect. The canonical path $\gamma_{i,j}$ is the union of (i) the subpath of the $t_i$-primary path between $t_i$ and $p$; and (ii) the subpath of the $t'_j$-primary path between $p$ and $t'_j$. 
%If there is a critical path between $t_i$ and $t'_j$, then the definition of the canonical path simply reduces to the critical path. %The following observation (whose proof appears in \Cref{apd: Proof of cl: canonical path well defined}) asserts that the canonical paths are well defined.

For every pair $t_i, t'_j$, we will eventually ensure that \emph{canonical path $\gamma_{i,j}$ is the shortest $t_i$-$t'_j$ path in $H$}, whose length is exactly $\dist_G(t_i,t'_j)$, their distance in $G$. 


\subsection*{Summary for the $2$-face emulators and generalization to $f$-face emulators}

At the end of this section, we briefly summarize the ideas for constructing the emulator graph for $2$-face instances. The overall idea is to extract important shortest paths from the input graph and let them support other shortest paths. Concretely, in a $2$-face instance
\begin{itemize}
    \item from each terminal, there are two inter-face terminal shortest paths that govern the shape of others from left and right sides, and we call them \emph{critical paths};
    \item to form a structure using only critical paths, we have to draw all left critical paths in parallel (and they become horizontal lines) and similarly draw all right critical paths in parallel (and they become vertical lines), thereby obtaining a grid-shape structure;
    \item to design routes for other shortest paths in this structure, we have to prioritize one direction (vertical/horizontal) in the grid over the other, thereby classifying the critical paths as primary and secondary ones, and eventually set the routes of other shortest paths as the concatenation of primary subpaths.
\end{itemize}

Since the whole structure consists of $O(k)$ critical paths, and every pair of critical paths intersect at most once, the total number of vertices in the graph is $O(k^2)$.

Our construction for $f$-face instances follows the same strategy as the $2$-face case: extract critical paths, draw them in parallel, let them intersect and form a structure, prioritize directions and guide the routes of other shortest paths. The main difference is that, from a terminal, we will need more than $2$ critical paths to govern the shapes of others. 
Specifically, we will have $O(f)$ critical paths from each terminal (see \Cref{fig: critical_example}), therefore having a total of $O(fk)$ critical paths. Their pairwise intersections form a structure on $O(f^2k^2)$ nodes, which will be our emulator structure.





\begin{figure}[h]
\centering
\includegraphics[scale=0.12]{figs/critical_example.jpg}
\caption{An illustration of $O(f)$ critical paths from a terminal in a $f$-face instance.\label{fig: critical_example}}
\end{figure}
























\iffalse
\begin{observation}\label{cl: canonical path well defined}
    For any two terminals $t_i$ and $t'_j$, the canonical path between $t_i$ and $t'_j$ uniquely exists.
\end{observation}

Furthermore, we have the property (whose proof appears in \Cref{apd: proof of 2-face cross-iff-cross property}) that two canonical paths intersect in our constructed graph iff their corresponding real shortest paths cross in the original graph.

\begin{observation}\label{obs: 2-face cross-iff-cross property}
    Any two canonical paths in $H_3$ cross iff their corresponding shortest paths cross in $G$.
\end{observation}
\fi



\iffalse
\subsection{Defining the Emulator}

Given a planar graph with terminal nodes lying on $2$ faces of the graph, we will show how to obtain a planar emulator of size $O(k^2)$. We first give an emulator which preserves shortest path distances between pairs of nodes on different faces. In this emulator, we will show that the shortest path distances between pairs of nodes on the same face is always (weakly) greater than their actual shortest path distance. As a result, we will be able to augment the emulator by adding a $1$-face emulator on the ``inside'' of each face; the resulting emulator will then preserve shortest path distances between all pairs of nodes. 

\paragraph{Structure of the Universal Emulator.} Recall that we have assumed without loss of generality that all terminals lie on either the outer or inner face. To define the emulator, consider all shortest paths between a fixed node $u_i$ on the outer face and all nodes $v_j$ on the inner face. Informally, these shortest paths either go to the left of the inner face or the right of the inner face (left and right here are very vaguely defined). At some consecutive pair of the nodes $v_j$ and $v_{j+1}$, there is a transition: $u\to v_j$ will be on the left and $u\to v_{j+1}$ will be on the right. We call this pair of paths a pair of critical path. Intuitively, these paths are critical because the shortest paths between $u_i$ and all nodes $v_j$ on the inner face lie ``between'' the two critical paths. Formally, we can define critical paths as follows:

\begin{definition}
    For a node $u_i$ on the outer face, we say that $u_i$ splits at the pair of nodes $v_j,v_{j+1}$ on the inner face if the cycle formed by composing the shortest paths $u_i$ to $v_j$, $v_j$ to $v_{j+1}$, and $v_{j+1}$ to $u_i$ contains the inner face. If $u_i$ splits at $v_j,v_{j+1}$, we say that the shortest paths $u_i\to v_j$ and $u_i\to v_{j+1}$ are critical paths between $u_i$ and the inner face.
\end{definition}

With our definition of critical paths, we can now define the emulator. For each node $u_i$ on the outer face, draw the two critical paths from $u_i$ to the inner face and let all of the critical paths intersect naturally such that no two pairs of paths intersect at the same point. \textcolor{red}{We may need to also have to draw the critical paths from the inner face to the outer face when the paths don't split nicely.} By adding nodes at every point where two critical paths intersect, we can make the resulting graph planar. Furthermore, since there are $O(k)$ critical paths coming out of each face, there are at most $O(k^2)$ intersections so our planar emulator indeed has size $O(k^2)$. It now remains to set the edge weights of the resulting emulator and prove that the shortest paths are preserved.


\subsection{Setting the Edge Weights}

\subsection{Augmenting the Emulator}
\fi
\section{The $f$-face Case: Constructing the Graph}
\label{sec: main}

In this section we provide the proof of \Cref{thm: main}. Recall that we are given a plane graph $G$ and a set $T$ of its vertices called terminals, such that all terminals lie on the boundaries of $f$ faces in its planar embedding, and the goal is to construct a planar emulator $H$ with $|V(H)|=O(f^2k^2)$. 
Similar to the way we handle the $2$-face case in \Cref{sec: warmup}, we will first show the graph structure and then assign weights to its edges.


\subsection{Simplifying the input graph}

In this subsection, we massage the input graph $G$ so that it has good properties that simplifies our construction of the emulator. Let $T$ be the set of terminals, and let $F_1,\ldots,F_f$ be the faces in the planar embedding of $G$ that contain all terminals, then we would like to ensure that

\begin{properties}{G}
\item the boundary of each face is a simple cycle that only contains terminals;\label{prop: G1}
\item the face cycles are vertex-disjoint, so each terminal lies on only one face among $F_1,\ldots,F_f$; and \label{prop: G2}
\item for distinct terminals $t_1,t_2,t_3,t_4$, if we denote by $P$ the $t_1$-$t_2$ shortest path, and denote by $Q$ the $t_3$-$t_4$ shortest path, then either
\begin{itemize}
    \item $P,Q$ are vertex-disjoint; or
    \item $P,Q$ intersect at a single internal vertex, and they cross over this intersection (that is, if we denote by $e_1,e_2$ the edges of $P$ and by $e_3,e_4$ the edges of $Q$ incident to the intersection, then the circular order of the edges $e_1,e_2,e_3,e_4$ entering this vertex is either $e_1,e_3,e_2,e_4$ or $e_1,e_4,e_2,e_3$).
\end{itemize}\label{prop: G3}
\end{properties}

For a face $F$, label the terminals on $F$ in a clockwise order as $t_1,\ldots,t_{r}$. To guarantee Property~\ref{prop: G1}, we simply add edges $(t_i,t_{i+1})$ with weight $\dist_G(t_i,t_{i+1})$ for each $1\le i\le r$ (using the convention $r+1=1$) inside face $F$, so now terminals $t_1,\ldots,t_{r}$ lie on a new face surrounded by the added edges. Clearly, the new face satisfies Property~\ref{prop: G1}, the boundary of the new face contains the same terminals as $F$, and the distance between any pair of terminals remains unchanged. Perform similar operations on other faces and eventually $G$ satisfies Property~\ref{prop: G1}.

Let $F_1,F_2$ be two faces and let $T_{1,2}$ denote the set of terminals lying on both $F_1$ and $F_2$. In order to guarantee Property~\ref{prop: G2}, we simply think of $T_{1,2}$ as terminals on $F_1$ but not on $F_2$. That is, when we add edges to form a new face (in the previous paragraph) for $F_2$, we simply ignore terminals in $T_{1,2}$. It is easy to verify that the new face to replace $F_2$ does not contain terminals in $T_{1,2}$. Perform similar operations on other pairs of faces and eventually $G$ satisfies Property~\ref{prop: G2}.

For achieving Property~\ref{prop: G3}, first we can ensure by standard techniques that between every pair of its vertices there is a unique shortest path connecting them in $G$ (for example, the lexicographic perturbation scheme in \cite{erickson2018holiest}).
As a corollary, the intersection between every pair $P,P'$ of shortest paths in $G$ is either empty or a subpath of both $P$ and $P'$.

In order to achieve Property~\ref{prop: G3}, we ``redraw'' the graph as follows. First, we remove from $G$ all edges from that does not participate in any shortest path connecting a pair of terminals. Then,
\begin{enumerate}
\item for each vertex $u$, create a tiny disc $\dset_u$ around $u$ (so all discs $\set{\dset_u}_{u\in V(G)}$ are disjoint);
\item for each edge $e=(u,v)$, create a thin strip $\sset_e$ along the image of $e$, with one end touching disc $\dset_u$ and the other end touching disc $\dset_v$, such that all strips $\set{\sset_e}_{e\in E(G)}$ are disjoint;
\item for each terminal shortest path $P=(u_0,u_1,\ldots,u_r)$, we define its \emph{channel} as the union of:\\ $\dset_{u_0},\sset_{(u_0,u_1)}, \dset_{u_1},\sset_{(u_1,u_2)},\ldots,\dset_{u_r}$;
\item for each terminal shortest path $P$ connecting $t$ to $t'$, we redraw it as a simple curve $\gamma_P$ from $t$ to $t'$ that lies entirely in its channel, such that there are no crossings in any strips, and all crossings happen within some disc;
\item repeatedly, for a pair $P,Q$ whose images $\gamma_P,\gamma_Q$ cross at least twice, find a pair $c,c'$ of crossings between them that appear consecutively, and uncross $\gamma_P,\gamma_Q$ at their segment between $c,c'$ (decreasing the total number of curve crossings by $2$), until there are no such pairs; and
\item view the curves as the drawing of the new graph, by placing a new vertex on each crossing, and viewing each segment between a pair of consecutive crossings as an edge;
\item for each new edge $(u',v')$, if its image passes through strips $e_1,e_2,\ldots,e_r$, then its new weight is defined to be $w'(u',v')=\sum_{1\le i\le r}w(e_i)$; if its image lies entirely in some disc, then its new weight $w'(u',v')=0$.
\end{enumerate}
 
\begin{figure}[h]
\centering
\subfigure
{\scalebox{0.07}{\includegraphics{figs/disengage_1.jpg}}}
\hspace{1.0cm}
\subfigure
{\scalebox{0.07}{\includegraphics{figs/disengage_2.jpg}}}
\caption{An illustration of replacing shortest paths with curves. The red vertex in the original graph (left) supports shortest paths $1$-$3$, $1$-$5$, $2$-$4$, $2$-$5$, $3$-$4$, $3$-$5$, $5$-$6$. The curves in its disc and its incident thin strips are shown on the right, where new vertices (purple) are placed on the created intersections. \label{fig: YD}}
\end{figure}    

See \Cref{fig: YD} for an illustration. It is easy to verify that the above process produces a new graph satisfying Property~\ref{prop: G3}, does not change the faces and the distance between any pair of terminals.



\subsection{The graph structure of the emulator}

Denote by $F_1,\ldots,F_f$ the faces in the plane embedding of $G$ where all terminals lie on. For each $1\le r\le f$, let $T_r$ be the set of terminals lying on face $F_r$.
%
We will construct several graphs and then stick them together to preserve terminal distances in $G$. Specifically, we construct 
\begin{itemize}
\item for each $1\le r\le f$, a one-face emulator $H_r$ with terminals in $T_r$ lying on a single face $F'_r$ in the same order as they lie on $F_r$, such that for every pair $t,t'\in T_r$, $\dist_{H_r}(t,t')=\dist_{G}(t,t')$; 
and in fact we will simply use the construction of \cite{ChangO20} and \cite{goranci2020improved}, which gives a construction of $H_r$ on $O(|T_r|^2)$ vertices;
and
\item a central graph $H^*$ that contains all terminals, where for each $1\le r\le f$, terminals in $T_r$ lie on a single face $F^*_r$ in the same order in which they lie on $F_r$ in $G$, and we will ensure that 
\begin{itemize}
    \item for each pair $t,t'$ of terminals lying on different faces, $\dist_{H^*}(t,t')=\dist_G(t,t')$;
    \item for each pair $t,t'$ of terminals lying on the same face, $\dist_{H^*}(t,t')\ge \dist_G(t,t')$.
\end{itemize}
\end{itemize}
The graphs $H^*,H_1,\ldots,H_f$ are then glued together in a similar way as in \Cref{apd: 2-face emulator}, by placing the drawings of $H_1,\ldots,H_f$ into the faces $F^*_1,\ldots,F^*_r$ of $H^*$, respectively, and identifying the copies of the each terminal. The correctness of the emulator follows immediately from the above properties.
Therefore, the only missing part is the construction of the central graph $H^*$.

\subsection{Constructing the graph $H^*$: preparation}
\label{sec: prepare}

The construction of graph $H^*$ is a generalization of  the $2$-face emulator construction in \Cref{sec: warmup}. 
The high-level idea for the $2$-face case is to keep, for each terminal $t$ on the outer face, the two critical paths, one named primary and the other named secondary, that govern the shapes of all shortest paths connecting $t$ to some terminal on the inner face, and let them cross in the same way as they cross in $G$. 
Moreover, we define canonical paths connecting every pair of terminals lying on different faces as the concatenation of subpaths of their primary paths,
and designate them as the shortest paths in the emulator. By doing this, we are able to ensure that two canonical paths cross iff their corresponding real shortest paths in $G$ cross.

%However, in order to convey intuitions, we made several simplifying assumptions in \Cref{sec: warmup}. 
We will construct $H^*$ using a similar approach. We first discuss on a high level (i) how to get rid of the simplifying assumptions made in \Cref{sec: warmup}; and (ii) what additional problems we will encounter for general $f$-face instances and how we plan to handle them.

The two assumptions we made in \Cref{sec: warmup} were: the numbers of terminals lying on the inner and the outer faces are the same, and the split location moves ``smoothly''. 
The advantage for these assumptions is that we only need to consider critical paths from terminals lying on the outer face, as they ``automatically include'' the critical paths from terminals on the inner face. Specifically, we could have defined critical paths from terminals on the inner face in the same way as the terminals on the outer face and added them into the structure of $H$, but in fact they were already there. Specifically, in \Cref{fig:direction}, for terminal $t'_1$: the critical path from $t_7$ to $t'_1$ and the critical path from $t_1$ to $t'_1$ are exactly the critical paths we wanted to add from $t'_1$; in other words, we could also say that $t'_1$ splits at $(t_7,t_1)$. Similarly, the critical paths from other terminal $t'_i$ were also there already.
%
However, this is no longer true if we remove the two simplifying assumptions. For example, many outer face terminals may split at the same location, leaving some inner face terminals not serving as the endpoint of any critical path, while others serving as many. Therefore, we do need to include critical paths from inner face terminals as well.

For $f$-face instances, when we look at all shortest paths from a terminal $t$ to other terminals, we will find that their shapes change more frequently, will need more than two of them to govern the rest. Luckily, we are able to show that there always exist $O(f)$ paths among them to govern the rest. This increases the number of total critical paths we need to add to $O(fk)$, which then leads to the $O(f^2k^2)$ bound of the size of $H^*$ (see \Cref{lem: critical} and \Cref{clm: size}). 

Another problem we encounter in handling $f$-face instances is that it becomes unclear what ``cross in the same way as they cross in $G$'' means. In \Cref{sec: warmup}, this simply meant that a pair of critical paths cross in $H$ iff they cross in $G$. However, for $f$-face instances, there can be more than three critical paths crossing each other, so the orders in which the crossings appear on these paths are crucial in the structure of $H^*$, and need to be handled with care. In fact, we cannot simply let the crossings appear in the same order as in $G$, but need to delicately move them, such that when we define canonical paths as the concatenation of certain critical subpaths, the ``canonical paths cross iff their corresponding shortest paths cross in $G$'' property is achieved.
%Our solution to this problems is still: the ordering of these critical paths intersecting each other should also follow the ordering they cross in $G$. 
All these need to be formalized, which we do now.

\paragraph{Critical Paths.}
For each $1\le r\le f$, denote $T_r=\set{t^r_1,\ldots,t^r_{|T_r|}}$, where the terminals are indexed according to the clockwise order in which they appear on the boundary of $F_r$. 
%For each $1\le j\le |T_r|$, denote by $a^r_j$ the boundary segment of face $F_r$ clockwise from $t^r_j$ to $t^r_{j+1}$ (here we use the convention that $|T_r|+1=1$). 
Let $t$ be a terminal not in $T_r$. 
We say that the $t$-$t^r_j$ shortest path in $G$ is \emph{equivalent to} the $t$-$t^r_{j+1}$ shortest path in $G$, iff the region enclosed by the union of (i) the $t$-$t^r_j$ shortest path; (ii) the $t$-$t^r_{j+1}$ shortest path; and (iii) boundary segment of face $F_r$ clockwise from $t^r_j$ to $t^r_{j+1}$, contains none of the faces in $\set{F_1,\ldots,F_f}$.
%
We say that the $t$-$t^r_j$ shortest path in $G$ is \emph{critical}, iff the $t$-$t^r_j$ shortest path in $G$ is not equivalent to
either the $t$-$t^r_{j+1}$ shortest path or the $t$-$t^r_{j-1}$ shortest path.

Intuitively, the $t$-$t^r_j$ shortest path is critical, iff the curve morphing from the $t$-$t^r_{j-1}$ shortest path to the $t$-$t^r_{j}$ shortest path and then to the $t$-$t^r_{j+1}$ shortest path has to ``scan over'' some other face and therefore incur a fundamental shape change. See \Cref{fig:critical} for an illustration.

\begin{figure}[h]
\centering
\includegraphics[scale=0.12]{figs/critical.jpg}
\caption{All critical paths (green) from terminal $t$ and two non-critical shortest paths (pink) from $t$ with their enclosed region with the boundary segment (which contains no faces from $\set{F_1,\ldots,F_f}$).\label{fig:critical}}
\end{figure}

We have the following lemma.

\begin{lemma}
\label{lem: critical}
For each terminal $t$, the number of critical paths from $t$ is $O(f)$.
\end{lemma}
\begin{proof}
Let terminal $t$ lie on $F_1$ without loss of generality. Consider the pairs $(P_1,P_2)$ of critical paths with an endpoint at $t$ such that $P_1$, $P_2$, and an edge enclose a region containing some face.
Define the set system $\mathcal{A}$ to be the collection of regions induced by such pairs $(P_1,P_2)$ with one endpoint at terminal $t$. We claim that $\mathcal{A}$ is a laminar family. Indeed, the boundary of two regions for different critical paths can never cross since two critical paths starting at $t$ can only have one intersection (at $t$), and no other intersections.
Equivalently, the set system $\mathcal{B}$ consisting of the set of faces in the regions enclosed induced by a pair of critical paths $(P_1,P_2)$ from $t$ must also be a laminar family. Since the universe of $\mathcal{B}$ is of size $f$, we have $|\mathcal{B}|\le 2f-1$. 
\begin{observation}
    Each set of faces $B\in\mathcal{B}$ corresponds to at most two different regions $A_1,A_2\in\mathcal{A}$.    
\end{observation}
\begin{proof}
    Let us say $A_1$ is a Type 1 region if $A_1$ has terminal endpoints $(t,t_1,t_2)$ where $t_1,t_2\in B$ and $A_2$ is a Type 2 region if $A_2$ has terminal endpoints $(t,t_1',t_2')$ where $t_1',t_2'\not\in B$. We will show that there is at most one region of each type (for each $B\in\mathcal{B}$).
    
    Suppose there is another Type 1 region $A\neq A_1$ for $B$, defined by two critical paths $P_1,P_2$ from $t$ to consecutive terminals on some face in $B$. Since both endpoints of $P_1$ (resp. $P_2$) are in $A_1$ and $P_1$ (resp. $P_2$) can't cross the boundary of $A_1$ twice, we have $P_1$ (resp. $P_2$) lies entirely in $A_1$ except for the endpoint $t$. But then $A$ doesn't contain $t_1$ or $t_2$, since that would mean $P_1$ or $P_2$ would touch the boundary of $A_1$, a contradiction. 
    
    Suppose there is another Type 2 region $A\neq A_2$ for $B$, defined by two critical paths $P_1,P_2$ from $t$ to consecutive terminals on some face not in $B$. Since $A$ doesn't contain any other faces, we have that $t_1',t_2'\not\in A$. But by laminarity of $\mathcal{A}$, this implies that $A\subseteq A_2$ since $A\cap A_2\neq\emptyset$ and $t_1'\not\in A$ so $A_2\not\subseteq A$. But there are no faces in $A_2$ not in $B$ (by definition of $B$), so the endpoints of $P_1,P_2$ must be $t_1',t_2'$ and $A=A_2$.
\end{proof}
Combining $|\mathcal{B}|\le 2f-1$ with the above observation, we have $|\mathcal{A}|\le 4f-2$. Since each region in $\mathcal{A}$ corresponds to at most $2$ critical paths, there are at most $8f-4$ critical paths starting at $t$.
\end{proof}

\paragraph{Primary and secondary paths, and canonical paths.}
The critical paths can be naturally paired according to their shapes.
Let $t$ be a terminal and consider a face $F_r$ such that $t\notin F_r$. For terminals $t^r_1,\ldots,t^r_{|T_r|}$ that lie on $F_r$, according to the equivalence relation defined on the shortest paths from $t$ to them in $G$, the circular ordering $(1,\ldots,|T_r|,1)$ can be partitioned into segments: 
\[[a_0+1,\ldots,a_1],[a_1+1,\ldots,a_2],\ldots,[a_{s}+1,a_0],\]
where for each $0\le i\le s$, the $t$-$t^r_{a_i+1}$ path, the $t$-$t^r_{a_i+2}$ path, ..., and the $t$-$t^r_{a_{i+1}}$ path are equivalent to each other, but not to the $t$-$t^r_{a_i}$ path or $t$-$t^r_{a_{i+1}+1}$ path.
By definition, the $t$-$t^r_{a_i+1}$ path and the $t$-$t^r_{a_{i+1}}$ path are critical paths, and we say that they are \emph{paired}. 
We call the boundary segment of $F_r$ going clockwise from $t^r_{a_i+1}$ to $t^r_{a_{i+1}}$ the \emph{governed segment} of these paired shortest paths.
%We say that both critical paths $t$-$t^r_{a_i+1}$ and $t$-$t^r_{a_{i+1}}$ govern all terminals lying on this governed segment.

For paired critical paths, we define one of them as \emph{primary}, and the other as \emph{secondary}, as follows. Consider paired critical paths $t$-$t^r_{a}$ and $t$-$t^r_{b}$, where $t^r_a$ and $t^r_b$ lie on face $F_r$ and $t$ lies on face $F_{r^*}$, and the $F_r$ boundary segment between $t^r_{a}$ and $t^r_{b}$ is going clockwise from $t^r_{a}$ to $t^r_{b}$. Now
\begin{itemize}
\item if $r^*<r$, then the $t$-$t^r_{a}$ path is primary, and the $t$-$t^r_{b}$ path is secondary; and
\item if $r^*>r$, then the $t$-$t^r_{b}$ path is primary, and the $t$-$t^r_{a}$ path is secondary.
\end{itemize}

Consider now a pair $t,t'$ of terminals lying on distinct faces $F_r,F_{r'}$, respectively, with $r<r'$. The \emph{canonical path} between terminals $t,t'$ are defined as follows.
Let $t'$-$t^r_{a}$, $t'$-$t^r_b$ be the paired critical paths whose governed segment $[t^r_{a},t^r_b]$ on face $F_r$ contains $t$, and let $t$-$t^{r'}_{a}$, $t$-$t^{r'}_b$ be the paired critical paths whose governed segment $[t^{r'}_{a},t^{r'}_b]$ on face $F_{r'}$ contains $t'$. So the $t$-$t^{r'}_{a}$ shortest path and the $t'$-$t^{r}_{b}$ shortest path must cross, and we denote their crossing by $p$.
The $t$-$t'$ canonical path is defined to be the concatenation of 
\begin{itemize}
\item the subpath of the $t$-$t^{r'}_{a}$ shortest path between $t$ and $p$; and 
\item the subpath of the $t'$-$t^{r}_{b}$ shortest path between $t'$ and $p$.
\end{itemize}
Additionally, we call the intersection $p$ the \emph{bend} of the canonical path $t$-$t'$. 
%Before moving on, we show a property of canonical paths between two faces which we will need later on.


\subsection{Constructing the graph $H^*$: the algorithm}

Similar to the construction of the structure of $H$ in \Cref{sec: warmup}, we can construct $H^*$ by (i) starting with the subgraph of $G$ induced by all critical paths from all terminals; (ii) suppressing all degree-$2$ vertices; and (iii) for every pair $t,t'$ of terminals, define its canonical path $\gamma_{t,t'}$ as in \Cref{sec: prepare}.
%It is easy to observe that all canonical paths are kept in this structure. 
However, this construction does not achieve the crucial property that a pair of canonical shortest paths cross iff their corresponding shortest paths in $G$ cross (for example, see \Cref{fig: misplace_1}).
Therefore, we need to modify the structure by an iterative process to achieve this.

\begin{figure}[h]
\centering
\includegraphics[scale=0.12]{figs/misplace_1.jpg}
\caption{An illustration of multiple crossings between the $t$-$t'$ critical path (the green dashed line between $t,t'$) and the $t_1$-$t'_1$ canonical path (red), and some other relevant critical paths.\label{fig: misplace_1}}
\end{figure}

Throughout, we maintain a graph $\hat H$, that is initialized as the structure obtained in the previous paragraph.
We also maintain a drawing of $\hat H$. Initially, the drawing of $\hat H$ is induced by the drawing of $G$ as $\hat H$ comes from $G$.
We will ensure that graph $\hat H$ always satisfies the following properties.


\begin{enumerate}
\item Graph $\hat H$ is an $f$-face instance that contains all terminals; for each $1\le r\le f$, the terminals in $T_r$ lie on $F_r$ in the same order in which they appear on $F_r$ in $G$.
\label{prop: f-face instance}
\item Graph $\hat H$ is the union of, for each critical path $P$ in $G$, a path $\rho_P$ connecting the same pair of terminals as $P$. For a pair $P,P'$ of critical paths whose endpoints are four distinct terminals, \label{prop: canonical intersection}
\begin{itemize}
    \item if $P,P'$ are vertex-disjoint, then $\rho_P,\rho_{P'}$ are also vertex-disjoint; and
    \item if $P,P'$ cross (at a single vertex), then $\rho_P,\rho_{P'}$ cross (also at a single vertex).
\end{itemize}
%\item For every pair $t_1,t_2$, the drawings of the $t_1$-$t_2$ canonical path in $\hat{H}$ and the $t_1$-$t_2$ shortest path in $G$ don't cross except at the endpoints.
%\label{prop: canonical-shortest-no-cross}
\item For every pair $t_1,t_2$ of terminals on different faces, the region enclosed by the $t_1$-$t_2$ canonical path in $\hat H$ and the $t_1$-$t_2$ shortest path in $G$ contains no face and no other terminals.
\label{prop: region}
\end{enumerate}

Clearly, Properties \ref{prop: f-face instance} and \ref{prop: canonical intersection} are satisfied by the initial $\hat H$. We show that Property \ref{prop: region} is also satisfied.

\begin{observation}
The initial graph $\hat H$ satisfies Property \ref{prop: region}.
\end{observation}
\begin{proof}
    Observe that by our construction of the critical and canonical paths, both the $t_1$-$t_2$ canonical path and the $t_1$-$t_2$ shortest path lie in the region enclosed by the critical paths outgoing from $t_1$ to the face $F_{t_2}$ where $t_2$ lies. Symmetrically, the two paths also lie in the region enclosed by the critical paths outgoing from $t_2$ to the face $F_{t_1}$ where $t_1$ lies. Call these regions $A_1$ and $A_2$, respectively. Observe that the only face which lies in region $A_1$ is $F_{t_2}$ and symmetrically, the only face which lies in region $A_2$ is $F_{t_1}$. Since the region enclosed by the $t_1$-$t_2$ canonical path and the $t_1$-$t_2$ shortest path lies in the intersection of $A_1$ and $A_2$, there must be no faces contained in it. Similarly, the only terminals which lie in region $A_1$ are $t_1$ and the terminals on $F_{t_2}$ and symmetrically, the only terminals which lie in $A_2$ are $t_2$ and the terminals on $F_{t_1}$. Since the region enclosed by the $t_1$-$t_2$ canonical path and the $t_1$-$t_2$ shortest path lies in the intersection of $A_1$ and $A_2$, the only terminals contained in the region are $t_1$ and $t_2$.
\end{proof}

Eventually, at the end of the iterative process, we will ensure that the resulting graph $\hat H$ satisfies the following additional property.

\begin{enumerate}[resume]
\item For a canonical path $\gamma_{t_1,t_2}$ between $t_1,t_2$ and another canonical path $\gamma_{t_3,t_4}$ between $t_3,t_4$ with $t_1,t_2,t_3,t_4$ being distinct terminals, $\gamma_{t_1,t_2}$ and $\gamma_{t_3,t_4}$ cross iff the $t_1$-$t_2$ shortest path in $G$ and the $t_3$-$t_4$ shortest path in $G$ cross.
\label{prop: intersect iff}
\end{enumerate}

%Throughout the process, we will also not create or remove any intersections between critical paths, but will only modify the orderings of intersections on critical paths, and along with it the drawing of $\hat H$ in the plane.

%\paragraph{Misplaced critical paths.}
%Let $t,t'$ be a pair of terminals such that there exists a critical path from $t$ to $t'$. We say that such a critical path is \emph{misplaced}, iff there exist a terminal $t_1$ lying on the same face with $t$ and right next to $t$ (either clockwise or anti-clockwise), and a terminal $t'_1$ lying on the same face with $t'$, right next to $t'$, and on the same side of the $t$-$t'$ critical path as $t_1$, such that the $t_1$-$t'_1$ canonical path intersects the $t$-$t'$ critical path (twice). 

%\paragraph{Properties of canonical paths.}



%\paragraph{Bad pairs of canonical paths.}
%We say a critical path $P$ and a non-critical canonical path $Q$ are a \emph{bad pair} if they intersect twice in $H$, whereas their corresponding shortest paths do not intersect in $G$.

%In the following observations, we will establish that these pairs are the only problems in the current definition of the emulator.

%We say that a pair $P,Q$ of canonical paths are a \emph{bad pair} if they intersect twice in $H$, whereas their corresponding shortest paths do not intersect in $G$. %In the following claims, we will show that bad pairs are the only pairs of canonical paths not satisfying Property \ref{prop: intersect iff}. 



%From these claims, we see that the only problems are canonical paths which intersect twice even though their corresponding shortest paths don't intersect; this is exactly the definition of a bad pair. Next, we will show that the root cause of the problem is that some critical paths intersect a canonical path twice (exactly the example illustrated in Figure \ref{fig: misplace_1}). 




% \begin{claim}
%     If $a$ is the number of intersections between two canonical paths in $\hat{H}$ and $b$ is the number of intersections between the corresponding shortest paths in $G$, we have $a\equiv b\text{ (mod 2)}$. \label{cl:modulo 2}
% \end{claim}
% \begin{proof}
%      Let $P_{1,2}$ and $P_{3,4}$ denote the shortest paths in $G$, let $Q_{1,2}$ and $Q_{3,4}$ denote the corresponding canonical paths in $H$ between the same terminals, and let us place these paths on the same drawing. Let $A$ denote the area inscribed by $P_{1,2}$ and $Q_{1,2}$ and $B$ denote the area inscribed by $P_{3,4}$ and $Q_{3,4}$. Finally, let $c$ denote the number of times $P_{3,4}$ intersects $Q_{1,2}$. 
     
%      By Property 3, we know that region $A$ or region $B$ doesn't include any face or any other terminals. Since the endpoints of $P_{3,4}$ are not in region $A$, path $P_{3,4}$ must cross the boundary an even number of times, so $b\equiv c\text{ (mod 2)}$. Since the endpoints of $Q_{1,2}$ are not in region $B$, path $Q_{1,2}$ must cross the boundary an even number of times, so $a\equiv c\text{ (mod 2)}$. Combining the two results gives $a\equiv b\text{ (mod 2)}$, as desired.
% \end{proof}



%Furthermore, we see that since pairs of critical paths and pairs of non-critical canonical paths can only intersect 0 or 1 times, the only possible problem is the intersection between critical and non-critical canonical paths. Given this, we see that the bad pairs defined previously are the only problem, and it suffices to fix them. In order to do so, we will define a process where at least one bad pair is fixed at each iteration. Furthermore, we will show that no new intersections are caused by the process, so no further problems are caused.



%\paragraph{Description of an iteration.}
%We now describe an iteration. If there are no misplaced critical paths in the current graph $\hat H$, we terminate the process and return it as $H^*$. Otherwise, pick any misplaced critical path. Assume that it connects $t$ to $t'$, and we let $t_1,t'_1$ be the terminals satisfying the definition above. Let $t_1$-$x$ path and $t'_1$-$y$ path be the primary paths from $t_1$ and $t'_1$ that generate the $t_1$-$t'_1$ canonical path, and denote by $p,q$ the intersections between $t$-$t'$ critical path and the $t_1$-$t'_1$ canonical path, with $p$ being closer to $t$ than $q$. We modify the image of the $t$-$t'$ critical path as follows. Starting from $t$, we keep its image until right before it enters $p$, and then from there we divert course and go along (close, but in parallel to) the image of the $t_1$-$t'_1$ canonical path until right before me meet $q$, and finally from $q$ we follow its old image until we eventually reach $t'$. So the new image of the $t$-$t'$ critical path does not intersect the $t_1$-$t'_1$ canonical path, but they are very close from $p$ to $q$. See \Cref{fig: misplace_2} for an illustration. In a similar way, we modify the images of all other critical paths that intersect twice with the $t_1$-$t'_1$ canonical path, to divert their images to go along with the $t_1$-$t'_1$ canonical path, such that all these new images \zihan{more picture}


\subsubsection{Definitions and notations}

\paragraph{Representation of critical paths and their subpaths.} The Property~\ref{prop: canonical intersection} ensures that, vaguely speaking, for any two non-terminal vertices $u,v$, the number of canonical paths that contain both of them is either $0$ or $1$ (since otherwise there will be two canonical paths crossing twice at $u,v$). 
Therefore, for notational convenience, when $u,v$ both lie on some canonical path $\rho$, we use $(u,v)$ to denote $\rho[u,v]$, the subpath of $\rho$ between $u$ and $v$. The notation is unambiguous even with $\rho$ being absent.
When we want to emphasize the direction of such a critical path/subpath, we will also use the notation $u\to v$ instead of $(u,v)$.
%
When $u,u'$ are terminals lying on the same face, we let $(u,u')$ be the face segment from $u$ to $u'$.
For a terminal $x$, we denote by $F_x$ the face that contains it.


\paragraph{Triangles.}
For three vertices $x,y,z$ where every pair of them appear on some canonical path, we denote by $(x,y,z)$ the triangle formed by $(x,y)$, $(y,z)$ and $(z,x)$.
We can also define triangles for $x,y,z$ where $x,y$ are terminals lying on the same face, and both pairs $x,z$ and $y,z$ lie on the same canonical path.

\paragraph{Areas.}\footnote{We will mostly use the notions areas to describe a type of ``morphing'' process that we will use later. Essentially, morphing corresponds to continuous deformation between a pair of homotopic simple curves (when terminal faces are viewed as holes). For curves $\alpha,\beta$ between terminals, we say that $\alpha$ is homotopic to $\beta$, or equivalently that $\alpha$ can be morphed into $\beta$, iff we can gradually change $\alpha$ to $\beta$, with the endpoints of $\alpha$ always staying on the same faces and the curve not scanning over any terminal faces. To avoid sophisticated formal definitions in topology, instead of saying that two curves are homotopic, we will just say that their enclosed area contains no terminal faces.} For a triangle $(x,y,z)$, we denote by $\textsc{Area}(x,y,z)$ its enclosed region. For critical paths $(x,y)$ and $(x',y')$ where $F_x=F_{x'}$ and $F_y=F_{y'}$, \emph{the area between paths $(x,y),(x',y')$}, denoted by $\textsc{Area}((x,y),(x',y'))$, is defined to be the area enclosed by (i) $(x,y)$; (ii) $(x',y')$; (iii) the segment of $F_x$ between $x,x'$; and (iv) the segment of $F_y$ between $y,y'$. That is, if $(x,y)$ and $(x,y')$ cross at $z$, then it is the union $\textsc{Area}(x,x',z)\cup\textsc{Area}(y,y',z)$; if $(x,y)$ and $(x',y')$ don't cross, then it is the quadrilateral $\textsc{Area}(x,y,y',x')$. See Figure \ref{fig: area_definition} for an illustration. %For any two non-crossing paths $P$ and $Q$ between the same two faces and a direction on $P$ (let's say from $x$ to $y$), let $\textsc{Area}(P,Q)$ denote the area enclosed by $P$ and $Q$ on the left of $P$ when going from $x$ to $y$ along $P$.

\begin{figure}[h]
\centering
\subfigure
{\scalebox{0.25}{\includegraphics{figs/wings_definition.png}}}
\hspace{1.0cm}
\subfigure
{\scalebox{0.25}{\includegraphics{figs/area_between_definition.png}}}
\caption{An illustration of $\textsc{Area}((x,y),(x',y'))$ (pink). \label{fig: area_definition}}
\end{figure}  

\paragraph{Relative directions between critical paths.} 
Fix two faces $F_1$ and $F_2$, where $F_1$ has lower face index than $F_2$. Consider the canonical path $\gamma$ between a vertex $u\in F_1$ and a vertex $v\in F_2$.
Assume this canonical path is formed by primary paths $(u,v')$ and $(v,u')$, and their paired secondary paths are $(u,v'')$ and $(v,u'')$, respectively.
The claim is that when we move from $u$ to $v'$ along the path $(u,v')$ and reaches $w$, the subpath $(w,v)$ exits from its left side (see Figure \ref{fig:orientation}).

%To see this, let $u\to v'$ and $v\to u'$ denote the two primary paths which cross at $w$, forming the canonical path between $u$ and $v$ and let $u\to v''$ and $v\to u''$ be the corresponding secondary paths crossing at $w'$ . 
This can be proved as follows.
Since $F_1$ has lower face index, the path $u\to v'$ leaves from the right side of the path $u\to v''$ at $u$. Similarly, the path $v\to u'$ leaves from the left side of $v\to u''$ at $v$. From the figure, this immediately implies that when moving $u$ to $v'$, we have $(w,v)$ must be on the left side. For ease of reference later, we state the claim formally as an observation:

\begin{figure}[h]
\centering
\includegraphics[scale=0.4]{figs/Figure_1a.png}
\caption{An illustration of the shape of the $u$-$v$ canonical path (red).\label{fig:orientation}}
\end{figure}
\begin{observation}
Fix two faces $F_1$ and $F_2$, where $F_1$ has a lower face index than $F_2$. Let $(u,v')$ be a primary path from $u$ to $F_2$ and $(v,u')$ be a primary path from $v$ to $F_1$, that form the $u$-$v$ canonical path with bend $w$. Then when we move along $u\to v'$, the subpath $(w,v)$ exits from its left.
    \label{obs:orientations}
\end{observation}

\paragraph{$1$-bend paths.} 
We say a path $Q$ between terminals $a$ and $b$, where $F_a$ has lower face index than $F_b$, is a \emph{$1$-bend path} with bend $c$, iff (i) $Q[a,c]$ is a subpath of some primary path $P_1$ from $a$ to $F_b$; (ii) $Q[b,c]$ is a subpath of some primary path $P_2$ from $b$ to $F_a$; and (iii) when moving along $P$ from $a$ to the endpoint on $b$, $(c,b)$ exits from its the left (i.e., $Q$ satisfies the property  in  \Cref{obs:orientations}). We say that a 1-bend path $Q$ formed by primary paths $P_1$ and $P_2$ is  \emph{safe}, iff $\textsc{Area}(P_1,P_2)$ contains no terminal faces. All canonical paths are safe 1-bend paths, but the converse is not true.


\paragraph{Bad pairs.} 
Let $P$ be a critical path and $Q$ be a safe 1-bend path. We say $(P,Q)$ is a \emph{bad pair}, iff (i) there exist faces $F_a,F_b$, such that both $P,Q$ connect a terminal on $F_a$ to a terminal on $F_b$; (ii) $P$ and $Q$ cross more than once; and (iii) the endpoints of $P$ and $Q$ are four distinct terminals. Note that such a pair $P$ and $Q$ cross exactly twice, since pairs of critical paths cross at most once and $Q$ is formed by two critical paths. We say that $(P,Q)$ is a \emph{canonical bad pair} iff $(P,Q)$ is a bad pair and $Q$ is a canonical path. 





\subsubsection{Description of an iteration}

While a canonical bad pair exists, our algorithm will iteratively decrease the number of bad pairs from $\hat H$ by moving some of the critical paths.
As canonical paths are defined as concatenations of subpaths of critical paths, they will also evolve during this iterative process. We will then show that the final $\hat{H}$ which contains no canonical bad pairs also must satisfy Property \ref{prop: intersect iff}. If in the current graph $\hat H$, there are no canonical bad pairs, then we return $\hat H$ and terminate the algorithm. Otherwise, the first step is to find a bad pair with some additional properties.

Let $(P,Q)$ be a bad pair. Suppose $P$ is a primary path from $a$ to $b$. %Let $F_a$ and $F_b$ denote the faces which $a$ and $b$ lie on, respectively, where $F_t$ denote the face that terminal $t$ lies on. 
Let $(c,f)$ and $(e,g)$ denote the primary paths forming $Q$, where $c$ lies on $F_a$ and $e$ lies on $F_b$. Let $d$ denote the bend of $Q$, and let $h$ and $i$ denote the two crossings between $P$ and $Q$, where $h$ is closer to $a$ and $i$ is closer to $b$. See Figure \ref{fig:notations} for an illustration.

\begin{figure}[h]
\centering
\includegraphics[scale=0.35]{figs/Figure_notations.png}
\caption{Notations for relevant nodes induced by a bad pair.\label{fig:notations}}
\end{figure}

% We say that $(P,Q)$ is a \emph{minimal pair}, iff
% \begin{itemize}
%     \item no critical path between $F_a$ and $F_b$ cross $(h,d)$ and $(d,i)$;
%     \item no primary path from $F_b$ to $F_a$ crosses $(h,i)$ and $(h,d)$;
%     \item no primary path from $F_a$ to $F_b$ crosses $(h,i)$ and $(i,d)$; and
%     \item no primary path from $F_b$ to $F_a$ crosses $(h,i)$ and $(i,d)$.
% \end{itemize}
% If only the first three conditions hold, then we say $(P,Q)$ is an \emph{almost minimal pair}.

We say that $(P,Q)$ is a \emph{minimal bad pair}, iff
\begin{itemize}
    \item no critical path between $F_a$ and $F_b$ crosses $(h,d)$ and $(d,i)$ and
    \item no primary path from $F_b$ to $F_a$ crosses $(h,i)$ and forms a safe 1-bend path with $(a,b)$.
    %\item no primary path $R$ from $F_b$ to $F_a$ crosses $(h,i)$ and forms a 1-bend path with $(a,b)$ such that $\textsc{Area}(R,(a,b))$ has no faces.
\end{itemize}


\paragraph{Step 1. Finding a minimal bad pair.} Given a canonical bad pair, we find a minimal bad pair through an iterative process as follows.
%
We maintain the bad pair $(\hat{P},\hat{Q})$ throughout the process, initialized as the given canonical bad pair $(P,Q)$, and repeat the following until convergence:
\begin{itemize}
    \item \textbf{Case 1. There exists a critical path $R$ between $F_a$ and $F_b$ which crosses $(h,d),(d,i)$.} We update $\hat P$ to be $R$, and then we update $a,b,\ldots,h,i,F_a,F_b$ with respect to $(\hat P,\hat Q)$ accordingly.
    \item \textbf{Case 2. There exists a primary path $R$ from $F_b$ to $F_a$ that crosses $(h,i)$ and forms a safe $1$-bend path with $(a,b)$.} Note that in this case, $R$ must also cross $(h,g)$, because $R$ must leave the $(h,g,b)$ region. We update $\hat Q$ to be the 1-bend path formed by $(c,g)$ and $R$. Then we update $a,b,\ldots,h,i,F_a,F_b$ with respect to $(\hat P,\hat Q)$ accordingly. Note that $\hat{Q}$ satisfies the property in Observation \ref{obs:orientations} since $R$ has the same structure as primary path $(e,f)$ which formed a 1-bend path with $(c,g)$, so $\hat{Q}$ is indeed a 1-bend path. Furthermore, by Observation \ref{obs:wing-no-face} below, $\hat{Q}$ is actually a safe 1-bend path. 
\end{itemize}

%It is a priori unclear that the process above converges. We first finish describing the algorithm assuming convergence, and give the convergence proof later. Thus, assume we have a minimal bad pair $(P,Q)$ given by the above process. 



% \subsubsection*{Step 1. Find an almost minimal pair}


% We will show that we can find a bad pair $(P,Q)$, such that either $(P,Q)$ is a minimal pair, or $P$ is not a primary path. We start by showing how to find an almost minimal pair, and use that procedure as a building block to find a minimal pair. 


% \begin{observation}\label{obs:almost-minimal-triangle}
% Given a bad pair $(P,Q)$, we can efficiently find a bad pair $(\hat P,\hat Q)$ such that either $(\hat P,\hat Q)$ is almost minimal or $\hat P$ is not a primary path.
% \end{observation}

% \begin{proof}
% Let $(P,Q)$ be a bad pair, where $P$ is a primary path connected terminal $a$ on face $F_a$ to terminal $b$ on face $F_b$.
% We perform iterations, starting with $(\hat P,\hat Q)=(P,Q)$.

% \textbf{Case 1. There exists a critical path $R$ between $F_a$ and $F_b$ which crosses $(h,d),(d,i)$.} 
% We update $\hat P$ to be $R$, and then we update $a,b,h,i,d,F_a,F_b$ with respect to $(\hat P,\hat Q)$ accordingly.

% \textbf{Case 2. There exists a primary path $R$ from $F_b$ to $F_a$ which crosses $(h,d),(h,i)$.}
% We update $\hat Q$ to be the 1-bend path formed by (i) the critical path that contains $(h,d)$; and (ii) path $R$.
% Then we update $a,b,h,i,d,F_a,F_b$ with respect to $(\hat P,\hat Q)$ accordingly.

% \textbf{Case 3. There exists a primary path $R$ from $F_a$ to $F_b$ which crosses $(h,i),(i,d)$.}
% We update $\hat Q$ to be the 1-bend path formed by (i) the critical path that contains $(i,d)$; and (ii) path $R$.
% Then we update $a,b,h,i,d,F_a,F_b$ with respect to $(\hat P,\hat Q)$ accordingly.

% This completes the description of an iteration.
% If the resulting $\hat P$ is not a primary path, then we terminate and return the current $(\hat P,\hat Q)$. If there still exists a path satisfying the condition of Cases 1--3, we perform an iteration similarly. Otherwise, we terminate and return $(\hat P,\hat Q)$.

% Observe that after each iteration, the new triangle $(h,i,d)$ is strictly contained within the previous triangle, so the process terminates efficiently. The resulting pair $(\hat P,\hat Q)$ is almost minimal, by definition.   
% \end{proof}    

% \subsubsection*{Step 2. Find a minimal pair}

% \begin{claim}\label{cl:minimal-triangle}
% Given a bad pair $(P,Q)$, we can efficiently find a bad pair $(\hat P,\hat Q)$ such that either $(\hat P,\hat Q)$ is minimal or $\hat P$ is not a primary path.
% \end{claim}
% \begin{proof}
% Similar to \Cref{obs:almost-minimal-triangle}, we perform iterations, starting with $(\hat P,\hat Q)=(P,Q)$.
%     \begin{itemize}
%         \item[1.] Apply \Cref{obs:almost-minimal-triangle} to the current bad pair $(\hat P, \hat Q)$ to obtain a pair $(P', Q')$.
%         If $P'$ is not primary, we terminate and return $(P',Q')$. Otherwise, update $a,b,h,i,d,F_a,F_b$ with respect to $(P', Q')$ accordingly.
%         \item[2.] Find the primary path $P''$ that, among all primary paths from $F_b$ to $F_a$ that cross $(h,i),(d,i)$, creates an intersection with $(d,i)$ closest to $d$ (if there are no such pairs, we terminate and output the current pair $(\hat P,\hat Q)$). Note that such a path (starting at $F_b$) must cross $(h,i)$ before $(i,d)$; otherwise, the structure of $P''$ and $(a,b)$ contradicts Observation \ref{obs:orientations}. Let $x$ denote its crossing with $(d,i)$ and let $y$ denote its crossing with the critical path containing $(h,d)$. 
%         See Figure [minimal] for an illustration.
%         \item[3.] Update $\hat P$ to be the critical path which $(d,x)$ lies on and let update $\hat Q$ to be the 1-bend path formed by the critical paths which $(d,y)$ and $(x,y)$ lie on.
%     \end{itemize}

%     It is clear that when the process terminates, the resulting pair $(\hat P,\hat Q)$ is minimal. Thus, it suffices to prove convergence of the process. We split this proof into two observations.



% For each $t>1$, let $\Delta_t$ be the triangle $(d,x,y)$ obtained after the $t^{th}$ iteration, and we call $(d,x)$ its bottom.

%     \begin{observation}
% \label{obs: cross bottom}
%         For each $t>1$, no primary paths between $F_a$ and $F_b$ cross the bottom of $\Delta_{t+1}$.
%     \end{observation}
%     \begin{proof}
%         We prove by contradiction. Assume first that a primary path from $F_a$ to $F_b$ crossed $(d,x)$ (the bottom of $\Delta_{t+1}$). Such a path must also cross either $(h,d)$ or $(h,i)$ in order to exit the triangle $(h,d,i)$. In either of these cases, it causes a contradiction to the almost minimality of the bad pair forming the $(h,i,d)$ triangle, since by definition no critical paths between $F_a$ to $F_b$ are allowed to cross both $(h,d)$ and $(d,i)$ or $(h,i)$ and $(d,i)$. Assume now that a primary path from $F_b$ to $F_a$ crossed $(d,x)$. Such a path must also cross either $(h,d)$ or $(h,i)$ in order to exit the triangle $(h,d,i)$. It may not cross $(h,d)$ by definition of almost minimal pair. It may not cross $(h,i)$ either, since otherwise it causes a contradiction to the definition of $x$, which is the closest-to-$d$ intersection between any $F_b$ to $F_a$ primary path and $(d,i)$. %Combining the two cases gives the desired claim.
%     \end{proof}

%     \begin{observation}
% \label{obs: triangle}
%         For each $t>1$, $\Delta_{t+1}\subseteq \Delta_{t}$.
%     \end{observation}
%     \begin{proof}
%         %We wish to show that the $(d,x,y)$ triangle lies inside $\Delta_t$. 
%         Let $\Delta_t=(d',x',y')$ and $\Delta_{t+1}=(d,x,y)$. Observe that each subpath $(d,x)$, $(x,y)$, and $(d,y)$ lies on some critical path, which means it can cross $\Delta_t$ at most twice. Thus, to show each side of $\Delta_{t+1}$ lies in $\Delta_t$ (thus implying that the triangle is a subset), it suffices to show that each of $d,x,y$ lies inside $\Delta_t$.
%         %\zihan{more explanation here: why three points guarantee the whole triangle, why a strict subset?} 
        
%         It is clear that $(d,x)$ lies inside $\Delta_t$ since $(d,h,i)$ lies inside $\Delta_t$ (due to the process in Observation \ref{obs:almost-minimal-triangle}). It thus remains to show that $y$ lies inside $\Delta_t$. From \Cref{obs: cross bottom}, no primary paths between $F_a$ and $F_b$ cross the bottom of $\Delta_{t}$. Thus, $P''$ must cross sides $(d',y'),(x',y')$. Observe that the crossing between $P''$ and $(d',y')$ is closer to $d'$ than the crossing between $P'$ and $(d',y')$; otherwise, $P''$ would cross $P'$ twice. But if $y$ was outside $\Delta_t$, then $P''$ would cross the critical path containing $(h,d)$ outside of $\Delta_t$. But this would imply that there is a crossing between $P''$ and $(d',y')$ which lies above that of $P'$, a contradiction since there is only one crossing between $P''$ and $(d',y')$.
%     \end{proof}

% When $\Delta_{t+1}=\Delta_t$, the algorithm has converged and we have found a minimal pair $(\hat P,\hat Q)$.
% From \Cref{obs: triangle}, the process terminates efficiently.
% \end{proof}

%The intuition of convergence is in each iteration, the triangle is moving ``upwards'' in the region. Specifically, in Case 1, $\hat{P}$ is updated so that the base of the triangle $(h,i)$ is higher than it was before. In Case 2, $\hat{P}$ isn't updated so the base of the triangle shrinks but doesn't move. The difficulty of formalizing this proof is showing that a cycle cannot occur. Even though the base of the triangle is always moving upwards, it is still possible for the triangle to wrap around either $F_a$ or $F_b$, and cause a cycle. We show below that this cannot happen. \george{change later maybe}
Now, we show that Step 1 indeed finds a minimal bad pair. We first need to define some language. For iteration $t$ of the process, let $a_t,b_t,\ldots,i_t$ denote the corresponding points in the graph. Let $\Delta_t=\textsc{Area}(h_t,d_t,i_t)$ and $W_t=\textsc{Area}(c_t,d_t,f_t)\cup\textsc{Area}(g_t,d_t,e_t)$.


\begin{observation}
    For any $t\ge0$, $\Delta_t$ and $W_t$ don't contain any terminal faces. \label{obs:wing-no-face}
\end{observation}
\begin{proof}
We prove the claim by induction on $t$. For the base case: $W_0$ and $\Delta_0$ don't contain any terminal faces since $Q$ is a canonical path. For the induction step: If $\hat{P}$ is updated, then $\Delta_{t+1}$ is a subset of $\Delta_t$ and $W_{t+1}=W_t$. By induction, $W_{t+1}$ and $\Delta_{t+1}$ don't contain any terminal faces. If $\hat{Q}$ is updated, let $R$ be the path which forms a 1-bend path with $(a_t,b_t)$ and causes the update to $\hat{Q}$. By our choice of $R$,  $\textsc{Area}(R,(a_t,b_t))$ contains no terminal faces. In {Figure \ref{fig:wing-induction}}, we see that in both cases, $\Delta_{t+1}\cup W_{t+1}$ is a subset of  $\Delta_{t}\cup W_t\cup\textsc{Area}(R,(a_t,b_t))$, completing the proof since the latter doesn't contain terminal faces. 
\end{proof}

\begin{figure}[h]
\centering
\subfigure
{\scalebox{0.22}{\includegraphics{figs/Figure_wing_induction_1.png}}}
\hspace{1.0cm}
\subfigure
{\scalebox{0.22}{\includegraphics{figs/Figure_wing_induction_2.png}}}
\caption{In the figures, the pink shaded region is $\textsc{Area}(R,(a_t,b_t))$, the yellow striped region is $\Delta_{t}\cup W_t$, and the region outlined in a red curve is $\Delta_{t+1}\cup W_{t+1}$. In both cases, $\Delta_{t+1}\cup W_{t+1}$ is a subset of the union $\Delta_{t}\cup W_t\cup\textsc{Area}(R,(a_t,b_t))$.}\label{fig:wing-induction}
\end{figure} 

\begin{observation}
    $\textsc{Area}((a,b),(a_t,b_t))$ doesn't contain any terminal face.
\end{observation}
\begin{proof}
    We first prove that $\textsc{Area}((a_t,b_t),(a_{t+1},b_{t+1}))$ doesn't contain any terminal faces. There are two cases based on if $\hat P$ or $\hat Q$ is updated. When $\hat{Q}$ is updated, we have that $(a_t,b_t)=(a_{t+1},b_{t+1})$ so there is nothing to prove. When $\hat P$ is updated, we have that $(a_{t+1},b_{t+1})$ crosses through the two sides, $(h_t,d_t)$ and $(d_t,i_t)$, of the triangle. In particular, this means that $(a_{t+1},b_{t+1})$ crosses $(c_t,g_t)$ and $(e_t,f_t)$, so the entirety of $(a_{t+1},b_{t+1})$ lies inside $W_t\cup\Delta_t$. Since $(a_t,b_t)$ also lies in $W_t$, we have that $\textsc{Area}((a_t,b_t),(a_{t+1},b_{t+1}))$ also does. But since $W_t$ and $\Delta_t$ don't contain any terminal faces by Observation \ref{obs:wing-no-face}, this implies the claim.

    Now, we complete the proof by induction. The base case follows by the above discussion, since $(a,b)=(a_0,b_0)$ by definition. Suppose that we have $\textsc{Area}((a,b),(a_t,b_t))$ contains no terminal faces. We proved above that $\textsc{Area}((a_t,b_t),(a_{t+1},b_{t+1}))$ also contains no terminal faces. Recall that the area between two paths can be viewed as the minimal region which is passed over in a continuous deformation between the two paths. Since such a continuous deformation between $(a,b)$ and $(a_t,b_t)$ can be composed together with one between $(a_t,b_t)$ and $(a_{t+1},b_{t+1})$ to form one between $(a,b)$ and $(a_{t+1},b_{t+1})$, we have that $\textsc{Area}((a,b),(a_{t+1},b_{t+1}))$ is a subset of the union of $\textsc{Area}((a,b),(a_{t},b_{t}))$ and $\textsc{Area}((a_t,b_t),(a_{t+1},b_{t+1}))$. Since the latter two areas don't contain any faces, neither does the former area, completing the proof. See Figure \ref{fig:casework} for an illustration of some examples.
    \begin{figure}[h]
    \centering
    \subfigure
    {\scalebox{0.33}{\includegraphics{figs/Figure_casework_1.png}}}
    \hspace{1.0cm}
    \subfigure
    {\scalebox{0.33}{\includegraphics{figs/Figure_casework_2.png}}}
    \hspace{1.0cm}
    \subfigure
    {\scalebox{0.33}{\includegraphics{figs/Figure_casework_3.png}}}
    \caption{An illustration of examples of $\textsc{Area}((a,b),(a_{t},b_{t}))$ in orange, $\textsc{Area}((a_t,b_t),(a_{t+1},b_{t+1}))$ in pink, and $\textsc{Area}((a,b),(a_{t+1},b_{t+1}))$ outlined in red. In each of the three examples, we can see that $\textsc{Area}((a,b),(a_{t+1},b_{t+1}))$ is a subset of $\textsc{Area}((a,b),(a_{t},b_{t}))\cup\textsc{Area}((a_t,b_t),(a_{t+1},b_{t+1}))$. \label{fig:casework}}
    \end{figure}      
    % We split the analysis into a few cases. 
    % \begin{itemize}
    %     \item Suppose $(a_t,b_t)$ doesn't cross $(a,b)$. Then $(a_{t+1},b_{t+1})$ either:
    %     \begin{itemize}
    %         \item crosses neither $(a_t,b_t)$ nor $(a,b)$.
    %         \item crosses both $(a_t,b_t)$ and $(a,b)$.
    %         \item crosses $(a_t,b_t)$ but not $(a,b)$.
    %     \end{itemize}
    %     \item Otherwise, suppose $(a_t,b_t)$ does cross $(a,b)$. Then $(a_{t+1},b_{t+1})$ either:
    %     \begin{itemize}
    %         \item crosses neither $(a_t,b_t)$ nor $(a,b)$.
    %         \item crosses both $(a_t,b_t)$ and $(a,b)$.
    %         \item crosses $(a_t,b_t)$ but not $(a,b)$.
    %         \item crosses $(a,b)$ but not $(a_t,b_t)$.
    %     \end{itemize}
    % \end{itemize}
    % In each of these cases, we see in \textcolor{red}{Figures Casework} that $\textsc{Area}((a,b),(a_{t+1},b_{t+1}))$ is a subset of the union of $\textsc{Area}((a,b),(a_{t},b_{t}))$ and $\textsc{Area}((a_t,b_t),(a_{t+1},b_{t+1}))$. Since the latter two areas don't contain any faces, neither does the former area, completing the proof.
\end{proof}

 


\begin{claim}\label{cl:feng-ding}
    Let $b'$ denote the terminal on $F_b$ that lies clockwise right next to $b$ let $S$ denote the critical path between $a$ and $b'$ (see left one of Figure \ref{fig:feng-ding}). Then $\Delta_t$ never crosses $S$.
\end{claim}
\begin{figure}[h]
\centering
\subfigure
{\scalebox{0.22}{\includegraphics{figs/Figure_feng_ding_1.png}}}
\hspace{1.0cm}
\subfigure
{\scalebox{0.22}{\includegraphics{figs/Figure_feng_ding_2.png}}}
\caption{In the left figure, we assert that the critical path $(a,b')$ must be the blue path. The other possibilities are illustrated by the red and orange drawings of $(a,b')$. The red path is not possible because it crosses $(c,g)$ twice and the yellow path is actually equivalent to the blue path. The right figure is for reference in the proof below.}\label{fig:feng-ding}
\end{figure} 
\begin{proof}
    First, observe that the region between $R=(a,b')$ and $(a,b)$ must contain a face $F$ because they are a pair of not equivalent critical paths. Next, suppose that some triangle crosses $S$. Consider the first such triangle; let's say it occurred at iteration $T$. Let $a_T,b_T,\ldots,i_T$ denote the analogous nodes induced by this triangle, such that $\textsc{Area}(h_T,d_T,i_T)$ is the triangle. Note that for the first time a triangle crosses $S$, we must have that $S$ crosses $(h_T,d_T)$ and $(d_T,i_T)$; otherwise the triangle would have crossed $S$ at iteration $T-1$ as well (see {Figure \ref{fig:triangle_crossing}, and the corresponding caption}).
\begin{figure}[h]
\centering
\subfigure
{\scalebox{0.22}{\includegraphics{figs/Figure_triangle_crossing_1.png}}}
\hspace{1.0cm}
\subfigure
{\scalebox{0.22}{\includegraphics{figs/Figure_triangle_crossing_2.png}}}
\hspace{1.0cm}
\subfigure
{\scalebox{0.22}{\includegraphics{figs/Figure_triangle_crossing_3.png}}}
\caption{In each of the figures, $(h,d,i)$ represents $\Delta_{T-1}$, the pink area represents $\Delta_{T}$, and the red curve represents the path $S$. In the top two figures, any path $S$ which crosses $\Delta_T$ twice must also cross $\Delta_{T-1}$ twice. In the bottom figure, a path $S$ which crosses $\Delta_T$ twice but doesn't cross $\Delta_{T-1}$ twice must look like the red path. In particular, note that the red path indeed crosses $(h_T,d_T)$ and $(d_T,i_T)$, as claimed in the proof.}\label{fig:triangle_crossing}
\end{figure} 

    We claim that the area $A$ between $f_T\to d_T\to g_T$ and $(a,b)$ contains the face $F$. The proof of the claim is illustrated in the right side of Figure \ref{fig:feng-ding}, and explained below. Since $(c_T,d_T)$ crosses $S$, we know that $(d_T,g_T)$ cannot cross $S$ again; otherwise, two critical paths would cross twice. Similarly, $(e_t,d_T)$ crosses $S$ so $(d_T,f_T)$ cannot cross $S$ again. Thus, $f_T\to d_T\to g_T$ doesn't cross $S$. Next, observe that $g_T$ is further than $b'$ from $b$ in the clockwise direction on $F_b$. % and $f_T$ is (trivially) further than $a$ from $a$ in the counter-clockwise direction on $F_a$.
    Since we have shown that $f_T\to d_T\to g_T$ doesn't cross $S$, the area between $f_T\to d_T\to g_T$ and $(a,b)$ must contain the area between $S$ and $(a,b)$. Since the latter area contains a face, the former area must as well.
    
    Note that the difference $A-\textsc{Area}((a_T,b_T),(a,b))$ is a subset of $W_T$. But from Observation \ref{obs:wing-no-face}, $W_T$ cannot contain a face, a contradiction. Hence, such a triangle $\textsc{Area}(h_T,d_T,i_T)$ cannot exist.
\end{proof}

\begin{observation}
    Let $U=\textsc{Area}((a,b),(a,b'))$. We have $\Delta_t\subseteq U$ for each iteration of Step 1.
\end{observation}
\begin{proof}
We prove by induction on $t$. The first triangle clearly lies inside $U$. For the induction step, there are two cases based on whether $\hat{P}$ or $\hat{Q}$ are modified. If $\hat{P}$ is modified, then the region enclosed by the triangle is a subset of that of the old triangle, so it must still lie inside $U$ by induction. If $\hat{Q}$ is modified, let $R$ denote the path which forms a safe 1-bend path with $(a,b)$ causing the update to $\hat{Q}$. In the new triangle, the three endpoints are $h_t$, the crossing between $(h_t,i_t)$ and $R$, and the crossing between $(h_t,g_t)$ and $R$. Label the endpoints of the new triangle $h_{t+1}$, $i_{t+1}$, and $d_{t+1}$ so that $h_{t+1}=h_t$,  The former two points lie inside the previous triangle, so they already lie in $U$. Since the triangle never crosses $R$, $(h_{t+1},d_{t+1})$ and $(d_{t+1},i_{t+1})$ can't cross $R$. We claim that $(h_{t+1},d_{t+1})$ and $(d_{t+1},i_{t+1})$ also can't cross $(a,b)$. If they did, Observation \ref{obs:orientations} would mean there are two critical paths which cross twice (see {Figure \ref{fig:stay_in_region}}). Since $R$ and $(a,b)$ are the boundary of the region $U$ and $(h_{t+1},d_{t+1})$ and $(d_{t+1},i_{t+1})$ never cross the boundary, each side of $\Delta_{t+1}$ lies inside $U$. This implies that  $\Delta_{t+1}$ lies inside $U$. 
\end{proof}

\begin{figure}[h]
\centering
\subfigure
{\scalebox{0.22}{\includegraphics{figs/Figure_stay_in_region_1.png}}}
\hspace{1.0cm}
\subfigure
{\scalebox{0.22}{\includegraphics{figs/Figure_stay_in_region_2.png}}}
\caption{The right figure illustrates what actually happens: the triangle $(h_{t+1},d_{t+1},i_{t+1})$ stays in the region $U$. The left figure illustrates what cannot happen: the triangle $(h_{t+1},d_{t+1},i_{t+1})$ crossing through $(a,b)$. This is not possible because it forces the subpath $(e_{t+1},d_{t+1})$ to cross $(a,b)$ twice, indicated by the red crossing nodes, a contradiction.\label{fig:stay_in_region}}
\end{figure} 

\begin{claim}\label{cl:step1-convergence}
    Step 1 converges efficiently to a minimal bad pair $(P,Q)$.
\end{claim}
\begin{proof}
    It suffices to prove Step 1 converges, since the $(\hat P,\hat Q)$ pair at convergence must be a minimal bad pair. Suppose the process didn't converge. Since %the triangle $\Delta_t$ always lies inside the region $U$ and 
    there are polynomially many triangles in $U$, triangles must repeat. Let $\Delta_{t_1}=\Delta_{t_2}$ be a pair of repeated triangles. Consider the union of areas $A_t:=\textsc{Area}((h_t,i_t),(h_{t+1},i_{t+1}))$ for $t_1\le t<t_2$. The border of this area starting at $h_{t_1}$ forms a loop which (1) starts and ends at $h_{t_1}$ and (2) stays in the region $U$. In particular, this implies that there must exist some iteration $t$ where $A_t$ crosses back through $(a_{t_1},b_{t_1})$ (see {Figure \ref{fig:loop}}). Let $t^*$ denote the first iteration between $t_1$ and $t_2$ where $A_{t^*}$ crosses back through $(a_{t_1},i_{t_1})$. By similar arguments in Claim \ref{cl:feng-ding} and {Figure \ref{fig:triangle_crossing}}, both $(h_{t^*},d_{t^*})$ and $(i_{t^*},d_{t^*})$ cross $(a_{t_1},b_{t_1})$. But Observation \ref{obs:orientations} implies such an $A_t$ cannot exist, since there would have to be two critical paths which cross each other twice (see discussion in the caption of Figure \ref{fig:loop}).
\end{proof}

\begin{figure}[h]
\centering
\subfigure
{\scalebox{0.22}{\includegraphics{figs/Figure_loop_1.png}}}
\hspace{1.0cm}
\subfigure
{\scalebox{0.22}{\includegraphics{figs/Figure_loop_2.png}}}
\caption{The left figure illustrates the sequence of areas $A_t$. In particular, there must exist some $A_{t*}$ which crosses back through $(a_{t_1},b_{t_1})$. We claim that such an $A_{t*}$ cannot exist. As illustrated in the right figure, such an $A_{t*}$ would imply that the corresponding $(e_{t*},d_{t*})$ would cross with $(a_{t_1},b_{t_1})$ twice, indicated by the red crossing nodes, a contradiction. \label{fig:loop}}
\end{figure} 


\paragraph{Step 2: Rerouting the minimal bad pair.}
First, we re-define some notation. Let $(a,b)$ denote the critical path $P$ and let $(c,g)$ and $(e,f)$ denote the critical paths which form $Q$, with $(c,g)$ and $(e,f)$ crossing at $d$ and $c$ lying on the same face as $a$. Let $h$ and $i$ denote the two crossings between $P$ and $Q$, where $h$ is closer to $a$ and $i$ is closer to $b$. From now on, we will refer to paths $P$ and $Q$ as $(a,b)$ and $c\to d\to e$, so that the notations $P$ and $Q$ can be used for arbitrary pairs of paths. For simplicity, we will denote $\textsc{Area}(h,i,d)$ as \emph{the triangle}. 

In order to fix the bad pair, we reroute $(a,b)$ to follow the path $(a,b)$ until it hits the triangle at $h$. When the path reaches $h$, we reroute it above the top angle of the triangle by first going along (but above) the $(h,d)$ subpath and then going along (but above) the $(d,i)$ subpath. When the rerouting hits $i$, we then follow the $(i,b)$ subpath of the original path $(a,b)$ until we reach $b$. This rerouting is illustrated in {Figure \ref{fig:2a}.} Clearly, the rerouted version of $(a,b)$ no longer crosses the $1$-bend path $c\to d\to e$ so we have fixed the bad pair. 

\begin{figure}[h]
    \centering
    \includegraphics[width=0.7\linewidth]{figs/Figure_2a.png}
    \caption{Illustrates the rerouting of $(a,b)$ over the triangle $(h,i,d)$.}
    \label{fig:2a}
\end{figure}

This completes the description of an iteration, and also completes the description of the algorithm. 

\subsubsection{Analysis of the algorithm}

\begin{observation}\label{obs:prop123}
    After one iteration, Properties \ref{prop: f-face instance}, \ref{prop: canonical intersection}, and \ref{prop: region} continue to hold.
\end{observation}
\begin{proof}
Property \ref{prop: f-face instance} continues to hold, as we have not modified the locations of any faces or terminals.
    %, and we maintain a planar embedding of $\hat H$ (each crossing in its drawing represents an intersection, which is a vertex in $\hat H$). 
%
    For Property \ref{prop: canonical intersection}, consider any pair of critical paths $P$ and $Q$. If neither $P$ nor $Q$ are the $(a,b)$ path, the claim is trivial since they are not rerouted. If $P=(a,b)$, then we claim $Q$ crosses $P$ after the rerouting if and if $Q$ crossed $P$ before. To see this, observe that it suffices to show that $Q$ crosses $(h,i)$ if and only if $Q$ crosses either $(d,h)$ or $(d,i)$. But this follows by minimality of the triangle since no path can cross both $(d,h)$ and $(d,i)$, and any path must cross the boundary of the triangle an even number of times. 
    
    Next, we show that Property \ref{prop: region} continues to hold. For any two terminals $t_1$ and $t_2$, let $A$ denote the region between the $t_1$-$t_2$ canonical path and the $t_1$-$t_2$ shortest path before the iteration and let $A'$ denote the region after the iteration. Observe that any canonical path is only modified around an area which is a subset of $\textsc{Area}(h,i,d)$ which doesn't contain any terminal faces or terminals (Observation \ref{obs:wing-no-face}). Consequently, the difference between $A$ and $A'$ doesn't enclose any terminal or face. Since $A$ didn't contain any face or any other terminals, this implies inductively that $A'$ also doesn't contain any face or any other terminals. Consequently, Property \ref{prop: region} also continues to hold. 
    %To see that the difference between $A$ and $A'$ doesn't enclose any face or terminal, observe that the canonical paths are only modified if they go through the region $A_0$ enclosed by the $t-t'$ critical path and the $t_1-t_1'$ canonical path. When these paths are modified, the resulting areas will then differ by at most $A_0$, by construction. But observe that $A_0$ does not contain any faces or terminals, so we are done. Applying the above proof to each iteration implies that Property \ref{prop: region} holds for the graph and critical paths at the end of the process.
\end{proof}

Next, we will show that Property \ref{prop: intersect iff} also holds for the final graph. We start by showing that the number of bad pairs decreases in each iteration in following arguments.

\begin{observation}\label{obs:same-faces}
If two safe 1-bend paths cross twice, their endpoints  lie on the same two terminal faces.
\end{observation}
\begin{proof}
    Let $P$ be a safe 1-bend path from $u$ to $v$ with bend $w$ and let $Q$ be a safe 1-bend path from $x$ to $y$ with bend $z$. Let $(u, v')$ and $(v, u')$ denote the critical paths forming $P$ and let $(x, y')$ and $(y, x')$ denote the critical paths forming $Q$. Assume for contradiction that the endpoint terminals of $P$ and $Q$ lie on at least three terminal faces. Then there must exist at least one endpoint of $P$ which lies on a different face than the other endpoints. Without loss of generality, assume $F_v\neq F_x,F_y$.

    Since $P$ and $Q$ cross twice and neither $\textsc{Area}(u,w,u')\cup \textsc{Area}(v,w,v')$ nor $\textsc{Area}(x,z,x')\cup \textsc{Area}(y,z,y')$ contains any terminal faces, $P$ crosses either $(x, y')$ or $(y, x')$ twice as well. Without loss of generality, let us say that $(u,w)$ and $(w,v)$ each crosses $(x,y')$. Since $F_v\neq F_x,F_y$ and $\textsc{Area}(u,w,u')\cup \textsc{Area}(v,w,v')$ doesn't contain any other terminal faces, we know that the terminals $x,y,x',y'$ don't lie in $\textsc{Area}(w,v,v')$. Combined with the fact that $(x,y')$ crosses $(w,v)$, this implies that $(x,y')$ also crosses $(w,v')$. But this means that $(x,y')$ crosses $(u,v')$ twice ($(u,w)$ and $(w,v')$), a contradiction. See Figure \ref{fig:two-face-reduction}.
    \begin{figure}[h]
        \centering
        \includegraphics[width=0.6\linewidth]{figs/Figure_two_face.png}
        \caption{An illustration of the proof of Observation \ref{obs:same-faces}. The two crossings which must occur between $(u,v')$ and $(x,y')$ are identified in red.}
        \label{fig:two-face-reduction}
    \end{figure}
\end{proof}

\begin{claim}\label{cl:canonical-path-crossing-triangle}
Before the path $(a,b)$ gets rerouted, any safe $1$-bend path $u\to w\to v$ crosses $(h,i)$ at least as many times as it crosses $(i,d)$ and $(h,d)$ combined.
\end{claim}
\begin{proof}
    Recall that $a\to b$ is the primary direction of the $(a,b)$ critical path. The safe 1-bend path can cross the triangle either 0, 2, or 4 times; it is clear that 6 times is never possible, or else there would necessarily exists two critical paths which cross twice. If the path crosses the triangle 0 times, the claim is trivial. If the path crosses the triangle 4 times, observe that the path cannot cross $(h,d)$ and $(i,d)$ a total of 3 or 4 times since that would contradict the minimality of the triangle as at least one primary path forming the 1-bend path would need to cross both $(h,d)$ and $(i,d)$. But if the path crosses $(h,d)$ and $(i,d)$ twice or less, it would cross $(h,i)$ twice or more, and the desired claim holds. 
\begin{figure}[h]
\centering
\subfigure
{\scalebox{0.22}{\includegraphics{figs/Figure_4a_1.png}}}
\hspace{1.0cm}
\subfigure
{\scalebox{0.22}{\includegraphics{figs/Figure_4a_2.png}}}
\hspace{1.0cm}
\subfigure
{\scalebox{0.22}{\includegraphics{figs/Figure_4a_3.png}}}
\hspace{1.0cm}
\subfigure
{\scalebox{0.22}{\includegraphics{figs/Figure_4a_4.png}}}
\caption{In all subfigures, $(u,v)$ and $(v,w)$ are represent by blue lines. The former three cases are not possible by a similar argument as in the caption of \ref{fig:triangle_crossing} via Observation \ref{obs:orientations}. The final case is not possible by minimality of the triangle.}\label{fig:4a}
\end{figure}     
    Finally, consider a safe 1-bend path which creates two crossings with the triangle. By Observation \ref{obs:same-faces}, we may assume without loss of generality that $F_a=F_u$ and $F_b=F_v$. We show that any safe 1-bend path crossing either $(h,d)$ twice, $(i,d)$ twice, or $(h,d)$ and $(i,d)$ contradicts the structure of the 1-bend paths in Observation \ref{obs:orientations}, or the minimality of the $(P,Q)$ bad pair. This is illustrated for the four different cases in the {Figure \ref{fig:4a}}. In the first case, the existence of such a 1-bend path contradicts the minimality of the $(P,Q)$ bad pair. In the remaining three cases, the $c\to d\to e$ and $u\to w\to v$ canonical paths are oriented differently while being between the same two faces, implying that at least one of them contradicts the Observation \ref{obs:orientations} by a similar argument as the proof of Claim \ref{cl:step1-convergence}. Thus, the only possibilities which remain for a 1-bend path to cross the triangle twice are: $(h,d)$ and $(h,i)$, $(h,i)$ and $(h,i)$, or $(h,i)$ and $(i,d)$, satisfying the claim.
    %\textbf{Case 2. A safe 1-bend path creates four crossings with the triangle.} \george{can simplify; only need to consider cases where 1-bend crosses left + right 3 or 4 times; but this is clearly not possible.} Again by Observation \ref{obs:same-faces}, we must have $F_a=F_u$ and $F_b=F_v$. We characterize what such paths must look like. In order for the safe 1-bend path to cross the triangle four times, each subpath of it must cross the triangle twice (since a critical path can cross the triangle at most twice). Furthermore, each subpath cannot cross the same side of the triangle twice, and also cannot cross both $(h,d)$ and $(i,d)$ by definition of minimality. With these constraints, there are only six possible structures of such a canonical path, illustrated in \textcolor{red}{Figure 5a of Planar Emulators Final.} 
    %We show that the first three structures cannot occur. Observe that the first three structures contradict Observation \ref{obs:orientations} once again. Specifically, the first three structures correspond exactly with the three cases considered where the 1-bend path crosses the triangle twice, so \textcolor{red}{Figure 4b of Planar Emulators Final} combined with Observation \ref{obs:orientations} implies that these structures are not possible. Thus, the final three structures are the only possible cases of a 1-bend path crossing the triangle four times. These correspond exactly to crossing: $(h,i)$, $(h,d)$, $(i,d)$, and $(h,i)$, or $(h,i)$, $(h,d)$, $(h,d)$, and $(h,i)$, or $(h,i)$, $(i,d)$, $(i,d)$, and $(h,i)$, satisfying the claim once again.
\end{proof}

Using the above claims and observations, we show that the number of bad pairs decreases.

\begin{figure}[h]
\centering
\subfigure
{\scalebox{0.22}{\includegraphics{figs/Figure_13a_1.png}}}
\hspace{1.0cm}
\subfigure
{\scalebox{0.22}{\includegraphics{figs/Figure_13a_2.png}}}
\caption{In both figures, the blue path represents $Q$ before the rerouting and the orange path represents $Q$ after the rerouting.}\label{fig:13}
\end{figure} 

\begin{observation}\label{obs:bad-pairs-decrease}
    The number of bad pairs $(P,Q)$ decreases after each iteration.
\end{observation}
\begin{proof}
    The $(a,b)$ critical path no longer crosses $c\to d\to e$, so we have fixed a bad pair. Consider any other pair of paths $P$ and $Q$, where $P$ is a critical path and $Q$ is safe 1-bend path. If neither $P$ nor $Q$ involves the $(a,b)$ path, then the claim is trivial as the paths $P$ and $Q$ are not changed in the rerouting. 
    If $P$ involves $(a,b)$, then $P=(a,b)$. To see that $P$ crosses $Q$ fewer times after the rerouting than before, it suffices to show that $Q$ crosses $(d,i)$ or $(d,h)$ at most as many times as it crosses $(h,i)$, which follows from Claim \ref{cl:canonical-path-crossing-triangle}, so $P$ and $Q$ cannot form a new bad pair after the rerouting.
%
    Finally, suppose $Q$ involves $(a,b)$. Let $(u,v)$ be the other primary path that forms $Q$, and let $w$ be the bend of $Q$. This case is illustrated in Figure \ref{fig:13}. Since $(a,b)$ is primary, $(a,b)$ and $c\to d\to e$ are a minimal pair, so no primary path $P'$ from $F_b$ to $F_a$ which crosses the bottom of the triangle can form a safe 1-bend path with $(a,b)$. In particular, this implies that the bend $w$ cannot lie on $(h,i)$, so it must either lie on $(a,h)$ or $(i,b)$. If $w$ is on $(a,h)$, then $Q$ is not changed after the rerouting since the $(a,h)$ subpath is not changed in the rerouting and $(u,w)$ is not rerouted (see right of Figure \ref{fig:13}). If $w$ is on $(i,b)$, then the rerouting of $Q$ only changes the $(h,i)$ subpath to go above the triangle along $(h,d)$ and $(d,i)$ (see left of Figure \ref{fig:13}). In this case, observe that $P$ always crosses $(h,i)$ the same number of times as $(h,d)$ and $(d,i)$ combined due to the minimality of the triangle, so the number of crossings between $P$ and $Q$ cannot increase so $P$ and $Q$ will not form a new bad pair after the rerouting.
\end{proof}





\iffalse
Next, we consider the case when $P$ is a primary path. We assume (without loss of generality) that $a\to b$ is the primary direction for the $(a,b)$ critical path. Unfortunately, when $P$ is a primary path, this rerouting may cause new problems with some other critical-canonical pairs of paths (see, e.g., \textcolor{red}{Figure 2b of Planar Emulators Final}. Thus, we also need to define a rerouting for all critical paths $S$ which cross through the $(h,d,i)$ region so that no new bad critical-canonical pairs are created. Since $(P,Q)$ is a minimal pair, there are no paths crosses both $(d,h)$ and $(d,i)$; thus, $S$ crosses either $(h,d)$ or $(i,d)$ along with $(h,i)$.



We will try to define the rerouting so that the rerouted paths cross if and only if their corresponding original paths cross. In fact, we will end up guaranteeing that the orderings of the crossings are preserved, so in fact the graph induced by the set of critical paths crossing the triangle is isomorphic before and after the rerouting. Informally, we reroute these paths around either the $h$-$i$-$d$ angle or the $i$-$h$-$d$ angle so that $S$ no longer crosses the triangle. Let us denote the region slightly below the $h-i$ subpath by $A$, where the ``height'' of $A$ is $\epsilon$-small meaning that a critical path crosses $A$ if and only if it crosses the $(h,i)$ subpath, and let $t_1,t_2,t_3,t_4$ denote the endpoints of region $A$ as illustrated in \textcolor{red}{Figure 3a of Planar Emulators Final}. We now define the rerouting of the paths $S$ which cross $(h,d)$ or $(i,d)$, along with $(h,i)$. 
%First, we define a partial rerouting to the boundary of $A$. That is, for a critical path $S$ crossing the triangle, we define the parts of the rerouting which are not within the region of $A$. Second, we specify how each of the paths $S$ are rerouted within the region $A$ so that the pairwise crossings are preserved. Our partial rerouting will have the property that two partial reroutings cross if and only if their original critical paths crossed outside of the $(h,b,g,d)$ region. Our rerouting within the region $A$ will have the property that two rerouting cross in $A$ if and only if they cross inside the $(h,b,g,d)$ region.

The rerouting is defined separately depending on whether or not $S$ crosses the $(h,g)$ subpath critical path. Informally, if $S$ crosses $(h,g)$, then we reroute it around the $d$-$h$-$i$ angle and if $S$ doesn't cross $(h,g)$, we reroute it around the $d$-$i$-$h$ angle. Let $S=(u,v)$ so that $u\to v$ crosses $(h,i)$ after crossing $(h,d)$ or $(i,d)$. The formal rerouting is defined in three subpaths in both of the cases as follows:
\begin{enumerate}[label=\roman*.]
    \item The first part of the rerouting defines the subpath of the new path starting at $u$ and ending when the path reaches the region $A$. This is done in two cases:
    \begin{itemize}
        \item[R1.] If $S$ crosses $(h,g)$, let $\ell$ denote the crossing between $S$ and $(h,g)$. The rerouting of $S$ starts at $u$, and follows the original path until almost reaching $\ell$. Then it goes around the angle $i$-$h$-$d$ by following the outside of $(\ell,h)$ and $(h,i)$ until reaching the area $A$. Finally, let $x(S)$ denote the point where the rerouting of $S$ crosses the boundary of $A$. When $S$ is clear from context, we also use $x$ to denote this point.

        For two paths $S_1$ and $S_2$ which are rerouted in this way, the path which crosses the $(h,g)$ subpath closer to $g$ is rerouted closer to the triangle. Formally, let $\ell_1$ and $\ell_2$ denote the respective crossings with the $(h,g)$ subpath. If $\ell_1$ is closer to $g$ on the $(h,g)$ subpath, then $S_1$ is rerouted closer to the triangle; otherwise, $S_2$ is rerouted closer to the triangle. All these paths are rerouted further from the triangle than the rerouting of the $(a,b)$ path.

        In particular, observe that the definition of how close a path is rerouted to the triangle implies that the ordering of the crossings $\ell(S)$ of $S$ with $g\to h$ are the exact same as $x(S)$ on $t_1\to t_2$.
        \item[R2.] Let $S=(u,v)$ be a critical path crossing $(h,i)$ and $(i,d)$, but not $(d,g)$; we will reroute it around the $h$-$i$-$d$ angle, but will need to make additional modifications to the path (all within the $(i,d,g,b)$ region). Let $j$ denote the crossing between $S$ and $(i,d)$. Before defining the rerouting, we define some language for talking about these critical paths. For two critical paths $S_1=(u_1,v_1)$ and $S_2=(u_2,v_2)$ of this type which cross in the region $(i,d,g,b)$, we say that $S_2$ crosses $S_1$ from below if $v_2$ lies closer than $v_1$ to $b$ when looking at the arc $(b,g)$ of the face. Now, we can define the rerouting.
    
        The rerouting of $S$ starts at $v$ and follows the path $S$ until another path $S_1$ crosses $S$ from below. The rerouting then follows $S_1$ instead (slightly above $S_1$ without ever crossing it) until another path $S_2$ crosses $S_1$ from below. This process repeats until the rerouting reaches the boundary of the triangle at the $(i,d)$ subpath (again, without touching it). Next, let $j'$ denote the location where the rerouting reaches the $(i,d)$ subpath. The rerouting now reroutes around the $h$-$i$-$d$ angle by following (the outside of) the $(j',i)$ subpath and then $(i,\ell)$ subpath. Finally, when it reaches the subpath $(u,j)$ of the original path $S$, it follows $S$ until reaching $u$.

        For two paths $S_1$ and $S_2$ rerouted this way, the path which crosses the $(g,b)$ subpath closer to $g$ is rerouted closer to the triangle. Formally, let $\ell_1$ and $\ell_2$ denote the respective crossings with the $(g,b)$ subpath. If $\ell_1$ is closer to $g$ than $\ell_2$ on the $(g,b)$ subpath, then $S_1$ is rerouted closer to the triangle; otherwise, $S_2$ is rerouted closer to the triangle. All these paths are rerouted further from the triangle than the rerouting of the $(a,b)$ path.

        Similarly to the type R1 paths, observe that the ordering of the crossings $\ell(S)$ of $S$ with $g\to d$ are the exact same as $x(S)$ on $t_4\to t_3$.        
    \end{itemize}
    \item The third part of the rerouting defines the subpath of the new path ending at $v$ and after the path leaves region $A$. Let $y(S)$ (or just $y$ when $S$ is clear from context) denote the location where $S$ crosses the $(t_2,t_3)$ subpath of the boundary of $A$. The rerouting of $S$ starts at $y(S)$ and ends at $v$, following the original path $S$. This same third subpath of the rerouting applies to both type R1 and R2 paths.
    \item The second part of the rerouting defines the subpath of the new path within the region $A$. We need to define a rerouting for all paths so that for each $S$, we have that the corresponding endpoints $x(S)$ and $y(S)$ are connected. We will show that there exists some rerouting within $A$ connecting each $x(S)$ with $y(S)$ so that for any two paths $S$ and $T$, we have that the reroutings of $S$ and $T$ within $A$ cross if and only if $S$ and $T$ cross within the $h$-$d$-$g$-$b$ region.

    Recall that the ordering of the $x(S)$ on $t_1\to t_2$ corresponds exactly to the ordering of $\ell(S)$ on $g\to h$ for R1 paths. Similarly, the ordering of $x(S)$ on $t_4\to t_3$ corresponds exactly to the ordering of $\ell(S)$ on $g\to d$ for R2 paths. Finally, the ordering of $y(S)$ on $t_2\to t_3$ corresponds exactly to the ordering of the crossings of $S$ on the $h\to i$ subpat. Since our goal is to give a routing within $A$ which preserves the crossing pattern between paths within the $h$-$d$-$g$-$b$ region, we can take the exact drawings of the paths connecting pairs of portals in $B$ and use them in $A$. This is possible because the orderings of $x(S)$ and $y(S)$ on $A$ corresponds exactly to those on the $h$-$d$-$g$-$b$ region. This gives our rerouting within $A$.
    %For the ease of language, given a path $S$ which crosses the triangle, consider the points where the $S$ crosses either $(h,i)$, $(h,d)$, $(d,g)$, or $(g,b)$; call these crossing points ``portals'' for the $h$-$b$-$g$-$d$ region. Let $B$ denote the $h$-$b$-$g$-$d$ region. Observe that for each $S$, there are exactly two portals for $B$. Furthermore, we claim that for each $S$, the portals on $B$ are (in some sense) isomorphic to the ones on $A$. Formally, consider the portals on subpath $(h,g)$ and portals on subpath $(t_1,t_2)$. There is a one-to-one correspondence between portals on these two subpaths which furthermore preserves the ordering of the paths $S$ which correspond to the portals. The same holds true for subpath $(h,i)$ with $(t_2,t_3)$ and subpath $(g,b)$ with $(t_3,t_4)$. This follows immediately by the choice of ordering in O1--O4.
    %Since the portals in $B$ correspond to the ones in $A$, this implies that there is some drawing of paths connecting portals in $A$ so that the crossings are preserved exactly. Specifically, we can take the exact drawings of the paths connecting pairs of portals in $B$ and use them in $A$. This gives our rerouting in $A$.
\end{enumerate}

%The first part of the rerouting is defined below in three different cases. Our description of the rerouting will be a full rerouting of the path $S$, but we will restrict the definition to only the part of the path outside the region $A$. As mentioned before, we will define the rerouting within the region $A$ separately later.
%\begin{itemize}
    %\item[R1.] Let $S=(u,v)$ be a critical path crossing $(h,d)$ and $(h,i)$; we will reroute it around the angle $i$-$h$-$d$. Let $\ell$ be the crossing between $S$ and $(h,i)$ and let $j$ be the crossing between $S$ and $(h,d)$. The rerouting of $S$ starts at $u$ and follows $(u,\ell)$ until (almost) reaching $\ell$. Then it goes around the angle $i$-$h$-$d$ by following the outside of $(\ell,h)$ and $(h,j)$. Finally, when it reaches the subpath $(j,v)$ of the original path $S$, it follows $S$ until reaching $v$. 
    %\item[R2.] Let $S=(u,v)$ be a critical path crossing $(h,i)$, $(i,d)$, and $(d,g)$; we will again reroute it around the angle $i$-$h$-$d$. Let $\ell$ be the crossing between $S$ and $(h,i)$, let $j$ be the crossing between $S$ and $(i,d)$, and let $k$ be the crossing between $S$ and $(d,g)$. The rerouting of $S$ starts at $u$ and follows $(u,\ell)$ until (almost) reaching $\ell$. Then it goes around the angle $i$-$h$-$d$ by following the outside of $(\ell,h)$ and $(h,k)$. Finally, when it reaches the subpath $(k,v)$ of the original path $S$, it follows $S$ until reaching $v$.
    %\item[R3.] Let $S=(u,v)$ be a critical path crossing $(h,i)$ and $(i,d)$, but not $(d,g)$; we will reroute it around the $h$-$i$-$d$ angle, but will need to make additional modifications to the path (all within the $(i,d,g,b)$ region). Let $\ell$ denote the crossing between $S$ and $(h,i)$ and let $j$ denote the crossing between $S$ and $(i,d)$. Before defining the rerouting, we define some language for talking about critical paths of type R3. For two critical paths $S_1=(u_1,v_1)$ and $S_2=(u_2,v_2)$ of type R3 which cross in the region $(i,d,g,b)$, we say that $S_2$ crosses $S_1$ from below if $v_2$ lies closer than $v_1$ to $b$ when looking at the arc $(b,g)$ of the face. Now, we can define the rerouting.
    
    %The rerouting of $S$ starts at $v$ and follows the path $S$ until another path $S_1$ crosses $S$ from below. The rerouting then follows $S_1$ instead (slightly above $S_1$ without ever crossing it) until another path $S_2$ crosses $S_1$ from below. This process repeats until the rerouting reaches the boundary of the triangle at the $(i,d)$ subpath (again, without touching it). Next, let $j'$ denote the location where the rerouting reaches the $(i,d)$ subpat. The rerouting now reroutes around the $h$-$i$-$d$ angle by following (the outside of) the $(j',i)$ subpath and then $(i,\ell)$ subpat. Finally, when it reaches the subpath $(u,j)$ of the original path $S$, it follows $S$ until reaching $u$.
%\end{itemize}
%To be formal about the rerouting process, we also need to define the relative closeness to the triangle for the rerouted paths which are defined to ``go along'' the boundary of the triangle. For example, for two paths $S_1$ and $S_2$ of type R1 which are both rerouted around angle $h$, we need to decide which one is rerouted closer to $(i,h)$ and $(h,d)$. This relative ordering is defined below:
%\begin{enumerate}
    %\item[O1.] $P=(a,b)$ is rerouted closer to the triangle than all other rerouted paths.
    %\item[O2.] For two paths $S_1$ and $S_2$ rerouted around the $i$-$h$-$d$ angle (of types R1 or R2), the path which crosses the $(h,i)$ segment closer to $i$ is rerouted closer to the triangle. Formally, let $\ell_1$ and $\ell_2$ denote the respective crossings with the $(h,i)$ segment. If $\ell_1$ is closer to $i$ on the $(h,i)$ segment, then $S_1$ is rerouted closer to the triangle; otherwise, $S_2$ is rerouted closer to the triangle.
    %\item[O3.] For two paths $S_1$ and $S_2$ rerouted around the $h$-$i$-$d$ angle (of type R3), the path which crosses the $(g,b)$ segment closer to $g$ is rerouted closer to the triangle. Formally, let $v_1$ and $v_2$ denote the respective crossings with the $(g,b)$ segment. If $u_1$ is closer to $g$ on the $(g,b)$ segment, then $S_1$ is rerouted closer to the triangle; otherwise, $S_2$ is rerouted closer to the triangle.
    %\item[O4.] For two paths, $S_1$ which is rerouted around the $i$-$h$-$d$ angle and $S_2$ which is rerouted around the $h$-$i$-$d$ angle, we have that $S_1$ is always rerouted closer to the triangle than $S_2$.
%\end{enumerate}
%The above relations suffice to define a total ordering on all rerouted paths. Next, we define the rerouting within the region $A$. 



% With this definition of rerouting in $A$, we have completed the description of the full rerouting of all paths $S$ which cross the triangle. It is easy to see the following:

% \begin{observation}
%     Let $\mathcal{S}$ denote the set of critical paths which cross the triangle and let $\mathcal{T}$ denote their rerouted counterparts. The planar graph induced by the crossings of $\mathcal{S}$ and $\mathcal{T}$ are isomorphic.\label{obs:isomorphic}
% \end{observation}

% Next, we begin proving properties of the iterative process.

% \begin{observation}
%     The number of intersections between pairs of canonical paths are preserved modulo 2.
%     \label{cl:monotone and modulo 2}
% \end{observation}
% \begin{proof}
%     For the second claim, let us fix two canonical paths $P$ and $Q$ before the iteration occurred, and let $P'$ and $Q'$ denote the corresponding paths after the iteration. We claim that the area $A$ enclosed by $P$ and $P'$ doesn't include any faces or any other terminals, and similarly for the area $B$ enclosed by $Q$ and $Q'$. Indeed, the region $A$ is always a subset of the region enclosed by the old $t$-$t'$ critical path and the new $t$-$t'$ critical path, which doesn't contain any faces or other terminals. A similar observation proves the same claim for region $B$. Since the number of times $Q$ crosses the boundary of region $A$ must be even since the endpoints of $Q$ lie outside of $A$, the number of intersections between pairs $(P,Q)$ and $(P',Q)$ is the same, modulo 2. Similarly, the number of times $P'$ crosses the boundary of region $B$ is even, so the number of intersections between the pairs $(P',Q)$ and $(P',Q')$ is the same, modulo 2. Thus, the number of intersections between pairs $(P,Q)$ and $(P',Q')$ is the same, modulo 2, as desired. 
% \end{proof}

% Finally, we show that this rerouting satisfies the desired properties of decreasing the number of bad pairs.

% \subsection{Constructing the graph $H^*$: analysis of the algorithm}


% \begin{observation}\label{obs:critical-canonical}
%     Suppose there is some canonical-canonical bad pair $(P_0,Q_0)$. Then we can find a critical-canonical bad pair $(P_1,Q_1)$ as well.
% \end{observation}
% \begin{proof}
%     There are two cases depending on the exact number of times $P_0$ and $Q_0$ cross in the drawing. Suppose $P_0$ and $Q_0$ cross three or more times\footnote{In fact, it can be shown that two canonical paths satisfying properties (P1)--(P3) can cross at most three times.}. Since $P_0$ consists of the concatenation of two subpaths of critical paths, $Q_0$ must cross one of the two subpaths twice. Let $P_1$ denote the critical path where said subpath comes from and let $Q_1=Q_0$; we have found a critical-canonical bad pair $(P_1,Q_1)$. 

%     Next, suppose that $P_0$ and $Q_0$ cross exactly twice. If any one subpath of either canonical path crosses the other canonical path twice, then we have a critical-canonical bad pair immediately. If not, then each subpath must cross exactly one other subpath from the other canonical path. There are two cases, depending on the number of faces the endpoints of $P_0$ and $Q_0$ lie on. 
    
%     Suppose that the endpoints of $P_0$ and $Q_0$ lie on more than two faces. Let $P_0=u\to v$ with bend $w$, and let $u\to v'$ denote one of the critical paths which form $P_0$. Assume (without loss of generality) that there are no other endpoints of $P_0$ or $Q_0$ which lie on the same face as $v$. Observe that the $(w,v)$ and $(w,v')$ segments are ``homotopic'' since there are no faces in the region between them. Since $(w,v)$ crosses $Q_0$, this implies that $(w,v')$ crosses $Q_0$ as well. But we know that $(u,w)$ crosses $Q_0$ as well, so the critical path $u\to v'$ crosses $Q_0$ twice, so we have found our desired pair $(P_1,Q_1)$ with $P_1=(u,v')$ and $Q_1=Q_0$. This is illustrated in \textcolor{red}{Figure 1b of Planar Emulators Final.}
    
%     If the endpoints of $P_0$ and $Q_0$ lie on the same two faces, there are two cases illustrated in \textcolor{red}{Figure 1c and 1d of Planar Emulators 11}. The case in \textcolor{red}{Figure 1d} is not possible, since it contradicts the property of the canonical paths stated in Observation \ref{obs:orientations}, as one of the two canonical paths must not satisfy the property. For the case in \textcolor{red}{Figure 1d}, observe that the critical path $(u,v')$ must cross $Q$ twice, so we can set $P_0=(u,v')$ and $Q_0=Q$ as the critical-canonical bad pair.
% \end{proof}


% \begin{claim}
%     A critical-critical pair of paths cross after the iteration if and only if they cross before.
%     \label{cl:crit-crit}
% \end{claim}
% \begin{proof}
%     Let $P$ and $Q$ denote the two critical paths. First, assume that neither $P$ nor $Q$ are one of the three critical paths $(a,b)$, $(f,e)$, of $(c,g)$; we will treat those cases separately. We have three main cases, each with some subcases, depending on whether or not $P$ and/or $Q$ are rerouted.

%     \textbf{Case 1: Neither $P$ nor $Q$ is rerouted.} This case is trivial, as $P$ and $Q$ don't change, so the number of crossings is clearly preserved.

%     \textbf{Case 2: One of $P$ or $Q$ is rerouted.} Without loss of generality, let us assume $P$ is rerouted. \begin{itemize}
%         \item[2a.] If $P$ is of type R1, then $P$ must cross the $(h,g)$ segment. There are two additional subcases based on whether $P$ crosses $(h,d)$ or $(d,g)$. 
        
%         If $P$ crosses $(h,d)$, then it is rerouted along the bottom left corner $i$ of the triangle and the area between $P$ and its rerouting is a subset of the triangle. Since $Q$ is not rerouted, we know that it does not cross the triangle. Thus, if $Q$ crossed $P$ before the iteration, the crossing must lie outside the triangle. These crossings are preserved since the segments of $P$ which lie outside the triangle are not modified. Furthermore, if $Q$ didn't cross $P$ before the iteration, then we claim that no new crossings are made by the rerouting of $P$. This follows because the only new segment in $P$ after rerouting occurs along the boundary of the triangle $(h,i,d)$ which $Q$ doesn't cross by assumption. Thus, the number of crossings between $P$ and $Q$ must be preserved.

%         If $P$ crosses $(d,g)$, then it is rerouted around the bottom left corner but the area between the original path and rerouting is no longer a subset of the triangle. This case is illustrated in \textcolor{red}{Figure 3b of Planar Emulators Final}. We wish to argue that $P$ crosses $Q$ if and only if the rerouting of $P$ crosses $Q$. To do this, it suffices to show that $Q$ crosses $(k,h)$ or $(h,\ell)$ if and only if $Q$ crosses $(\ell,k)$; this is the only part of $P$ which is different after rerouting. Since $Q$ doesn't cross the triangle, it suffices to show that $Q$ crosses $(d,k)$ if and only if it crosses $(j,k)$. But this is clear because $Q$ must cross the boundary of the $(d,k,j)$ region an even number of times, $Q$ doesn't cross $(d,j)$ because it doesn't cross the triangle, and furthermore $Q$ can only cross $(d,k)$ or $(k,j)$ at most once since they are both critical paths. 
%         \item[2b.] If $P$ crosses $(h,i)$, $(i,d)$, but not $(d,g)$, then $P$ is of type R3 so it gets rerouted around the bottom right corner. This case is illustrated in \textcolor{red}{Figure 4a of Planar Emulators Final}. Unfortunately, we cannot proceed in the exact same way as in 2a because there are further modifications to $P$ within the $(i,d,g,b)$ region. In order to argue that $Q$ crosses the rerouting of $P$ whenever $Q$ crosses the original $P$, we need to show that $Q$ crosses $(\ell,j)$ or $(j,k)$ if and only if $Q$ crosses $(\ell,i)$, $(i,j)$, $(j,k')$, or $(k',k)$; this again is the only part of $P$ which is modified. Since we know that $Q$ doesn't cross the boundary of the triangle $(h,i,d)$, it suffices to prove that $Q$ crosses $(j,k)$ if and only if $Q$ crosses $(j,k')$ or $(k',k)$. 

%         If $Q$ doesn't enter the $(i,d,g,b)$ region, then we are done since $P$ is not modified outside this region. For the remainder of this paragraph, we assume $Q$ does cross the $(i,d,g,b)$ region. If $Q$ doesn't cross $(j,k)$, then it must cross $(k,k'')$; we need to show that $Q$ doesn't cross $(k,k')$ in this case. Observe that if $Q$ crosses $(k,k')$ in this case, it must cross $(k,k')$ at least twice (by our standard parity argument); we need to show this is not possible. If $Q$ does cross $(j,k)$, it clearly must cross $(k,k')$ at least once; we need to show that $Q$ doesn't cross $(k,k')$ three times in this case. We use the same argument for both cases: we show that $Q$ cannot cross $(k,k')$ two or more times. If $(k,k')$ was a segment of a critical path, this would be trivial. Unfortunately, due to the rerouting process within the $(i,d,g,b)$ region, $(k,k')$ can consist of multiple segments $k=k_0$-$k_1$-$\ldots$-$k_{r-1}$-$k_r=k'$.

%         We now argue that $Q$ cannot cross two of these segments. Suppose for contradiction that $Q$ crossed two of the segments, $(k_s,k_{s+1})$ and $(k_t,k_{t+1})$ with $s<t$, consecutively. Let $x$ and $y$ denote the respective crossings with $(k_s,k_{s+1})$ and $(k_t,k_{t+1})$ and let $S_s$ denote the critical path which contains the segment $(k_s,k_{s+1})$; we will show that $S_s$ crosses $Q$ at least twice in this case, giving a contradiction. Specifically, we will show that $S_s$ doesn't cross the rerouting of $P$ a second time. In particular, this implies that $S_s$ enters the region inscribed $(x,k_{s+1},\ldots,k_{t},y)$, implying that $S_s$ must cross the boundary again to leave the region. Since the boundary of this region consists of either segment $(x,y)$ from $Q$ or parts of the rerouting of $P$, we see that $S_s$ must cross $Q$ again, as desired. 
        
%         To see that $S_s$ doesn't cross the rerouting of $P$ a second time, observe for $z=s+1,\ldots,r$, we have that $S_{z-1}$ lies below $S_{z}$ after (to the left of) their crossing $k_z$; if this were not the case, then $S_z$ and $S_{z-1}$ would be two critical paths which crossed twice. Thus, $S_s$ cannot cross the rerouting of $P$ a second time, or else it would cross some $S_z$ with $z\in\{s+1,\ldots,r\}$ twice, giving a contradiction. That completes the proof.
%     \end{itemize}
%     \textbf{Case 3: Both $P$ and $Q$ are rerouted.} Since the induced subgraph on the rerouted paths is isomorphic before and after rerouting (Observation \ref{obs:isomorphic}), clearly the crossings must be preserved.
%     \iffalse
%     \begin{itemize}
%         \item[3a.] $P$ and $Q$ are of type R1 and R1, respectively. We need to consider the cases of $P$ and $Q$ crossing and not crossing in the triangle separately. This is illustrated in \textcolor{red}{Figures on Page 12 of Planar Emulators 10.} In 12a, there are no crossings on before rerouting (left). The green path has higher priority than the blue path, so there are also no crossings after the rerouting (right). In 12b, $P$ and $Q$ cross before the rerouting (left). The green path has higher priority, so it is rerouted closer to the angle, and causes a crossing after the rerouting (right).
%         \item[3b.] $P$ and $Q$ are of type R2 and R2, respectively. We need to consider the cases of $P$ and $Q$ crossing and not crossing in the $b$-$h$-$d$-$g$ region separately. This is illustrated in \textcolor{red}{Figures on Page 13 of Planar Emulators 10.} In 13a, there are no crossings on before rerouting (left). The green path has higher priority than the blue path, so there are also no crossings after the rerouting (right). In 13b, $P$ and $Q$ cross before the rerouting (left). The blue path has higher priority, so it is rerouted closer to the angle, and causes a crossing after the rerouting (right).
%         \item[3c.] $P$ and $Q$ are of type R3 and R3, respectively. We need to consider the cases of $P$ and $Q$ crossing and not crossing in the triangle separately. This is illustrated in \textcolor{red}{Figures on Page 14 of Planar Emulators 10.} In 14a, there are no crossings on before rerouting (left). The green path has higher priority than the blue path since it crosses $(g,b)$ closer to $g$, so it is rerouted closer to the $h$-$i$-$d$ angle. Furthermore, we know that the rerouting of the green path is always ``higher'' than the rerouting of the blue path, and touches $(d,i)$ closer to $d$. Hence, there are no crossings after the rerouting (right). In 14b, $P$ and $Q$ cross before the rerouting (left). The green path has higher priority since it crosses $(g,b)$ closer to $g$, so it is rerouted closer to the angle, and causes a crossing after the rerouting (right). Furthermore, there are no additional crossings because the green path is always rerouted higher than the blue path within the $(d,i,b,g)$ region.
%         \item[3d.] $P$ and $Q$ are of type R1 and R2, respectively. We need to consider the cases of $P$ and $Q$ crossing and not crossing in the $b$-$h$-$d$-$g$ region separately. This is illustrated in \textcolor{red}{Figures on Page 15 of Planar Emulators 10.} In 15a, there are no crossings on before rerouting (left). The green path has higher priority than the blue path, so there are also no crossings after the rerouting (right). In 15b, $P$ and $Q$ cross before the rerouting (left). The blue path has higher priority, so it is rerouted closer to the angle, and causes a crossing after the rerouting (right).
%         \item[3e.] $P$ and $Q$ are of type R1 and R3, respectively. We need to consider the cases of $P$ and $Q$ crossing and not crossing in the $b$-$h$-$d$-$g$ region separately. This is illustrated in \textcolor{red}{Figures on Page 16 of Planar Emulators 10.} In 16a, there are no crossings on before rerouting (left). Since the green path is on the left and is rerouted around the left angle while the blue path is on the right and is rerouted around the right angle, there are also no crossings after the rerouting (right). In 16b, $P$ and $Q$ cross before the rerouting (left). Since the green path is on the right and is rerouted around the left angle while the blue path is on the left and is rerouted around the right angle, this causes a crossing after the rerouting (right).
%         \item[3f.] $P$ and $Q$ are of type R2 and R3, respectively. We need to consider the cases of $P$ and $Q$ crossing and not crossing in the $b$-$h$-$d$-$g$ region separately. This is illustrated in \textcolor{red}{Figures on Page 17 of Planar Emulators 10.} In 17a, there are no crossings on before rerouting (left). Since the blue path is on the left and is rerouted around the left angle while the green path is on the right and is rerouted around the right angle, there are also no crossings after the rerouting (right). In 17b, $P$ and $Q$ cross before the rerouting (left). Since the blue path is on the right and is rerouted around the left angle while the green path is on the left and is rerouted around the right angle, this causes a crossing after the rerouting (right).
%     \end{itemize}\fi
    
%     \textbf{Case 4: $P$ is $(a,b)$ or $(f,e)$ or $(c,g)$.} Each of these cases follows essentially by construction. Suppose $P=(a,b)$; if $Q=(f,e)$ or $Q=(c,g)$, we clearly see that $P$ and $Q$ cross before and after the rerouting. Suppose $Q$ was any other critical path which wasn't rerouted. Then $Q$ doesn't cross the boundary of the triangle $(h,i,d)$, so the number of crossings between $Q$ and $P$ are preserved since the only modifications to $P$ are along the boundary of the triangle (the same argument in Case 2a applies here). Next, suppose $P=(f,e)$; if $Q$ is also not rerouted, then the claim is trivial. Suppose $Q$ is rerouted; then it is rerouted differently based on the four cases of $Q$ crossing (1) $(h,i)$, $(i,d)$, but not $(f,d)$, (2) $(h,i)$, $(i,d)$, and $(f,d)$ (3) $(h,i)$, $(h,d)$, but not $(d,g)$, and (4) $(h,i)$, $(h,d)$ and $(d,g)$. In the first case, it's clear from construction that $P$ and $Q$ don't cross before or after the rerouting. In the remaining cases, it's clear that $P$ and $Q$ cross exactly once both before and after the rerouting. Thus, the crossings between $P$ and $Q$ are always preserved exactly. The case of $P=(c,g)$ is similarly argued by construction of the rerouting.
% \end{proof}

% Next, we will analyze the number of crossings between any pair of critical-canonical paths. In order to simplify the analysis, we will first classify the types of canonical paths which can cross the $(h,i,d)$ triangle.

% \begin{claim} Suppose P1-P3 hold. The following hold true for canonical paths in $\hat{H}$:
%     \begin{itemize}
%         \item Two critical paths cross at most 1 time.
%         \item One critical path crosses one non-critical canonical path at most 2 times.
%         \item Two canonical paths cross at most 3 times. 
%     \end{itemize}\label{cl:maximum-number-of-crossings}
% \end{claim}
% \begin{proof}
%     The first two claims are simple. First, observe that the critical paths are the same as their corresponding shortest paths in $G$, so they can only cross either 0 or 1 times. Second, observe that a critical path $P$ is simply a shortest path in $G$ and a canonical path $Q$ is a concatenation of two shortest paths in $G$ (these shortest paths may not be between terminals). Since the number of crossings between $P$ and $Q$ is equal to the number of crossings between $P$ with each part of $Q$ and each part of $Q$ crosses $P$ at most once, we have that there are at most 2 crossings. The remainder of the proof focuses on the third claim.

%     Suppose for contradiction that two non-critical canonical paths $P$ and $Q$ crossed 4 (or more) times.  Let $P_1,P_2$ denote the two components of the canonical path $P$, and define $Q_1,Q_2$ similarly. Since $P_1,P_2$ and $Q_1,Q_2$ are components of shortest paths in $G$, each pair $(P_i,Q_j)$ can cross each other at most once. Since the pairs $(P_1,P_2)$ and $(Q_1,Q_2)$ cannot cross by construction, the 4 times $P$ and $Q$ cross must be between each pair $(P_i,Q_j)$ for $i,j\in[2]$ (i.e., each of these pairs must cross). We will show that this is not possible. 

%     First, we define some more notation. Let $u_1,u_2$ denote the endpoints of $P$ and $v_1,v_2$ denote the endpoints of $Q$. For each $i,j\in[2]$, let $F_{u_i},F_{v_j}$ denote the respective faces which terminals $u_i,v_j$ lie on. Without loss of generality, assume $u_i\in P_i$ and $v_j\in Q_j$ for $i,j\in[2]$. For each $j\in[2]$, let $Q_j'$ denote the remaining segment of the critical path which $Q_j$ is part of. Let $A_1$ and $A_2$ denote the areas enclosed by $Q_1$, $Q_2'$, $F_{v_1}$ and $Q_1'$, $Q_2$, $F_{v_2}$, respectively. Finally, let $w$ denote the common point of $P_1$ and $P_2$. 

%     Observe that $A_1$ and $A_2$ are disjoint regions, so there is at least one region $A_i$ such that $w\not\in A_i$. Without loss of generality, assume $i=1$ (i.e., $w\not\in A_1$). We will argue that $u_1$ and $u_2$ lie in the region $A_1$, which would give a contradiction since $u_1$ and $u_2$ are endpoints of a canonical path, implying that they lie on separate faces. To see that $u_1\in A_1$, recall we have shown that $P_1$ crosses both $Q_1$ and $Q_2$. Furthermore, observe that the path $P_1$ starts at $w$, which is outside of $A_1$. If the other endpoint $u_1$ was outside of $A_1$, then it must be the case that $P_1$ crosses the boundary of $A_1$ an even number of times. The only possible places where the boundary of $A_1$ can be crossed by a path are along $Q_1$ or $Q_2'$. We know that $P_1$ crosses $Q_1$ at least once. By properties of critical paths, we know that $P_1$ cannot cross $Q_1$ again. By the same properties, we know that $P_1$ cannot cross $Q_2'$ since $P_1$ crosses $Q_2$, and $Q_2$ and $Q_2'$ lie on the same critical path. Consequently, the other endpoint $u_1$ of $P_1$ must lie in $A_1$ so it must lie on face $F_{v_1}$. The same argument applies to $u_2$ as well, finishing our proof by contradiction. 
% \end{proof}

% \begin{observation}\label{obs:canonical-path-crossing-triangle}
%     Before the iteration occurs, a canonical path either doesn't cross the triangle, crosses one of the three following pairs: $(h,d)$ and $(d,i)$, $(h,i)$ and $(h,i)$, or $(h,i)$ and $(i,d)$, or crosses one of the three following quadruples: $(h,i)$, $(h,d)$, $(i,d)$, and $(h,i)$, or $(h,i)$, $(h,d)$, $(h,d)$, and $(h,i)$, or $(h,i)$, $(i,d)$, $(i,d)$, and $(h,i)$. In particular, any canonical path crossing the boundary of the $(h,i,d)$ triangle must also cross the $(h,i)$ segment.
% \end{observation}
% \begin{proof}
%     Recall that we assume that $a$ is the primary side of the $(a,b)$ critical path. The canonical path can cross the triangle either 0, 2, or 4 times; it is clear that 6 times is never possible. If the canonical path crosses twice, we will show that none of the following cases occur: a canonical path which crosses $(h,d)$ twice, $(i,d)$ twice, or $(h,d)$ and $(i,d)$. If the canonical path crosses the triangle four times, we will show that the only possibility is if it crosses $(h,i)$, then $(h,d)$, then $(d,i)$, and then $(h,i)$ again. 

%     First, we consider the case of two crossings with the triangle. Suppose for contradiction that the canonical path $u\to w\to v$ crosses $(h,d)$ and $(i,d)$. For a terminal node $t$, let $F_t$ denote the face it lies on. Also let $u\to v'$ and $v\to u'$ denote the critical paths which form the canonical path $u\to w\to v$. Suppose that at least one of $u$ or $v$ lies on a different face from $a$ and $b$. Without loss of generality, assume $F_v\neq F_a,F_b$. Then since $(w,v)$ crosses the $(d,e)$ segment, we know that $(w,v')$ also crosses $(d,e)$. Thus, we have a contradiction since $u\to v'$ crosses another critical path twice. Similar arguments prove that a canonical paths which cross $(h,d)$ twice or $(i,d)$ twice must lie on the same two faces as $a$ and $b$. 
    
%     Thus, we may assume without loss of generality that $F_a=F_u$ and $F_b=F_v$, and derive a contradiction for this case. We claim that this is still not possible for any canonical path crossing either $(h,d)$ twice, $(i,d)$ twice, or $(h,d)$ and $(i,d)$. Specifically, we show that this is a contradiction to the structure of the canonical paths in Observation \ref{obs:orientations}. This is illustrated for the three different cases in the \textcolor{red}{Figure 4b of Planar Emulators Final}; the $c\to d\to e$ and $u\to w\to v$ canonical paths are oriented differently while being between the same two faces, implying that at least one of the two canonical paths contradict the Observation \ref{obs:orientations}.

%     Next, we consider the case of a canonical path with four crossings with the triangle. There are only four possible structures of such a canonical path, illustrated in \textcolor{red}{Figure 5a of Planar Emulators Final.} We show that the first three structures cannot occur. The first structure has the same structure as the case of a canonical path crossing $(h,d)$ and $(i,d)$; the previous proof showing this is not possible can show that this structure is not possible either. The second structure would induce two canonical paths with four total crossings, which we have already shown is not possible (Claim \ref{cl:maximum-number-of-crossings}). Finally, the third structure is not possible since there is a critical path which crosses $(h,d)$ and $(i,d)$; this contradicts the minimality of the triangle we found in Claim \ref{cl:minimal-triangle}. Thus, the fourth structure is the only possible case of a canonical path crossing the triangle four times.
% \end{proof}

% \begin{claim}
%     The number of crossings in a critical-canonical pair of paths does not increase.\label{cl:canon-canon}
% \end{claim}
% \begin{proof}
%     Let $P$ and $Q$ be an arbitrary critical-canonical pair of paths. Suppose that $P$ and $Q$ crossed 1 or 2 times before the iteration. We know that the number of times $P$ and $Q$ cross is preserved modulo 2 after an iteration occurs, by Observation \ref{cl:monotone and modulo 2}. But since $P$ is critical and $Q$ is canonical, they can cross at most two times after the iteration since two critical paths cross at most once (Claim \ref{cl:crit-crit}) and canonical paths are made of two segments from critical paths. This implies that the number of crossings between $P$ and $Q$ must be preserved. Hence, we only need to consider $P$ and $Q$ which don't cross before the iteration, and we must show that they still don't cross after the iteration ends. Like in the previous proof, we will first assume that $(a,b)$, $(c,g)$ and $(e,f)$ are not used in either $P$ or $Q$, and treat these cases separately afterwards. Now, we begin the casework. 

%     \textbf{Case 1: Neither $P$ nor $Q$ is rerouted.} Like before, this case is trivial since $P$ and $Q$ are not modified, immediately implying that the number of crossings must be preserved.

%     \textbf{Case 2: $P$ is rerouted, $Q$ is not.} 
%     \begin{itemize}
%         \item[2a.] If $P$ crosses $(h,i)$ and $(h,d)$, then $P$ is of type R1 and it is rerouted along the bottom left corner $i$ of the triangle and the area between $P$ and its rerouting is a subset of the triangle. The proof in this case is the exact same as that of case 2a in Claim \ref{cl:crit-crit}; we repeat it here for completeness. Since $Q$ is not rerouted, we know that it does not cross the triangle. Thus, if $Q$ crossed $P$ before the iteration, the crossing must lie outside the triangle. These crossings are preserved since the segments of $P$ which lie outside the triangle are not modified. Furthermore, if $Q$ didn't cross $P$ before the iteration, then we claim that no new crossings are made by the rerouting of $P$. This follows because the only new segment in $P$ after rerouting occurs along the boundary of the triangle $(h,i,d)$ which $Q$ doesn't cross by assumption. Thus, the number of crossings between $P$ and $Q$ must be preserved.
%         \item[2b.] If $P$ crosses $(h,i)$, $(i,d)$, and $(d,g)$, then $P$ is of type R2 and it is rerouted around the bottom left corner. This case is illustrated in \textcolor{red}{Figure 5b of Planar Emulators Final}, and is similar to case 2b in Claim \ref{cl:crit-crit}. We need to argue that if $Q$ doesn't cross $(j,k)$, it cannot cross $(d,k)$ either. Suppose for contradiction that this did occur, and $Q$ crossed $(d,k)$ twice despite not crossing $(j,k)$ or $(j,d)$. The only way this is possible is for the bend $w$ of $Q$ to lie in the $(d,k,j)$ region and the endpoints ($u$ and $v$) of $Q$ to be outside the region (on the other side of the $(d,k)$ segment). But we have argued that this structure cannot occur for canonical paths. Indeed, if at least one of $u$ or $v$ doesn't lie on the faces $F_a$ and $F_b$, then there is a pair of critical paths which crosses twice. If both $u$ and $v$ lie on one of $F_a$ and $F_b$, then we contradict Observation \ref{obs:orientations}. 
%         \item[2c.] If $P$ crosses $(h,i)$, $(i,d)$, but not $(d,g)$, then $P$ is of type R3 so it gets rerouted around the bottom right corner. This case is illustrated in \textcolor{red}{Figure 4a of Planar Emulators Final}, and the proof, again, is similar to that of case 2c in Claim \ref{cl:crit-crit}. In order to argue that $Q$ crosses the rerouting of $P$ whenever $Q$ crosses the original $P$, we need to show that $Q$ crosses $(\ell,j)$ or $(j,k)$ if and only if $Q$ crosses $(\ell,i)$, $(i,j)$, $(j,k')$, or $(k',k)$; this again is the only part of $P$ which is modified. Since we know that $Q$ doesn't cross the boundary of the triangle $(h,i,d)$, it suffices to prove that $Q$ crosses $(j,k)$ if and only if $Q$ crosses $(j,k')$ or $(k',k)$. If $Q$ doesn't enter the $(i,d,g,b)$ region, then we are done since $P$ is not modified outside this region. Thus, we assume $Q$ does cross the $(i,d,g,b)$ region. 
        
%         If $Q$ doesn't cross $(j,k)$, then it must cross $(k,k'')$; we need to show that $Q$ doesn't cross $(k,k')$ in this case. Observe that if $Q$ crosses $(k,k')$ in this case, it must cross $(k,k')$ at least twice (by our standard parity argument); we need to show this is not possible. If $Q$ does cross $(j,k)$, it clearly must cross $(k,k')$ at least once; we need to show that $Q$ doesn't cross $(k,k')$ three times in this case. We use the same argument for both cases: we show that $Q$ cannot cross $(k,k')$ two or more times. If $(k,k')$ was a segment of a critical path, this would be trivial. Unfortunately, due to the rerouting process within the $(i,d,g,b)$ region, $(k,k')$ can consist of multiple segments $k=k_0$-$k_1$-$\ldots$-$k_{r-1}$-$k_r=k'$.

%         We now argue that $Q$ cannot cross two of these segments. Suppose for contradiction that $Q$ crossed two of the segments, $(k_s,k_{s+1})$ and $(k_t,k_{t+1})$ with $s<t$, consecutively. Let $x$ and $y$ denote the respective crossings with $(k_s,k_{s+1})$ and $(k_t,k_{t+1})$ and let $S_s$ denote the critical path which contains the segment $(k_s,k_{s+1})$; we will show that $S_s$ crosses $Q$ at least twice in this case, giving a contradiction. Specifically, we will show that $S_s$ doesn't cross the rerouting of $P$ a second time. In particular, this implies that $S_s$ enters the region inscribed $(x,k_{s+1},\ldots,k_{t},y)$, implying that $S_s$ must cross the boundary again to leave the region. Since the boundary of this region consists of either segment $(x,y)$ from $Q$ or parts of the rerouting of $P$, we see that $S_s$ must cross $Q$ again, as desired. 
        
%         To see that $S_s$ doesn't cross the rerouting of $P$ a second time, observe for $z=s+1,\ldots,r$, we have that $S_{z-1}$ lies below $S_{z}$ after (to the left of) their crossing $k_z$; if this were not the case, then $S_z$ and $S_{z-1}$ would be two critical paths which crossed twice. Thus, $S_s$ cannot cross the rerouting of $P$ a second time, or else it would cross some $S_z$ with $z\in\{s+1,\ldots,r\}$ twice, giving a contradiction. Finally, observe that $S_s$ and $Q$ must cross an odd number of times because $S_s$ and $Q$ are paths connecting two different opposite pairs of sides of the quadrilateral $(i,d,g,b)$. Since $S_s$ and $Q$ cross at least twice, they must also cross at least three times. This gives us our desired contradiction, as a critical path and a canonical path cannot cross three times. That completes the proof.
%     \end{itemize}
%     \textbf{Case 3: $Q$ is rerouted, $P$ is not.} Let $Q=(u,v)$, and let $w$ denote its bend. 
%     \begin{itemize}
%         \item[3a.] If $Q$ crosses the triangle four times, we treat this case separately here. Recall that in this case, the two critical paths forming $Q$ must both cross the $(h,i)$ segment and $Q$ must be a canonical path between faces $F_a$ and $F_b$. There are two subcases, depending on whether the bend $w$ lies within the $(h,d,g,b)$ region or not.

%         If $w$ doesn't lie in the $(h,d,g,b)$ region, then it is clear that both of the critical paths forming $Q$ cross the $(h,g)$ segment and hence are of type R1. If $w$ does lie in the $(h,d,g,b)$ region, it must actually lie in the $(i,d,g,b)$ region. In this case, the reroutings of the critical paths forming $Q$ will cross within the $A$ region. These two cases and the corresponding reroutings of $Q$ are illustrated in \textcolor{red}{Figure 8b and 8c of Planar Emulators Final}. In both subcases, we can see that the only part of the rerouted $Q$ which was not part of the original $Q$ goes along the boundary of the triangle. Since we know that the critical path $P$ isn't rerouted, it does not cross the triangle. Thus, if $P$ has no crossings with the original $Q$, it must also not have any crossings with the rerouted $Q$.
        
        
%         For the remaining cases, we may assume $Q$ only crosses the triangle twice.
%         \item[3b.] If $w$ does not lie in the $(g,d,h,b)$ region, we claim that the proof reduces to analyzing the crossings between critical-critical pairs. Indeed, if the bend doesn't lie within the region, this implies that one segment of the critical path (let's say $(u,w)$) lies completely outside the $(g,d,h,b)$ region, which implies that it doesn't change during the rerouting. Furthermore, the location of $w$ also doesn't change during the rerouting since it is also not in the $(g,d,h,b)$ region. Thus, analyzing crossings between $P$ and $Q$ reduces to analyzing crossings between $P$ and the $(w,v)$ segment, and hence reduces to the analysis between two critical paths in Claim \ref{cl:crit-crit}.
%         \item[3c.] If $w$ lies in the $(h,i,d)$ triangle and $Q$ crosses $(h,i)$ and $(h,d)$ or $Q$ crosses $(h,i)$ twice, we claim that $Q$ is only modified within a neighborhood of the $(h,i,d)$ triangle (see \textcolor{red}{Figure 7b and 8a of Planar Emulators Final}). In other words, we have that the area between $Q$ and it's rerouting is (essentially) contained within the triangle. This implies that if $P$ doesn't cross $Q$ before, it cannot cross $Q$ after the rerouting either or else $P$ would have to cross the region between $Q$ and the rerouting of $Q$, which is a subset of the triangle, a contradiction.
%         \item[3d.] If $Q$ crosses $(h,i)$, $(i,d)$, and $(d,j)$, then $Q$ will be rerouted around the $i$-$h$-$d$ angle. We wish to argue that $P$ crosses $Q$ if and only if $P$ crosses the rerouting of $Q$. To do this, it suffices to show that $P$ crosses $(k,h)$ or $(h,\ell)$ if and only if $P$ crosses $(\ell,k)$; this is the only part of $Q$ which is different after rerouting. Since $P$ doesn't cross the triangle, it suffices to show that $P$ crosses $(d,k)$ if and only if it crosses $(j,k)$. But this is clear because $P$ must cross the boundary of the $(d,k,j)$ region an even number of times, $P$ doesn't cross $(d,j)$ because it doesn't cross the triangle, and furthermore $P$ can only cross $(d,k)$ or $(k,j)$ at most once since they are both critical paths.
%         \item[3e.] If $Q$ crosses $(h,i)$, $(i,d)$, but not $(d,j)$, then $Q$ will be rerouted around the $h$-$i$-$d$ angle. This case is similar to case 2c above, and is illustrated in \textcolor{red}{Figure 9a of Planar Emulators Final}. Let $v$ denote the endpoint of $Q$ which lies on the $(g,b)$ segment, and let $v\to u'$ denote one of the critical paths forming $Q$. Now, we prove that $P$ crosses $Q$ before the rerouting if and only if they cross after the rerouting. If $P$ doesn't enter the $(i,d,g,b)$ region, then the claim is obvious since $Q$ is only modified within this region and along the boundary of the triangle. If $P$ does enter the $(i,d,g,b)$ region, it must enter from the $(d,g)$ segment and furthermore, it cannot cross the boundary of $(i,d,g,b)$ again. This is true because $P$ doesn't cross the triangle, and also doesn't cross $Q$. Thus, one endpoint of $P$ must lie on the $(g,b)$ segment.

%         In the figure, we see that the rerouting of $Q$ remains unchanged outside of a small neighborhood of the $(h,d,g,b)$ region. Thus, $P$ crosses the rerouting of $Q$ only if $P$ crosses the rerouting of $Q$ within the $(h,d,g,b)$ region. But the rerouting of $Q$ within this region is the same as the rerouting of the $v'\to u$ critical path within this region. Since $P$ doesn't cross $Q$, it is clear that $P$ also doesn't cross $v'\to u$ within $(h,d,g,b)$. We wish to show that $P$ also doesn't cross the rerouting of $v'\to u$, which would imply our desired claim. In Claim \ref{cl:crit-crit}, we showed that two critical paths cross after the rerouting if and only if they cross before the rerouting. Since the rerouting of $v'\to u$ is only modified within the $(h,d,g,b)$ region, this implies our desired claim that $P$ doesn't cross the rerouting of $v'\to u$.
%     \end{itemize}
%     \textbf{Case 4: Both $P$ and $Q$ are rerouted.} Since the induced subgraph on the rerouted paths is isomorphic before and after rerouting (Observation \ref{obs:isomorphic}), clearly the crossings must be preserved.

%     \textbf{Case 5: $(a,b)$ is used in either $P$ or $Q$.} Suppose $P=(a,b)$; we need to show that if $Q$ doesn't cross $(a,b)$, the rerouting of $Q$ doesn't cross the rerouting of $(a,b)$. Since $Q$ doesn't cross $(a,b)$, we know that $Q$ doesn't touch the triangle and also doesn't change after the iteration. Since the rerouting of $P$ goes along the boundary of the triangle, $Q$ and $P$ still don't cross after the rerouting.

%     Now, suppose $(a,b)$ is used in $Q$. Let $v$ be the other endpoint of $Q$, with $v\to u'$ being the other critical path forming $Q$, so that $Q$ is a canonical path from $a\to v$ with bend $w$. We know that $v\in F_b$ and $u'\in F_a$, by definition of canonical paths. There are two cases, depending on whether the bend $w$ lies on the $(a,h)$ or $(i,b)$ segments (recall that the triangle is minimal, so the bend $w$ cannot lie on the $(h,i)$ segment). In the former case, $Q$ is not modified in the rerouting. Furthermore, since $P$ is only modified within a neighborhood of the triangle, the number of crossings between $P$ and $Q$ does not change. In the latter case, $Q$ is rerouted ``above'' the triangle. Since we can assume $P$ doesn't cross the triangle (or else it must cross $Q$), we have that $P$ and $Q$ still don't cross after the rerouting. 
    
%     %In \textcolor{red}{Figure}, we see that in each of these cases, $Q$ is only changes within a neighborhood of the triangle. Thus, if $P$ doesn't touch the triangle, the claim is easy. If $P$ does touch the triangle, recall that we are assuming $P$ and $Q$ don't cross. This means that $w$ must either lie closer to $a$ on the $(a,b)$ segment in comparison to 

%     \textbf{Case 6: $(c,g)$ or $(e,f)$ are used in either $P$ and $Q$.} %If only one of $(c,g)$ and $(e,f)$ is used in $P$ and $Q$, then it is trivial since the subgraphs on the rerouted paths combined with $(c,g)$ (or with $(e,f)$) is still isomorphic before and after the rerouting.

%     Suppose $P=(c,g)$ (or symmetrically $P=(e,f)$). Recall that we can assume $Q$ doesn't cross $P$. If $Q$ doesn't cross the triangle, then neither $P$ nor $Q$ are modified in the rerouting, and we are done. If $Q$ does cross the triangle, either $Q$ crosses $(h,i)$ and $(i,d)$ or crosses $(h,i)$ twice. In the former case, one endpoint of $Q$ must be on the $(b,g)$ segment, implying that the rerouting of $Q$ will be ``around'' the $h$-$i$-$d$ angle and will still not cross $P$, as desired. In the latter case, the rerouting of $Q$ is rerouted under $(h,i)$ and also will still not cross $P$, finishing the proof. \textcolor{red}{See Figure 6a and 6b in Planar Emulators Final.}

%     Otherwise, assume $Q$ uses only $(e,f)$. Let $R$ denote the other critical path which forms $Q$. If $R$ doesn't cross the triangle, then $Q$ is not modified in the rerouting. Furthermore, any critical path $P$ is only modified within the $(b,h,d,g)$ area, so the number of crossings between $P$ and $Q$ are preserved. If $R$ does cross the triangle, then it must cross $(h,i)$ and $(i,d)$. If $P$ doesn't cross the triangle, then the number of crossing between $P$ and $Q$ doesn't change since $Q$ is only modified within a region around the triangle. If $P$ does cross the triangle, then there is only one case illustrated in \textcolor{red}{Figure 7a of Planar Emulators Final}, and there are still no crossings after the rerouting. The proof for the case where $Q$ uses only $(c,g)$ is symmetric.

%     Finally, suppose $Q$ uses both $(c,g)$ and $(e,f)$. Since we can assume that $P$ doesn't cross $Q$, we must have that $P$ doesn't cross the triangle and hence is not rerouted. Since $Q$ also is not rerouted, $P$ and $Q$ still do not cross after the rerouting.
% \end{proof}



% \iffalse
% \paragraph{Description of an iteration.}
% We now describe an iteration. 
% If there are no bad pairs in the current graph $\hat H$, we terminate the process and return it as $H^*$. Otherwise, pick any bad pair of canonical paths. Let $t,t'$ denote the endpoints of the critical path and let $t_1,t_1'$ denote the endpoints of the non-critical path. Let $t_1$-$x$ path and $t'_1$-$y$ path be the primary paths from $t_1$ and $t'_1$ that generate the $t_1$-$t'_1$ canonical path, and denote by $p,q$ the intersections between $t$-$t'$ critical path and the $t_1$-$t'_1$ canonical path, with $p$ being closer to $t$ than $q$.
% We modify the image of the $t$-$t'$ critical path as follows.
% Starting from $t$, we keep its image until right before it enters $p$, and then from there we divert course and go along (close, but in parallel to) the image of the $t_1$-$t'_1$ canonical path until right before me meet $q$, and finally from $q$ we follow its old image until we eventually reach $t'$. So the new image of the $t$-$t'$ critical path does not intersect the $t_1$-$t'_1$ canonical path, but they are very close from $p$ to $q$. See \Cref{fig: misplace_2} for an illustration. In a similar way, we modify the images of all other critical paths that intersect twice with the $t_1$-$t'_1$ canonical path, to divert their images to go along with the $t_1$-$t'_1$ canonical path, such that all these new images \zihan{more picture}

% Graph $\hat H$ is modified accordingly. The changes only lie in the ordering of the intersections between the $t$-$t'$ critical path and other critical paths. For the new $t$-$t'$ critical path, we reorder its intersections with other critical paths according to the order of the crossings appearing on its new image. 
% For other critical paths, we move its intersection with the $t$-$t'$ critical path to the new location where their crossing is. This completes the description of an iteration.

% \begin{figure}[h]
% \centering
% \subfigure[A misplaced $t$-$t'$ critical path (shown in the green dashed curve). The $t_1$-$t'_1$ critical path is shown in red. Their intersections are shown in pink/black boxes.]
% {\scalebox{0.1}{\includegraphics{figs/misplace_1.jpg}\label{fig: misplace_1}}}
% \hspace{0.5cm}
% \subfigure[After this iteration: The new image of the $t$-$t'$ critical path (orange) and the new location of the intersections between it and the $t'_1$-$y$ and the $t_1$-$x$ critical paths.]
% {\scalebox{0.1}{\includegraphics{figs/misplace_2.jpg}\label{fig: misplace_2}}}
% \caption{An iteration of the process: modifying the $t$-$t'$ critical path. \label{fig: misplace}}
% \end{figure}
% \fi

% \fi

Finally, we can prove our desired result about the graph output by the above process.

\begin{claim}
    After the process terminates, the graph and its critical paths satisfies Properties \ref{prop: f-face instance}, \ref{prop: canonical intersection}, \ref{prop: region}, and \ref{prop: intersect iff}.
\end{claim}
\begin{proof}
    Properties \ref{prop: f-face instance}, \ref{prop: canonical intersection}, and \ref{prop: region} hold for the final graph by inductively applying \Cref{obs:prop123}. It remains to prove that Property \ref{prop: intersect iff} also holds. Observe that the number of bad pairs in the beginning of the process is finite. By \Cref{obs:bad-pairs-decrease}, we know that the number of bad pairs decreases by at least 1 each iteration. Since the process continues as long as there is still a canonical bad pair remaining, the final graph has no canonical bad pairs. We show that this implies that there are no pairs of canonical paths $(P,Q)$ which cross more than once. 

    %First, we show that the endpoints of $P$ and $Q$ must lie on the same pair of faces. This is illustrated in \textcolor{red}{Figure 1b of Planar Emulators Final.} Suppose that the endpoints of $P$ and $Q$ lie on more than two faces. Let $P=u\to v$ with bend $w$, and let $u\to v'$ and $v\to u'$ denote the critical paths which form $P$. Let $Q=x\to y$ with bend $z$ be the other canonical path, and let $x\to y'$ and $y\to x'$ denote the critical paths which form $Q$. Since the endpoints of $P$ and $Q$ lie on at least three faces, there must be at least one endpoints of both $P$ and $Q$ such that no other endpoint lies on the same as them. Let $v$ and $y$ be these two endpoints.
    
    %Observe that the $(w,v)$ and $(w,v')$ segments are ``homotopic,'' meaning that any critical path we have already labelled crosses $(w,v)$ if and only if it crosses $(w,v')$, since there are no faces in the region between them and no other endpoints lying on the same face as $v$ and $v'$. Similarly, $(z,y)$ and $(z,y')$ are ''homotopic.'' We use this to give our contradiction.    If $(x,z)$ doesn't cross $Q$, then $(y,z)$ crosses both $(u,w)$ and $(w,v)$, implying that $(y,x')$ crosses $(u,w)$ and $(w,v')$, giving a contradiction since $(y,x')$ crosses $(u,v')$ twice. If $(y,z)$ doesn't cross $Q$, one can symmetrically prove that $(x,y')$ must cross $(u,v')$ twice. The remaining case is when both $(x,z)$ and $(y,z)$ cross $Q$ at least once. In this case, both $(x,z)$ and $(y,z)$ must cross $(u,v')$ since it either crosses $(u,w)$ or crosses $(w,v)$ which would imply that it crosses $(w,v')$. But this implies that $(y,z')$ also crosses $(u,v')$ so $(x,y')$ crosses $(u,v')$ twice, a contradiction.
    
    Suppose for contradiction that such a pair $(P,Q)$ existed. We will find a canonical bad pair $(P',Q')$, a contradiction. If one of $P$ or $Q$ was already a critical path, this would be trivial. Assume $P$ and $Q$ are both non-critical canonical paths. By Observation \ref{obs:same-faces} and the fact that canonical paths are safe 1-bend paths, we have that the endpoints of $P$ and $Q$ lie on the same two faces. Next, we find $(P',Q')$. If any one segment of either canonical path crosses the other canonical path twice, then we have a canonical bad pair immediately. If not, then each segment must cross exactly one other segment from the other canonical path. There are two cases illustrated in \textcolor{red}{Figure 1c and 1d of Planar Emulators Final}. The case in \textcolor{red}{Figure 1c} is not possible, since it contradicts the property of the canonical paths stated in Observation \ref{obs:orientations}, as one of the two canonical paths must not satisfy the property. For the case in \textcolor{red}{Figure 1d}, observe that the critical path $(u,v')$ must cross $Q$ twice, so we can set $P'=(u,v')$ and $Q'=Q$ as canonical bad pair.
\end{proof}




\iffalse
\begin{claim}
After this iteration, the $t$-$t'$ critical path is no longer misplaced, and no new misplaced critical path has been created.
\end{claim}
\begin{proof}
Essentially, the modification of the $t$-$t'$ critical path is the ``divert course along the $t_1$-$t'_1$ canonical path'', so after this iteration the $t$-$t'$ critical path and the $t_1$-$t'_1$ canonical path no longer intersect.
On the other hand, if we denote by $t_2,t'_2$ the terminals lying right next to $t,t'$ respectively but on the other side of the $t$-$t'$ path, then the $t_2$-$t'_2$ canonical path is the concatenation of the \zihan{more picture}
Therefore, the $t$-$t'$ critical path is no longer misplaced.
\end{proof}



\begin{claim}
    After an iteration, the following are true:
    \begin{itemize}
        \item the $t$-$t'$ critical path and $t_1$-$t_1'$ canonical path are no longer a bad pair. 
        \item the number of intersections between pairs of canonical paths are preserved modulo 2.
    \end{itemize} \label{cl:monotone and modulo 2}
\end{claim}
\begin{proof}
    For the first claim, observe that the modification of the $t$-$t'$ critical path is the ``divert course along the $t_1$-$t'_1$ canonical path'', so after this iteration the $t$-$t'$ critical path and the $t_1$-$t'_1$ canonical path no longer intersect. On the other hand, if we denote by $t_2,t'_2$ the terminals lying right next to $t,t'$ respectively but on the other side of the $t$-$t'$ path, then the $t_2$-$t'_2$ canonical path is the concatenation of the \zihan{more picture.} Therefore, the $t$-$t'$ critical path and $t_1$-$t_1'$ canonical path are no longer a bad pair. \textcolor{red}{This was taken from Zihan's writing before.}

    For the second claim, let us fix two canonical paths $P$ and $Q$ before the iteration occurred, and let $P'$ and $Q'$ denote the corresponding paths after the iteration. We claim that the area $A$ enclosed by $P$ and $P'$ doesn't include any terminal faces or any other terminals, and similarly for the area $B$ enclosed by $Q$ and $Q'$. Indeed, the region $A$ is always a subset of the region enclosed by the old $t$-$t'$ critical path and the new $t$-$t'$ critical path, which doesn't contain any terminal faces or other terminals. A similar observation proves the same claim for region $B$. Now, the remainder of the proof follows similarly to that of Claim \ref{cl:modulo 2}. Since the number of times $Q$ crosses the boundary of region $A$ must be even since the endpoints of $Q$ lie outside of $A$, the number of intersections between pairs $(P,Q)$ and $(P',Q)$ is the same, modulo 2. Similarly, the number of times $P'$ crosses the boundary of region $B$ is even, so the number of intersections between the pairs $(P',Q)$ and $(P',Q')$ is the same, modulo 2. Thus, the number of intersections between pairs $(P,Q)$ and $(P',Q')$ is the same, modulo 2, as desired. 
\end{proof}
\fi

% Using the above observations, we also prove Property \ref{prop: canonical-shortest-no-cross}.

% Let $(P_1,Q_1),\ldots,(P_N,Q_N)$ denote the sequence of bad pairs which are rerouted by the algorithm. Fix a critical path $P$. Let $Q_{i_1},\ldots,Q_{i_K}$ denote subsequence of 1-bend paths which $P$ is rerouted around (i.e., $P_{i_j}=P$ for each $j\in[K]$). Note that $Q_{i_j}$ may be a rerouted version of the canonical path from the original skeleton graph. Let the original drawing be denoted by $Q_{i_j}^0$. Similarly, let the original drawing of $P$ be denote by $P^0$.

% \begin{observation}
%     $P^0$ and $Q_{i_j}^0$ cross twice for each $j\in[K]$.
% \end{observation}
% \begin{proof}
    
% \end{proof}


% \begin{claim}
%     Rerouting $P^0$ around $Q_{i_1},\ldots,Q_{i_K}$ is equivalent to rerouting $P^0$ around $Q_{i_1}^0,\ldots,Q_{i_K}^0$.
% \end{claim}
% \begin{proof}
%     We prove the claim inductively for all critical paths $P$ and corresponding 1-bend paths $Q_{i_1},\ldots,Q_{i_K}$ simultaneously. The base case is easy for $j=0$ and $j=1$. Fix $j\ge 2$, and assume $P_{i_j}$ is equivalent to $P^0$ rerouted around $Q_{i_1}^0,\ldots,Q_{i_{j-1}}^0$. We want to inductively show that rerouting $P_{i_j}$ around $Q_{i_j}$ and rerouting $P_{i_j}$ around $Q_{i_j}^0$ is equivalent.

%     Suppose for contradiction that the two reroutings were not equivalent. This would imply that $Q_{i_j}^0$ was modified by a rerouting
% \end{proof}

% \george{below proof doesn't work}

% \begin{claim}
%     After one iteration, Property \ref{prop: canonical-shortest-no-cross} continues to hold.
% \end{claim}
% \begin{proof}
%     Consider any canonical path $Q$ between terminals $u$ and $v$ with bend $w$. Let $P_1$ from $u$ to $v'$ and $P_2$ from $v$ to $u'$ denote the critical paths that form $Q$. Finally, let $P$ denote the shortest path between $u$ and $v$. If $P_1$ or $P_2$ are not modified in the iteration or the modification doesn't affect $(u,w)$ and $(w,v)$, we trivially have the desired claim for $Q$. Suppose that we modify $P_1$ in the iteration, and the modification affects $(u,w)$. To show $Q$ doesn't cross $P$ after one iteration, it suffices to show that the rerouting of $P_1$ doesn't cross $P$.

%     Let the triangle around which we are rerouting be denoted by $(h,i,d)$. Let $(c,g)$ and $(e,f)$ denote the two critical paths containing $(h,d)$ and $(i,d)$, respectively. Observe that when going from $u$ to $v'$ along $P$, the triangle can never lie on the right side. If this were the case, by Observation \ref{obs:same-faces}, we have that the endpoints of the $c$-$d$-$e$ and $u$-$v$-$w$ canonical paths lie on the same face. But this contradicts the orientation property of 1-bend paths from Observation \ref{obs:orientations} (see \textcolor{red}{Figure}). 
    
%     Given that the triangle can't lie on the right side of $P$, we want to show the triangle doesn't cross $P$. For example, we want to preclude the possibility illustrated in \textcolor{red}{Figure TODO}. We first make the following observation:

%     \begin{observation}
%         A secondary critical path and a canonical path never form a bad pair.
%     \end{observation}
%     \begin{proof}
        
%     \end{proof}

%     Now, we can finish the proof. Suppose for contradiction that the triangle crossed $P$. Let $P_3$ denote the secondary critical path paired with $P_1$ connecting $u$ to $v''$. By the above observation, $P_3$ is never rerouted during the iterative process. In particular, we always have the property that $P_3$ is in the clockwise direction of $P$ by definition of the pairs of critical paths. Since the triangle crosses $P$ and lies on the left side of $P$, this implies that the triangle crossing $P$ must also cross $P_3$. But we just showed that it is not possible for a canonical path to cross a secondary critical path twice, a contradiction.  
% \end{proof}
\fi


It the end, we prove the following corollary which gives bounds on the size of $H^*$ and $H$.

\begin{corollary}
\label{clm: size}
$|V(H^*)|\le O(f^2k^2)$, and therefore $|V(H)|=O(f^2k^2)$.
\end{corollary}
\begin{proof}
All non-terminal vertices in $H^*$ are crossings between pairs of critical paths, so it suffices to count the number of crossings.
We have shown that the number of critical paths from each terminal $t$ is $O(f)$, so the total number of critical paths is $O(fk)$. Since two critical paths have at most one crossing, the total number of crossings is $O(f^2k^2)$. 

Recall that $H$ is obtained by sticking graphs $H^*,H_1,\ldots,H_f$ together.
From the construction, for $1\le r\le f$, $|V(H_r)|=O(|T_r|^2)$. Therefore,
\[|V(H)|\le |V(H^*)|+\sum_{1\le r\le f}O(|T_r|^2)=O(f^2k^2)+ O\bigg(\sum_{1\le r\le f}|T_r|\bigg)^2=O(f^2k^2)+O(k^2)=O(f^2k^2).\qedhere\]
\end{proof}


\iffalse
\tianyi{I think there are other properties of $H^*$ that need to be stated here. Basically, any other path that crosses the region should also cross the original shortest path.}

\begin{claim}
    
\end{claim}

\fi







\iffalse
The structure of our emulator will generalize the construction given for the 2-face case, but the setting of edge weights will differ significantly. Instead of explicitly setting the weights of each edge, we will show that for our given emulator, there exists a set of edge weights such that the pairwise distances between terminals is preserved. \george{probably needs more}


\subsection{Defining the Emulator}

Our construction will follow a similar outline as the 2-face case. We will define a subset of the pairwise shortest paths between the terminals to be \emph{critical} paths. The set of critical paths from a vertex $u$ will again have the property that they capture the structural information of all shortest paths from $u$. We then define the emulator $H$ to be the \textcolor{red}{subgraph} of $G$ which are induced by the intersections of the critical paths. 

\begin{lemma}
    There exists a planar graph of size $O(f^2k^2)$ with the following property:
\end{lemma}

\paragraph{Constructing the Emulator.} 
We first define a generalized notion of critical paths. \george{more}

\begin{definition}
    Define critical paths.
\end{definition}

With the definition of critical paths, we can now define the emulator. \george{define the emulator.}



\paragraph{Properties of the Emulator.} We now discuss the key properties of the emulator which will be needed later in setting the edge weights. \george{much more to write}

\begin{definition}
    Define canonical paths; may need to first define primary/secondary critical paths.
\end{definition}

\begin{observation}
    Two canonical paths in the emulator cross if and only if their corresponding shortest paths in the original graph cross. \label{obs:cross}
\end{observation}
\fi

























\section{The $f$-Face Case: Setting the Edge Weights}
\label{sec: edge weight}

As a last step, we need to set the edge weights in $H^*$ to preserve the distances between all interface terminal pairs. This is done in the following lemma, which, combined with the properties of graphs $H_1,\ldots,H_f$ and our construction of $H^*$, immediately implies that $H$ preserves the distances between all pairs of terminals in $G$.

\begin{lemma}
\label{lem: f face weight}
There is an edge weight function $w: E(H^*)\to \mathbb{R^+}$, such that
\begin{itemize}
    \item for each pair $t,t'$ of terminals lying on different faces, $\dist_{H^*}(t,t')= \dist_G(t,t')$; and
    \item for each pair $t,t'$ of terminals lying on the same face, $\dist_{H^*}(t,t')\ge \dist_G(t,t')$.
\end{itemize}
\end{lemma}


The remainder of this section is dedicated to the proof of \Cref{lem: f face weight}.

\subsection{The framework}
\label{subsec: weight overview}

\iffalse
We will now show that there exists a choice of edge weights for the emulator described in the previous section such that the shortest paths between all terminals nodes are preserved exactly. In order to show this, we will write the set of constraints that the edge weights must satisfy as a system of linear inequalities and argue that this linear system is feasible. Our proof of feasibility will be based on the conditions of Farkas' Lemma, and will critically use the property of the structure of our emulator given in Observation \ref{obs:cross}. The main result of the section is the following lemma:
\begin{lemma}
    For a planar graph $G$ and its corresponding planar emulator $H$ described in the previous section, there exists a set of edge weights for $E_H$ such that
    \begin{itemize}
        \item we have $d_H(u,v)=d_G(u,v)$ for all pairs of terminals $u$ and $v$ on different faces
        \item we have $d_H(u,v)\ge d_G(u,v)$ for all pairs of terminals $u$ and $v$ on the same face
    \end{itemize}\label{lem:edge-weights}
\end{lemma}
\fi


Instead of giving clean formula for directing computing edge weights as in \cite{ChangO20}, here we adopt a different approach: we characterize the required properties of such an edge weight function by an LP, and then show its existence by proving the feasibility of the LP. %Before describing the LP, we first introduce some notions.
Such an approach is recently used in \cite{chen2025path}, and we believe it will find more applications in distance-based graph problems.

\subsubsection*{Step 1. Computing edge weights by an LP}


We now describe the LP. The variables are $\set{x_e\mid e\in E(H^*)}$, where each $x_e$ represents the length we give to the edge $e$ in $E(H^*)$. There is no objective function, as we only care about its feasibility.
%
\begin{eqnarray*}
	\mbox{(LP-$H^*$)}	\quad
&\sum_{e\in E(Q)}x_e\geq \dist_G(t,t') &\forall t,t', \forall \text{ simple path }Q\text{  from }t \text{ to  }t'\\
&\sum_{e\in E(Q)}x_e\leq \dist_G(t,t') &\forall t,t', \forall \text{ canonical path }Q\text{ from }t \text{ to to }t'\\
	&x_e\geq 0&\forall e\in E(H^*)
\end{eqnarray*}

It is easy to verify that \Cref{lem: f face weight} is true iff (LP-$H^*$) is feasible.

\iffalse
\paragraph{Linear Constraints.} First, we write the set of linear constraints which the edge weights must satisfy. For each edge $e$ in the emulator $H$, let $x_e$ be the variable denoting the weight of the edge. In addition to the non-negativity constraint $x_e\ge 0$ for all edges $e$, the weights must satisfy the following:

\begin{enumerate}
    \item For any simple path $Q$ in the emulator between terminals $u$ and $v$, the length of the path in the emulator $H$ must be greater than the shortest path distance between $u$ and $v$ in the original graph $G$:
    \begin{align}
        \sum_{e\in Q}x_e\ge d_G(u,v)\label{const:simple-paths}
    \end{align}
    We note that this constraint also applies to simple paths between terminals $u$ and $v$ on the same face. This will be important later on when we need to augment the emulator by adding a $1$-face emulator to the inside of each face.
    \item For any canonical path $P$ in the emulator connecting terminals $u$ and $v$, the length of the path in the emulator $H$ must be the shortest path distance between $u$ and $v$ in the original graph $G$:
    \begin{align}
        \sum_{e\in P}x_e\le d_G(u,v)\label{const:canonical-paths}
    \end{align}
    Combined with the constraints in (\ref{const:simple-paths}), the constraint above enforces that the shortest path distances are preserved exactly by the length of canonical paths.
\end{enumerate}
It is clear that the feasibility of the linear system is equivalent to the existence of edge weights $x_e$ satisfying the conditions in Lemma \ref{lem:edge-weights}. For the remainder of the subsection, let $Ax\le b$ with $x\ge 0$ represent the above system of linear constraints in canonical form; we will show that it is feasible.

\paragraph{Farkas' Lemma.}
To show that the above linear system is feasible, we will apply Farkas' Lemma, which gives a simpler characterization of the feasibility of the linear system. \george{one more sentence here, giving a handwavy explanation of the statement of Farkas' Lemma.}

\begin{lemma}[Farkas' Lemma]
Exactly one of the two assertions holds:
\begin{enumerate}[label=\alph*.]
    \item $Ax\le b$ has a solution with $x\ge 0$
    \item $A^Ty\ge 0$ has a solution with $b^Ty<0$ and $y\ge 0$.
\end{enumerate}\label{lem:farkas}
\end{lemma}

We will try to prove via contradiction that the latter assertion cannot hold for our linear system, implying that the system must have a feasible solution. For this, we will need the following simpler but equivalent condition for the latter assertion in Farkas' Lemma. 
\fi


\subsubsection*{Step 2. Reducing the feasibility of the LP to the non-existence of certain flows}

In order to show that (LP-$H^*$) is feasible, we use Farkas' Lemma, which takes the dual of (LP-$H^*$) and convert the distance-based problem into a flow-based problem.
For preparation, we introduce some notion on flows.

\newcommand{\pset}{{\mathscr{P}}}
\newcommand{\cost}{\mathsf{cost}}

\paragraph{Terminal flows.} 
%In order to better interpret the latter condition of Farkas' lemma for our linear system, we need to introduce some definitions for terminal flows. 
Let $L$ be a graph with $T\subseteq V(L)$, and let $\mathcal{P}$ be the collection of paths in $L$ with both endpoints lying in the terminal set $T$. A \emph{terminal flow} $F: \mathcal{P}\to \mathbb{R}^+$ assigns each path $P\in\mathcal{P}$ with a flow value $F(P)$. 
For each edge $e\in E(L)$, we define $F_e=\sum_{P\in \mathcal{P}: e\in E(P)}F(P)$ as the total amount of flow of $F$ sent through edge $e$. We say that a flow $F$ \emph{dominates} another flow $F'$ if for each edge $e\in E(L)$, $F_e\ge F'_e$. 
For each pair $t,t'\in T$, we define $F_{t,t'}=\sum_{P \text{ connects } t\text{ to }t'}F(P)$ as the total amount of flow in $F$ sent between $t,t'$. 
%We say that a pair $F,F'$ of flows are \emph{equivalent} if for each pair $t,t'\in T$, $F_{t,t'}=F'_{t,t'}$. 
We define the \emph{cost} of $F$ as $\cost(F)=\sum_{t,t'\in T}F_{t,t'}\cdot \dist_G(t,t')$.
Note that, in the remainder of this section we will consider terminal flows in either $L=H^*$ or $L=G$, but we always use the terminal distances $\dist_G(t,t')$ in $G$ to define their cost.


\iffalse
\begin{claim}
    Let $Ax\le b$ be the canonical form of the linear program above, and suppose the latter assertion in Farkas' Lemma holds. Then there must exists multi-sets $\mathcal{P},\mathcal{P}^\prime$ of paths in the emulator $H$ such that:
    \begin{enumerate}
        \item $\mathcal{P}$ contains only canonical paths (including critical paths)
        \item $\mathcal{P}^\prime$ contains only simple paths
        \item $\mathcal{P}^\prime$ uses a subset of the edges of $\mathcal{P}$ (that is, $E_H(\mathcal{P}^\prime)\subseteq E_H(\mathcal{P})$)
        \item $D(\mathcal{P})<D(\mathcal{P}^\prime)$, where $D(\mathcal{P})=\sum_{p\in \mathcal{P}: \text{p connects }t,t'}D(t,t')$
    \end{enumerate}\label{cl:farkas}
\end{claim}
\fi

The following claim is a simple corollary of Farkas' lemma. Its proof is provided in \Cref{apd: Proof of clm: feasibility or flow}.

\begin{claim}
\label{clm: feasibility or flow}
%Let $Ax\le b$ be the canonical form of the linear program above, and suppose the latter assertion in Farkas' lemma holds. Then 
Either \textnormal{(LP-$H^*$)} is feasible, or    
there exist terminal flows $F,F'$ in $H^*$, such that:
    \begin{itemize}
        \item $F$ only assigns non-zero values to canonical paths;
        \item $F$ dominates $F'$; and
        \item $\cost(F)<\cost(F')$.
    \end{itemize}\label{cl:farkas}
\end{claim}


\subsubsection*{Step 3. Proving the non-existence of certain flows}

From now on, we focus on proving that the terminal flows $F,F^\prime$ satisfying the conditions in Claim \ref{cl:farkas} cannot exist. Note that this implies the feasibility of (LP-$H^*$) (by \Cref{clm: feasibility or flow}) and then \Cref{lem: f face weight}.

We will assume for contradiction that such terminal flows exist, and eventually deduce a contradiction from them.
Note that, if, instead of having terminal flows $F,F'$ described in \Cref{clm: feasibility or flow}, we have terminal flows $\hat F,\hat F'$ in $G$, such that
\begin{properties}{P}
    \item $\hat F$ only assigns non-zero values to the shortest paths connecting terminals on different faces;\label{prop: shortest}
    \item $\hat F$ dominates $\hat F'$; and \label{prop: dominates}
    \item $\cost(\hat F)<\cost(\hat F')$, \label{prop: cost}
\end{properties}
then we get a contradiction easily. This is because, if we denote, for each edge $e\in E(G)$, by $\ell_e$ its length, then for each path $P$ in $G$ connecting a pair $t,t'$ of terminals, $\sum_{e\in E(P)}\ell_e\ge \dist_G(t,t')$, so
\[
\begin{split}
\cost(\hat F') =\sum_{t,t'\in T}\hat F'_{t,t'}\cdot \dist_G(t,t') & =\sum_{t,t'\in T}  \sum_{P \text{ connects } t\text{ to }t'}\hat F'(P)\cdot\dist_G(t,t')  \\
& \le \sum_{t,t'\in T}  \sum_{P \text{ connects } t\text{ to }t'}\hat F'(P)\cdot\sum_{e\in E(P)}\ell_e  \\
& =\sum_{e\in E(G)} \ell_e \cdot \sum_{P: e\in E(P)}\hat F'(P)  \\
& =\sum_{e\in E(G)} \ell_e \cdot \hat F'_e  \\
& \le \sum_{e\in E(G)} \ell_e \cdot \hat F_e   \quad\quad\quad\quad\quad\quad\quad\quad\quad\text{($\hat F$ dominates $\hat F'$)}
\\
& =\sum_{e\in E(G)} \ell_e \cdot \sum_{P: e\in E(P)}\hat F(P)  \\ 
& = \sum_{t,t'\in T}  \sum_{P \text{ connects } t\text{ to }t'}\hat F(P)\cdot\sum_{e\in E(P)}\ell_e  \\
& =\sum_{t,t'\in T}  \sum_{P \text{ connects } t\text{ to }t'}\hat F(P)\cdot\dist_G(t,t')  \\
& =\sum_{t,t'\in T}\hat F_{t,t'}\cdot \dist_G(t,t')
= \cost(\hat F) < \cost(\hat F'),
\end{split}
\]
where the last but one line used Property \ref{prop: shortest} that flow $\hat F$ only assigns non-zero values to the shortest paths connecting terminals on different faces (and so $\dist_G(t,t')=\sum_{e\in E(P)}\ell_e$ as $P$ is the shortest path connecting $t$ to $t'$).

Therefore, it remains to convert the given terminal flows $F,F'$ in $H^*$ into terminal flows in $G$ with Properties \ref{prop: shortest}-\ref{prop: cost}.
We first construct a flow $\hat F^*$ in $G$ as follows. For each pair $t,t'$ of terminals, we send $F_{t,t'}$ units of flow between $t,t'$ along the shortest path in $G$ connecting $t$ and $t'$. Clearly, this flow $\hat F^*$ satisfies Property \ref{prop: shortest}, and Property \ref{prop: cost} as by definition $\cost(F)=\cost(\hat F^*)$.
We need the following claim, whose proof %utilizes Wye-Delta transformations, and 
is provided later.

\begin{claim}
\label{clm: routable}
There exists a flow $\hat F'$ in $G$, s.t. $\hat F'_{t,t'}=F'_{t,t'}$ for all $t,t'\in T$, and $\hat F^*$ dominates $\hat F'$.
\end{claim}

 
Let $\hat F'$ be the flow given by \Cref{clm: routable}.
Since $\hat F'_{t,t'}=F'_{t,t'}$ for every pair $t,t'\in T$, 
\[\cost(\hat F')=\sum_{t,t'}\hat F'_{t,t'}\cdot \dist_G(t,t')=\sum_{t,t'} F'_{t,t'}\cdot \dist_G(t,t')=\cost(F')<\cost(F)=\cost(\hat F^*).\]
Since $\hat F^*$ dominates $\hat F'$, they satisfy Properties \ref{prop: shortest}-\ref{prop: cost}. From the above discussion, the existence of such a pair $(\hat F^*,\hat F')$ of flows in $G$ causes a contradiction, and therefore implies \Cref{lem: f face weight}.



%If in a process of modifying $\hat F$, the values $\set{\hat F_{t,t'}}_{t,t'\in T}$ never change, then its cost $\cost(\hat F)$ never changes, either. Similarly, if in a process of modifying $\hat F'$, the values $\set{\hat F'_{t,t'}}_{t,t'\in T}$ never change, then its cost $\cost(\hat F')$ never changes, either.


The remainder of this section is dedicated to the proof of \Cref{clm: routable}. To construct the flow $\hat{F}'$ in $G$ satisfying Properties \ref{prop: shortest}-\ref{prop: cost}, we will start with the original flows $F, F'$ in $H^*$ and morph them into flows $\hat{F}, \hat{F}'$ in $G$. Before going into details, we first explain our strategy at a high level. Imagine that we draw $H^*$ and $G$ on the same plane, so that the image of terminal vertices in $T$ are the same in $H^*$ and $G$. Then, we will go over all pairs of terminals $t, t'\in T$ and try to gradually change the drawing of the canonical path $\gamma_{t, t'}$ in $H^*$ to the drawing of their shortest path in $G$, and in the meantime we will update flows $F, F'$ in $H^*$ while maintaining Properties \ref{prop: shortest}-\ref{prop: cost}.
Eventually, $H^*$ will become the subgraph of $G$ induced by all critical paths, and then its flow $F'$ will automatically become a flow in $G$, which is the $\hat F'$ we want.

We now present the details. Our morphing process consists of two phases: graph splitting and flow morphing.

\subsection{Phase 1 in proving \Cref{clm: routable}: graph splitting}

%\paragraph{Flow packing.} 
For a pair $\gamma,\gamma'$ of canonical paths, if their endpoints are $4$ distinct terminals, then they are either vertex-disjoint or crossing at a single point. However, if $\gamma,\gamma'$ share an endpoint $t$, then it is possible that they share a subpath from $t$. Recall that our goal is to morph canonical paths $\gamma,\gamma'$ to their corresponding shortest paths $P,P'$ in $G$, and $P,P'$ may well be internally disjoint -- sharing only the endpoint $t$.
The goal of this preparation phase is to split the flows $F,F'$ and the drawing of all canonical paths in $H^*$ such that all canonical paths are edge-disjoint. 

We first prove the following lemma, whose purpose will be clear shortly.

\begin{lemma}\label{interval-packing}
    Consider two multi-sets of intervals of positive integers $\aset = \{[1, a_1], [1, a_2], \ldots, [1, a_k]\}$ and $\bset = \{[b_1, c_1], [b_2, c_2], \ldots, [b_l, c_l]\}$, such that for any integer $x$, the number of intervals in $\aset$ that contains $x$ is at least the number of intervals in $\bset$ that contains $x$. Then, there is a $\bset$-partition $\bset = \bset_1\cup \bset_2\cup \cdots \cup \bset_k$, such that for each $1\leq i\leq k$:
    \begin{itemize}
        \item all intervals in $\bset_i$ are disjoint; and
        \item the union of intervals in $\bset_i$ is contained in $[1, a_i]$.
    \end{itemize}
\end{lemma}
\begin{proof}
    This is proved by induction on the sum $a_1+a_2+\cdots+a_k$. The base case $a_1+a_2+\cdots+a_k = 1$ is trivial. In general, assume $a_1 = \max_{1\leq i\leq k}\{a_i\}$. If $a_1 > \max_{1\leq j\leq l}\{c_j\}$, then replace $a_1\leftarrow a_1-1$ and apply induction. Otherwise, without loss of generality, assume $c_1 = a_1$. Then, add $[b_1, c_1]$ to $\bset_1$ and apply induction on the instance $\aset' = \{[1, b_1-1], [1, a_2], \ldots, [1, a_k]\}, \bset' = \{[b_2, c_2], \ldots, [b_l, c_l]\}$.
\end{proof}

\paragraph{Flow packing.}
As $F$ may only use canonical paths, it admits a natural splitting: for each pair $s,t$ of terminals, it sends $F_{s,t}$ units of flow along the canonical path $\gamma_{s,t}$.
But $F'$ can use any path. To better understand its structure, we decompose it as follows.
For each $F'$-flow path $P$ (say from $s$ to $t$), we start from $s$ and follow the route of $P$. Since $H^*$ consists of critical paths, $P$ must start along some critical path $\rho$. Let $\rho[s,x]$ be the maximal subpath from $s$ shared by $\rho$ and $P$. We say $\rho[s,x]$ is the first segment of $P$, and denote it by $\beta_{s,x}$.
We then continue at $x$, find the next critical path that $P$ follows, collect another segment $\beta_{x,y}$. We repeat the same operation until we reach the endpoint $t$ of $P$, obtaining a collection $\set{\beta_{x,y}\mid x,y\in P}$ of segments.
We do the same for every other $F'$-flow path. 

%Initially, define $H\leftarrow H^*$, $\hat{F}\leftarrow F, \hat{F}'\leftarrow F'$. Next, 
For each canonical path $\gamma_{s, t}$ with bend $r$, we will compute a set $\Gamma_{s, r} = \{\beta_{x, y}\mid x, y\in \gamma_{s, r}\}$ where each $\beta_{x, y}$ is a segment obtained by decomposing $F'$, such that
\begin{itemize}
    \item each edge in $\gamma_{s,t}$ is contained in no more than $F_{s,t}$ segments in $\beta_{x, y}$; and
    \item the sets $\set{\Gamma_{s,r}\mid s,t\in T, r\in \gamma_{s,t}}$ partition all the segments obtained from decomposing $F'$.
\end{itemize}

\begin{figure}[h]
\centering
\includegraphics[scale=0.5]{figs/flow-pack2.png}
\caption{\label{fig:flow-pack2} Each canonical path covers a prefix of the critical path $\rho_{s,t}$. Since $F$ dominates $F'$, we can treat each $F$-flow path from $s$ to $t_i$ as an interval $[1, a_i]$, and each $F'$-segment on $\rho_{s, t}$ as an interval $[b_j, c_j]$. Then we can apply \Cref{interval-packing} to assign maximal flow subpaths to canonical paths.}
\end{figure}

%Initially, let all sets $\Gamma_{s, r} = \emptyset$. Building the sets $\Gamma_{s, r}$ consists of the following steps. 
Fix any terminal vertex $s$, pick a critical path $\rho_{s, t}$, and consider all canonical path $\gamma_{s, t_1}, \gamma_{s, t_2}, \ldots, \gamma_{s, t_k}$ formed by $\rho_{s, t}$ and some other critical path. So by construction of $H^*$, all vertices $t_1, t_2, \ldots, t_k$ are on the same face as $t$.
%
Let $r_i$ be the bend of $\gamma_{s, t_i}, \forall 1\leq i\leq k$, and assume $r_1, r_2, \ldots, r_k$ are indexed according to their distances from $s$ in the increasing order. %Go over all the unit flow paths in $\hat{F}'$ and 
Let $\Gamma_s$ be the set of $F'$-segments which are subpaths of $\rho_{s, t}$. By definition of $F$, we send $F_{s, t_i}$ units of flow from $s$ to $r_i$ along $\rho_{s, t}$, and the union of all these flows should dominate the flow of $\Gamma_s$. If we view each flow of $F_{s, t_i}$ from $s$ to $r_i$ as an interval, and each flow in $\Gamma_s$ also as an interval, then by applying \Cref{interval-packing} we can build the sets $\Gamma_{s, r_i}, 1\leq i\leq k$ with the desired properties. See \Cref{fig:flow-pack2} as an example.


\paragraph{Graph splitting.} 
We now describe the iterative process for splitting the canonical paths in $H^*$. Initialize $H\leftarrow H^*$, $\hat{F}\leftarrow F$.
Go over each canonical path $\gamma_{s, t}$ in $H$ with bend $r$ and perform the following operations. See \Cref{fig:flow-split1}.

\begin{figure}[h]
\centering
\includegraphics[scale=0.47]{figs/flow-split1.png}
\caption{An illustration of splitting a canonical path: Before (left) and after (right).\label{fig:flow-split1}}
\end{figure}

\begin{enumerate}[leftmargin=*]
    \item Add a new vertex $r'$ infinitesimally close to $r$ in the region enclosed by $\gamma_{s, t}$ and $\delta_{s, t}$ (the shortest $s$-$t$ path in $G$). For convenience, $r'$ will be called a \emph{copy} of $r$.

    \item Draw two new curves $(s,r')$ and $(t,r')$  infinitesimally close to the original drawing of $(s,r), (t,r)$ respectively. Denote $\gamma'_{s, t} = s\to r'\to t$.

    For any previous canonical path $\gamma_{s, t_1}$ with bend $r_1$ lying between $s, r$ on $sv$, we need to adjust the curve $(s,r')$ such that $(s,r')$ crosses the new curve $(r_1',t_1)$. See \Cref{fig:flow-split3}.
    
    For each path $\rho$ currently in $H$ which crossed $(s,r)$ at a vertex $x$ before this change to $H$, create a new intersection $x'$ between $(s,r')$ and $\rho$; this is done symmetrically for $(t,r')$. Note that $\rho$ might be created in previous iterations and not necessarily in $H^*$.
    
    \item To update the flow $\hat{F}$, originally we sent $\hat{F}_{s, t}$ units of flow from $s$ to $t$ along the path $s\to r \to t$. In the new graph, we reroute this $\hat{F}_{s, t}$ units of flow through the path $\gamma'_{s, t} = s\to r'\to t$.
\end{enumerate}
\begin{figure}[h]
\centering
\subfigure[New crossing between $(s,t')$ and $(r'_1,t'_1)$.]
{\scalebox{0.54}{\includegraphics{figs/flow-split3.png}}\label{fig:flow-split3}}
\hspace{0.7cm}
\subfigure[Rerouting $F'$: augment $\alpha$ with path $(w,w')$.]
{\scalebox{0.46}{\includegraphics{figs/flow-trace.png}}\label{fig:flow-trace}}
\caption{An illustration of the construction of $H$ and the new $F'$-routing on it.}
\end{figure} 

\iffalse
\begin{figure}[h]
    \centering
    \includegraphics[scale=0.5]{figs/flow-split3.png}
\caption{\label{fig:flow-split3}}
\end{figure}
\fi

When all iterations are finished, remove from $H$ all the original critical paths; that is $H\leftarrow H\setminus H^*$. Finally, we need to show that we can reroute $F'$, still satisfying Property \ref{prop: dominates}.

\begin{claim}
The demands $\set{F'_{s,t}}_{s,t\in T}$ can be routed in the current graph $H$ while dominated by $\hat{F}$.
\end{claim}
\begin{proof}
Recall that we have decomposed the flow $F'$ in $H^*$ into segments, now we concatenate them in the updated graph $H$. %Consider any unit flow in $F'$ from $s$ to $t$, which has been decomposed into segments $\beta_{s,x_1},\beta_{x_1,x_2},\ldots,\beta_{x_l,t}$.
%For each segment $\beta_{x_i,x_{i+1}}$, let $\Gamma_{\bar s,\bar r}$ be the set it belongs to in the Flow packing step. Recall that we have drawn a new curve $(\bar s',\bar r')$ to reflect this set. We denote the points on it corresponding to $x_i,x_{i+1}$ by $x'_i,x'_{i+1}$ respectively, and the segment $\beta_{x_i,x_{i+1}}$ is now rerouted as the subpath of $(\bar s',\bar r')$ between $x'_i,x'_{i+1}$.
%\iffalse
Consider any unit flow in $F'$ from $s$ to $t$.
We will iteratively construct a $s$-$t$ flow path $\alpha$ in $H$ using the flow sets $\Gamma_{s, *}$ constructed before. Throughout, we will maintain
    \begin{itemize}[leftmargin=*]
        \item $u$ as a varying terminal;    
        \item $x, y$ as vertices in $H^*$ which is on a critical path $\rho$ starting from $u$, such that $\beta_{x,y}$ is a segment of the unit $s$-$t$ flow; and        
        \item $r$ as a vertex on $\rho$ with $\beta_{x, y}\in \Gamma_{u, r}$; and $w$ as a vertex on $(s,r')$ in $H$, where $r'$ is a copy of $r$.
    \end{itemize} 

At the beginning, let $u, x\leftarrow s$. In an iteration, assume the flow $s$-$t$ flow travels from the segment $\beta_{x,y}$ to a segment $\beta_{y,z}$.
%$on a critical path $\rho_1$ of $u$, and then turn to another critical path $\rho_2$ from terminal $v$ at the crossing $y$ between $\rho_1, \rho_2$. Call the flow sub-path between $y, z$ as $\beta_{y, z}$. Assume $\beta_{s, t}$ travels along the critical path $\rho_2$ from $y$, and leaves $\rho_2$ until it reaches vertex $z$.
Assume we have constructed the prefix of $\alpha$ from $s$ to $u$.
\iffalse
    \begin{figure}[h]
        \centering
        \includegraphics[scale=0.5]{figs/flow-trace.png}
        \caption{\label{fig:flow-trace}}
    \end{figure}
\fi
    Assume $\beta_{y,z}\in \Gamma_{v, p}$. So $x, y$  lie between $u, r$ and $y, z$ lie between $v, p$. Then, by construction of $H$, $u$ must be connecting to a copy $r'$ of $r$, and $v$ must be connecting to a copy $p'$ of $p$. Therefore, the two paths cross, say at $w'$. We extend $\alpha$ at its current end by the path $(w,w')$.
    If $z = t$, then we can complete $\alpha$ by sending this flow to $t$ through $(w',t)$. Otherwise, we update the variable $u\leftarrow v, x\leftarrow y, y\leftarrow z, r\leftarrow p$, and move on to the next iteration. See \Cref{fig:flow-trace}.
\end{proof}



\subsection{Phase 2 in proving \Cref{clm: routable}: flow morphing}


%\zihan{first 3 paragraph reroute}

%Our goal is to modify the planar drawing of $H$ and morph it to $G$. In the meantime we will not change $\hat{F}$, and argue that the demand of $\hat{F}'$ is always routable while still dominated by $\hat{F}$ every time we morph the planar drawing of $H$. In what follows, we say a path is an canonical path as a flow path in $\hat{F}$.

%To morph the drawing of $H$ into $G$, we imagine that we draw $H$ and $G$ on the same plane so that the terminal vertices in $H$ and $G$ are mapped to the same corresponding locations in both drawings. For each pair of terminals $s,t\in T$, we want to morph the drawing of the canonical path $\gamma_{s,t}$ in $H$ to the drawing of the corresponding shortest path $\delta_{s,t}$ in $G$. We will crucially make use of the fact that there are no faces in the region enclosed by $\gamma_{s,t}$ and $\delta_{s,t}$.

Recall that in graph $H$ constructed above, all canonical paths are edge disjoint. 
In this subsection, we will gradually morph $H$ to $G$, by morphing each canonical path to their corresponding shortest path, thereby converting flows $\hat F, \hat F'$ in $H$ to flows in $G$, completing the proof of \Cref{clm: routable}.

Throughout the morphing process, $H$ will always be the union of, for each pair $s,t\in T$, a canonical path $\gamma_{s,t}$. We will always let $\hat F$ be the flow that,  for all pairs $s,t\in T$,  sends $F_{s,t}$ units through path $\gamma_{s,t}$. On the other hand, instead of maintaining the other flow $\hat F'$, we will simply ensure that, over the morphing process, the demands $\set{F'_{s,t}}_{s,t\in T}$ is always routable in the graph $H$ where for each pair $s,t\in T$, the edges of $\gamma_{s,t}$ have capacity $F_{s,t}$. Note that eventually this routable property in $G$ immediately implies \Cref{clm: routable}.



%An $\textsc{M-Area}$ is the region (not containing a face) enclosed by a shortest (sub)-path $\delta_{u,v}$ and canonical (sub)-path $\gamma_{u,v}$ such that $\delta_{u,v}$ and $\gamma_{u,v}$ don't intersect, except at $u$ and $v$. 
\paragraph{Morphing areas (M-Areas).}
For a canonical path $\gamma$ and two points $u,v\in\gamma$, denote by $\gamma_{u,v}$ the subpath of $\gamma$ between $u$ and $v$. For a shortest path $\delta$ in $G$ and two points $u,v\in\delta$, denote by $\delta_{u,v}$ the subpath of $\delta$ between $u,v$.
Let $u$ and $v$ be points lying on both $\gamma$ and $\delta$. We say that \emph{$\gamma,\delta$ form an \textsc{M-Area} between $u,v$}, iff 
\begin{itemize}
    \item $\delta_{u,v}$ and $\gamma_{u,v}$ don't cross internally; and
    \item the area enclosed by $\gamma_{u,v}$ and $\delta_{u,v}$ (i.e., the \textsc{M-Area}) does not contain any terminal faces.
\end{itemize} 
An \textsc{M-Area} is said to be \emph{minimal} if there is no other \textsc{M-Area} which is a strict subset of it. For an \textsc{M-Area}, we will identify it based on the two subpaths $\gamma_{u,v}$ and $\delta_{u,v}$ that enclose it.
%\zihan{more discussion on $u,v$ being a terminal.}

\begin{observation}\label{obs:marea-top-bottom}
    Any canonical path crossing a minimal \textsc{M-Area} must cross both $\delta_{u,v}$ and $\gamma_{u,v}$ exactly once. 
\end{observation}
\begin{proof}
    A canonical path can't cross $\gamma_{u,v}$ twice since two canonical paths can't cross twice. This holds at the beginning and continues to hold due to Observation \ref{obs:flow-path-cross-iff-cross} later on. A canonical path $\gamma'$ also can't cross $\delta_{u,v}$ twice or else there would be a smaller \textsc{M-Area} formed by $\gamma'$ and $\delta$.
    %\george{This proof is assuming no problems with $u$ and $v$ being terminals; in particular, it assumes that two canonical paths don't cross twice, which is true if we distinguish crossings and intersections of paths.}
\end{proof}
%A triangle (with respect to $\gamma_{u,v}$ and $\delta_{u,v}$) is any two canonical paths $\gamma_b$ and $\gamma_c$ which have a crossing in the \textsc{M-Area}.

We now describe the iterative morphing process, which continues iff there are still \textsc{M-Areas}.
In each iteration, consider a minimal \textsc{M-Area} enclosed by $\gamma_{u,v}$ and $\delta_{u,v}$. Let $a$ be the crossing of two canonical paths $\gamma_b$ and $\gamma_c$ which cross in the \textsc{M-Area}, and let $b$ and $c$ be the crossings between the pairs $\gamma_{u,v},\gamma_b$ and $\gamma_{u,v},\gamma_c$, respectively. We say that the crossing is \emph{closest} (to $\gamma_{u,v}$) if no other canonical paths cross either $(a,b)$ or $(a,c)$. We want to morph $\gamma_{u,v}$ to $\delta_{u,v}$ over the \textsc{M-Area}, and we do this iteratively by routing $\gamma_{u,v}$ around closest crossings.

We now describe how to find a closest crossing within the minimal \textsc{M-Area}. While there exist two canonical paths which cross in in the \textsc{M-Area}, we do the following. Start with two flow paths $\gamma_b$ and $\gamma_c$ which have a crossing $a$ in the \textsc{M-Area} and cross $\gamma_{u,v}$ at $b,c$ respectively. To find a closest crossing, repeat the following until convergence:
\begin{itemize}[leftmargin=*]
    \item If there exists a canonical path which crosses $(a,b)$, update $\gamma_c$ to be canonical path which crosses $(a,b)$ closest to $b$. Update $a$ and $c$ to be the new crossings of $\gamma_c$ with $\gamma_b$ and $\gamma_{u,v}$, respectively.
    \item If there exists a canonical path which crosses $(a,c)$, update $\gamma_b$ to be canonical path which crosses $(a,c)$ closest to $c$. Update $a$ and $b$ to be the new crossings of $\gamma_b$ with $\gamma_c$ and $\gamma_{u,v}$, respectively.
\end{itemize}
At each iteration after the first, we claim the area of the triangle $(a,b,c)$ decreases. Indeed, consider $\gamma_b$ and $\gamma_c$ in any later iteration. When we are about to update $\gamma_c$, we know there are no canonical paths crossing the current $(a,c)$ by construction from the previous iteration. Thus, the new path $\gamma_c$ must cross the current $(a,b)$ and $(b,c)$, forming a triangle which is a proper subset of the previous triangle $(a,b,c)$.
Similarly, when we are about to update $\gamma_b$, there are no canonical paths crossing the current $(a,b)$, so $\gamma_b$ crosses the current $(a,c)$ and $(c,b)$, and the area of the triangle decreases again. Since the area of the triangle strictly decreases each iteration and there is a finite number of triangles, the process converges and we find a closest crossing. See Figure \ref{fig:minimal-crossing} for an illustration.

\begin{figure}[h]
\centering
\subfigure
{\scalebox{0.32}{\includegraphics{figs/Figure_minimal_crossing_1.png}}}
\hspace{1.0cm}
\subfigure
{\scalebox{0.32}{\includegraphics{figs/Figure_minimal_crossing_2.png}}}
\caption{In the left figure, $(a,b,c)$ is the current triangle with crossing at $a$. The red signs along $(a,c)$ indicate that the new $\gamma_c$ can never cross $(a,c)$. Hence, the new $\gamma_c$ (drawn in orange) must cross $(a,b)$ and $(b,c)$, decreasing the size of the triangle $(a,b,c)$. The right figure is symmetric.}\label{fig:minimal-crossing}
\end{figure} 

Upon convergence, we either have (i) a minimal \textsc{M-Area} such that no two canonical paths cross inside the \textsc{M-Area} or (ii) a closest crossing inside the minimal \textsc{M-Area}. Let $\gamma$ denote the canonical path which $\gamma_{u,v}$ is part of; we now describe the morphing of the canonical path $\gamma$ in the two cases. In case (i), the rerouting of $\gamma$ follows $\gamma$ until reaching $u$. Then it follows along $\delta_{u,v}$ until reaching $\gamma$. Finally, it follows along $\gamma$ until reaching the other terminal. In case (ii), the rerouting of $\gamma$ follows (infinitesimally close to) $\gamma$ until reaching $b$. The rerouting then follows (infinitesimally close to) the subpath $(b,a)$ and then $(a,c)$ until (almost) hitting $\gamma$. Finally, it continues to follow (infinitesimally close to) $\gamma$ until reaching the other terminal. Both of these reroutings are illustrated in {Figure \ref{fig:reroute}}. This completes the description of an iteration.

\begin{figure}[h]
\centering
\subfigure
{\scalebox{0.343}{\includegraphics{figs/Figure_reroute_1.png}}}
\hspace{1.0cm}
\subfigure
{\scalebox{0.32}{\includegraphics{figs/Figure_reroute_2.png}}}
\caption{In both cases, the red path illustrates the morphing of $\gamma_{u,v}$.}\label{fig:reroute}
\end{figure} 

In summary, our morphing process repeats the following while there is a non-empty \textsc{M-Area}:
\begin{itemize}
    \item[1.] Find a minimal \textsc{M-Area}, determined by $\gamma_{u,v}$ and $\delta_{u,v}$.
    \item[2.] Repeat the following while there are two canonical paths crossing in the \textsc{M-Area}.
    \begin{itemize}
        \item Find a closest crossing $a$ between $\gamma_b$ and $\gamma_c$ lying in the \textsc{M-Area}.
        \item Reroute $\gamma_{u,v}$ around the closest crossing $a$ via $(a,b)$ and $(a,c)$ (see Figure \ref{fig:reroute}, right).
        \item Observe that $\gamma_{u,v}$ and $\delta_{u,v}$ still define a minimal \textsc{M-Area}.
    \end{itemize}
    \item[3.] Reroute $\gamma_{u,v}$ to $\delta_{u,v}$ (see Figure \ref{fig:reroute}, left).
\end{itemize}

We first prove some properties of the morphing process. In particular, \Cref{obs:flow-path-cross-iff-cross} implies that Observation \ref{obs:marea-top-bottom} continues to hold throughout the morphing process.
\begin{observation}\label{obs:flow-path-cross-iff-cross}
    At any iteration of the morphing process, two canonical paths cross at most once.
\end{observation}
\begin{proof}
This holds at the beginning by Property \ref{prop: intersect iff}. After an iteration, since only one canonical path $\gamma$ was modified, it suffices to prove that any other flow path $\gamma'$ crosses $\gamma$ the same number of times before and after the iteration. In case (i), Observation \ref{obs:marea-top-bottom} implies that $\gamma'$ crosses $\gamma_{u,v}$ if and only it crosses $\delta_{u,v}$. In case (ii), it suffices to prove that $\gamma'$ crosses $(a,b)$ or $(a,c)$ if and only if it crosses $(b,c)$. This is since no canonical paths cross $(a,b)$ or $(a,c)$ as $a$ is a closest crossing, and also by induction, a canonical path can never cross $(b,c)$ twice before this iteration.
\end{proof}

\begin{observation}\label{obs:flow-shortest-no-face}
    At any iteration of the morphing process, the area between any canonical path and its corresponding shortest path in $G$ contain no terminal faces.
\end{observation}
\begin{proof}
This holds at the start of the morphing process. After each iteration, the change of the area between any canonical path and its corresponding shortest path is a subset of the \textsc{M-Area} in this iteration. Since an \textsc{M-Area} doesn't contain any terminal faces, the area still contains no terminal faces.
\end{proof}

We now show that the process converges.

\begin{claim}
    The iterative morphing process converges, and eventually every canonical path is morphed to the corresponding shortest path in $G$.
\end{claim}
\begin{proof}
    %For the first property, observe that the claim holds for the canonical paths before this iteration. Since only $\gamma$ was modified, it suffices to prove that any other flow path $\gamma'$ crosses $\gamma$ after the iteration if and only if it crosses before the iteration as well. In case (1), it suffices to prove $\gamma'$ crosses $\gamma_{u,v}$ if and only it crosses $\delta_{u,v}$. This is clear since $\gamma'$ can't cross $\gamma_{u,v}$ twice since both are flow paths and $\gamma'$ can't cross $\delta_{u,v}$ twice or there would be a smaller \textsc{M-Area}. In case (2), it suffices to prove that $\gamma'$ crosses $(a,b)$ or $(a,c)$ if and only if it crosses $(b,c)$. But this is trivial as no flow paths cross $(a,b)$ or $(a,c)$ by minimality of the triangle, and a flow path can never cross $(b,c)$ twice as two flow paths cross at most once.
    %For the second property, consider the flow path $\gamma$ in $\hat{F}$ which was morphed. Consider the two cases. In case (1), the subpath $\gamma_{u,v}$ of $\gamma$ is rerouted to $\delta_{u,v}$. If the morphing of $\gamma$ crosses the corresponding shortest path while $\gamma$ doesn't, this would imply that $\delta_{u,v}$ crosses the shortest path corresponding to $\gamma$ twice by parity, a contradiction. In case (2), $\gamma$ is rerouted around the closest crossing. If the morphing of $\gamma$ crosses the corresponding shortest path while $\gamma$ doesn't, this implies that the triangle must cross the shortest path. But this is not possible since the tip of the triangle is inside the minimal \textsc{M-Area}, and the shortest path corresponding to $\gamma$ lies (weakly) outside of the minimal \textsc{M-Area}.
    %For the second property, it suffices to prove that for the rerouted flow path $\gamma$, the change in area between it and the corresponding shortest path $\delta$ is a subset of the area between $\gamma$ and $\delta$. In both cases, suppose that the area between $\gamma_{u,v}$ and $\delta_{u,v}$ is on the left side when going from $u$ to $v$; it suffices to show that the area between $\gamma$ and $\delta$ is also on the right side when going from $u$ to $v$. Suppose for contradiction that this was not the case. Then $\delta$ would cross the shortest path containing $\delta_{u,v}$ twice, since the area between $\delta$ and $\gamma$ is known to not contain a face by induction (see \textcolor{red}{Figure Nameless}), a contradiction. Note that the third property is (inductively) a corollary of the second.
    First, we prove that we successfully morph $\gamma_{u,v}$ to $\delta_{u,v}$ for each minimal \textsc{M-Area} we find. Consider the number of canonical path crossings in the \textsc{M-Area} as the potential function. After each iteration of morphing the canonical path $\gamma$ around a closest crossing, this number decreases by 1. Since the total number of such points is finite, the process converges. Next, we prove that the entire process converges. Consider $\Psi_1$ defined as the total number of crossings between canonical paths and shortest paths and $\Psi_2$ defined as the total number of canonical paths which haven't been morphed to their corresponding shortest paths. When we morph $\gamma_{u,v}$ to $\delta_{u,v}$, either $\Psi_1$ decreases by at least 1 when at least one of $u,v$ is a non-terminal or $\Psi_2$ decreases by 1 when both $u,v$ are terminals and we finish morphing the canonical path between $u$ and $v$. Since $\Psi_1$ and $\Psi_2$ are both finite at the start, the process converges.


    Upon convergence, we know that there are no more \textsc{M-Area}s. We want to show that this implies that all canonical paths are morphed to their corresponding shortest paths. Suppose for contradiction this isn't the case. Let $\gamma$ be a canonical path which hasn't been morphed to its corresponding shortest path $\delta$. Let $x$ and $y$ denote two crossings between $\delta$ and $\gamma$ which are consecutive on $\delta$. If they are also consecutive on $\gamma$, then we are done. Otherwise, let $x'$ and $y'$ be two crossings between $\gamma_{x,y}$ and $\delta$ which are consecutive on $\delta$ and $x'\neq x$ or $y'\neq y$. Observe that the number of crossings between $\gamma_{x',y'}$ and $\delta$ is strictly less than the number of intersections between $\gamma_{x,y}$ and $\delta$ because $\gamma_{x',y'}$ is a strict subpath of $\gamma_{x,y}$ and in particular, doesn't include at least one of the intersections $x$ or $y$. Hence, recursively applying this process gives a pair of points $x,y$ such that $\delta_{x,y}$ and $\gamma_{x,y}$ don't intersect. Since the area surrounded by two canonical subpaths is a subset of the area between canonical paths, which from Observation \ref{obs:flow-shortest-no-face} contains no terminal faces, this is an \textsc{M-Area}. 
\end{proof}

%By the above claim, we know that we can iteratively apply this morphing process to the flow paths $\hat{F}$ and the process converges. Upon convergence, we would have that the drawings of the canonical paths are exactly the same as the drawings of corresponding shortest paths. 
%Now, it suffices to show that $\hat{F}'$ can be suitably rerouted as well to complete the proof.

The next claim shows the routability of demands $\set{F'_{s,t}}_{s,t\in T}$, and completes the proof of \Cref{clm: routable}.

\begin{claim}
After morphing a canonical path over a crossing, the demands $\set{F'_{s,t}}_{s,t\in T}$ is still routable in the updated graph $H$ where for each pair $s,t$, edges in the canonical $\gamma_{s,t}$ have capacity $F_{s,t}$.
\end{claim}
\begin{proof}
The routability holds at the beginning. It suffices to show that it holds after each iteration.
%It suffices to prove that the capacitated graph defined by the flow $\hat{F}$ is flow equivalent before and after the morphing of a flow path in one iteration. 
In case (i), where there are no crossings in the \textsc{M-Area}, the graph defined by the flow $\hat{F}$ is exactly the same as before, so the claim holds by induction. In case (ii), we show that the updated graph and the old graph  are flow-equivalent, via Wye-Delta transformations.
% \begin{figure}[h]
% \centering
% \subfigure
% {\scalebox{0.07}{\includegraphics{figs/YD_D.jpg}}}
% \hspace{1.0cm}
% \subfigure
% {\scalebox{0.07}{\includegraphics{figs/YD_Y.jpg}}}
% \caption{Illustrating the Wye-Delta transformation. \label{fig: YD}}
% \end{figure}    
    \begin{observation}[Wye-Delta Transformation]
        The following operations on a capacitated instance $(H,T)$ with edge capacities $c(\cdot)$ preserve flow equivalence (i.e., a demand on terminals $T$ is routable in $H$ iff it is routable in the new graph $H$ after the operation)
        \begin{itemize}
            \item Wye-Delta transformation: let $x$ be a degree-three non-terminal with neighbors $u,v,w$. Remove $x$ along with its incident edges, and add edges $(u,v)$, $(u,w)$, $(v,w)$ with respective capacities $[c(x,u)+c(x,v)-c(x,w)]/2$, $[c(x,u)+c(x,w)-c(x,v)]/2$, and $[c(x,v)+c(x,w)-c(x,u)]/2$.
            \item Delta-Wye transformation: let $u,v,w$ denote endpoints of a triangle. Remove the edges of the triangle and add a new non-terminal $x$ with new edges $(x,u)$, $(x,v)$, and $(x,w)$ and respective capacities $c(u,v)+c(u,w)$, $c(v,u)+c(v,w)$, and $c(w,u)+c(w,v)$.
        \end{itemize}
    \end{observation}
\begin{figure}[h]
    \centering
    \includegraphics[width=0.8\linewidth]{figs/YD_transform.png}
    \caption{An illustration of the Delta-Wye transformation (applying on the left graph to obtain the right graph), where $a,b,c$ represent the edge capacities of $(u,v)$, $(u,w)$, and $(v,w)$, respectively. Applying Wye-Delta transformation on the right graph gives us the left graph.}
    \label{fig:enter-label}
\end{figure}
    \begin{proof}
        Let $G$ denote the original graph and $H$ denote the graph after the Wye-Delta/Delta-Wye transformation. Let $\Delta$ and $Y$ denote the induced subgraphs of $G$ and $H$ on $\{u,v,w\}$ and $\{u,v,w,x\}$, respectively. First, we reduce showing terminal-flow equivalence on the entire graph to showing terminal-flow equivalence of the $\Delta$ and $Y$ structures. Fix a demand $\mathbf{b}\in\mathbb{R}^T$ on the terminals. We want to show that $\mathbf{b}$ is routable in $G$ if and only if it is routable in $H$, assuming that $\Delta$ and $Y$ are terminal flow equivalent. Indeed, take a flow $f$ in $G$ routing the demand and let $f''$ denote the flow $f$ restricted to the edges $(u,v)$, $(v,w)$, and $(u,w)$. The flow $f-f'$ no longer completely routes $\mathbf{b}$; it routes $\mathbf{b}-\mathbf{b}'$ for some residual demand $\mathbf{b}'\in\mathbb{R}^T$, which is supported on $\{u,v,w\}$. Note that $\mathbf{b}'$ can be routed in $\Delta$, specifically via $f'$. Since $\Delta$ and $Y$ are flow equivalent, there must be a flow $f''$ in $Y$ which routes $\mathbf{b}'$. Thus, the flow $f-f'+f''$ routes $\mathbf{b}-\mathbf{b}'+\mathbf{b'}=\mathbf{b}$ in $H$. The flow $f-f'$ is supported on $H-Y$ and $f''$ is supported on $Y$, so $f-f'+f''$ has congestion $1$, as desired. The reduction in the other direction is similar.
        
        %To prove flow equivalence, we need to show that any terminal demand routable in $G$ is also routable in $H$. Observe that other than the $\Delta$ and $Y$ (induced) subgraphs, $G$ and $H$ have the same edges and edge capacities. Suppose for now that $\Delta$ and $Y$ are flow equivalent, where the terminals are $u$, $v$, and $w$. Then, we claim that for any terminal demand and any flow in $G$ routing the demand, we can route the demand in $H$ as well. For a flow path which doesn't go through the $\Delta$ or $Y$ structures, we can route the flow the same way in $G$ and $H$. For each flow path which does go through the $\Delta$ or $Y$ structures, it induces a (labeled) demand on some pair of terminals $u$, $v$, and $w$. But since $\Delta$ and $Y$ are flow equivalent, this collection of demands on pairs between $u$, $v$, and $w$ can be routed in $\Delta$ if and only if they can be routed in $Y$. Thus, it remains to show that $\Delta$ and $Y$ are flow equivalent.

        It remains to show that $Y$ and $\Delta$ are terminal-flow equivalent. Since both the $Y$ and $\Delta$ structures are Okamura-Seymour instances (planar graphs with all terminals lying on the outer boundary), the flow-cut gap is $1$~\cite{okamura1981multicommodity}. Therefore, it suffices to prove terminal-cut equivalence of $\Delta$ and $Y$ to conclude the flow equivalence claim. Consider any partition of the terminals $\{u,v,w\}$ in $\Delta$ and $Y$. We want to show that the value of the minimum cut separating the terminals in $\Delta$ and $Y$ are the same. By symmetry, it suffices to show that the minimum cut separating $u$ from $v$ and $w$ is the same in $Y$ and $\Delta$. And this is indeed the case since the minimum cut in $Y$ is $c(x,u)$ and the minimum cut in $\Delta$ is $c(u,v)+c(u,w)=[c(x,u)+c(x,v)-c(x,w)]/2+[c(x,u)+c(x,w)-c(x,v)]/2=c(x,u)$, as desired. For the Delta-Wye transformation, we can similarly argue that it suffices to show that the minimum cut separating $u$ from $v$ and $w$ are the same in $\Delta$ and $Y$. This is also the case since the minimum cut in $\Delta$ is $c(u,v)+c(u,w)$ and the minimum cut in $Y$ is $c(x,u)=c(u,v)+c(u,w)$. 
    \end{proof}

\begin{figure}[h]
\centering
\subfigure
{\scalebox{0.08}{\includegraphics{figs/YD_11.jpg}}}
\hspace{0.3cm}
\subfigure
{\scalebox{0.08}{\includegraphics{figs/YD_12.jpg}}}
\hspace{0.3cm}
\subfigure
{\scalebox{0.08}{\includegraphics{figs/YD_13.jpg}}}
\hspace{0.3cm}
\subfigure
{\scalebox{0.08}{\includegraphics{figs/YD_21.jpg}}}
\hspace{0.3cm}
\subfigure
{\scalebox{0.08}{\includegraphics{figs/YD_22.jpg}}}
\hspace{0.3cm}
\subfigure
{\scalebox{0.08}{\includegraphics{figs/YD_23.jpg}}}
\caption{In the above sequence, we show that locally the old graph $H$ flow-equivalent to the blue structure. In the below sequence, we show that the new graph $H$ after the rerouting is flow-equivalent to the blue structure. But the two blue structures are exactly the same (modulo capacity-$0$ edges), so we can conclude the old and new graphs $H$ are flow-equivalent. \label{fig: YD_2}}
\end{figure}

    
    The only difference between the old and new graphs $H$ is whether $\gamma$ crosses $\gamma_b$ and $\gamma_c$ before or after $\gamma_b$ and $\gamma_c$ cross each other. The topology of the remainder of the graph is exactly the same, so it suffices to prove flow equivalence locally.
    Let $a,b,c$ denote the capacities of edges in paths $\gamma,\gamma_b,\gamma_c$, respectively. 
    In Figure \ref{fig: YD_2}, we show that the old and new graphs $H$ are locally flow-equivalent via Wye-Delta transformations. 
\end{proof}




% \paragraph{Graph morphing.} Next we are going to modify the planar drawing of $H$ and morph it to $G$. In the meantime we will not change $\hat{F}$, and argue that $\hat{F}'$ is always routable while still dominated by $\hat{F}$ every time we morph the planar drawing of $H$.

% \begin{definition}
%     For two different paths $\gamma'_{s, t}$ and $\gamma'_{u, v}$ in $H$, they are \emph{truly intersecting}, if their drawings on the plane intersect, plus that there is a vertex in $H$ sitting at this intersection point.
% \end{definition}

% \begin{claim}\label{intersect-once}
%     Each path $\gamma'_{u, v}$ can have at most one true intersection with $\gamma'_{s, t}$.
% \end{claim}
% \begin{proof}
%     \tianyi{Need properties of $H^*$.}
% \end{proof}

% During our morphing process, when we move the drawing of some paths $\gamma'_{s, t}$, we may create new intersections with some other paths $\gamma'_{u, v}$ on the plane, but we will not add more vertices to $H$ and put them at the new intersection vertices. Therefore, we will not create new true intersections. Also, as we will see, we will not remove true intersections neither.

% Go over each pair of terminals $s, t\in T$. We will gradually modify the drawing of $\gamma'_{s, t}$ and change to the drawing of $\pi_{s, t}$ in the end. In the meantime, let $A_{s, t}$ be the planar region enclosed by $\gamma'_{s, t}$ and $\pi_{s, t}$, including the boarder $\gamma'_{s, t}$. The general idea of morphing is to reduce the number of true intersections in $A_{s, t}$ while preserving flow dominance. More specifically, we will ensure that the region of $A_{s, t}$ only shrinks monotonically.

% Suppose that any path $\gamma'_{u, v}$ that truly intersect with $\gamma'_{s, t}$ do not have any true intersections inside the region $A_{s, t}$, then we continuously transform $\gamma'_{s, t}$ to $\pi_{s, t}$, and there is no need to change the flow $\hat{F}'$ since the graph $H$ itself does not change at all. See \Cref{fig:morph1}.

% \begin{figure}[h]
% \centering
% \includegraphics[scale=0.5]{figs/morph1.png}
% \caption{\label{fig:morph1}}
% \end{figure}

% Next, suppose that there exists a path $\gamma'_{u, v}$ that truly intersect with $\gamma'_{s, t}$ at vertex $x$, plus that it has another true intersection inside the region $A_{s, t}$. Let $y$ be the true intersection on $\gamma'_{u, v}$ inside $A_{s, t}$ that is the closest to $x$ along $\gamma'_{u, v}$. By \Cref{intersect-once}, $y$ cannot be on $\gamma'_{s, t}$ as well.

% Assume $y$ is the intersection with path $\gamma'_{p, q}$, then we can move $x$ along $\gamma'_{u, v}$ to reach somewhere slightly below $y$ by morphing the drawing of $\gamma'_{s, t}$. In this case, although the planar embedding has changed, we do not need to change the flow $\hat{F}'$. So Property \ref{prop: dominates} is preserved. See \Cref{fig:morph2}.

% By repeating this process, we can gradually reduce the number of true intersections strictly inside $A_{s, t}$. So in the end, we can morph $\gamma'_{s, t}$ to be $\pi_{s, t}$ without violating Property \ref{prop: dominates}. 

% \tianyi{Look suspicious as I am not even using Wye-Delta transform, but I had some trouble finding the triangles.}

% \begin{figure}[h]
% \centering
% \includegraphics[scale=0.5]{figs/morph2.png}
% \caption{\label{fig:morph2}}
% \end{figure}

%Finally, let us assume that $\gamma'_{p, q}$ truly intersects $\gamma'_{s, t}$ at vertex $z$, and this is the case for all choices of $\gamma'_{u, v}$ that truly intersects with $\gamma'_{s, t}$. Let us find such a path $\gamma'_{u, v}$ where the area of the enclosed region by $\{\gamma'_{s, t}, \gamma'_{u, v}, \gamma'_{p, q}\}$ is minimized. \tianyi{Having some trouble with area minimization since it might not be well-defined.}


%From the above observations, it suffices to modify $\hat F$ so that it eventually satisfies Property \ref{prop: shortest}, while keeping the values $\set{\hat F_{t,t'}}_{t,t'\in T}$ unchanged and ensuring the routability of demand $D$ along the way. This is done via an iterative process as follows.

\iffalse


\fi

\iffalse
Although the flow $\hat F^*$ is constructed in one-shot in Step 3, in order to prove \Cref{clm: routable}, we view it as being constructed iteratively.
As a pre-processing step, we first transform the terminal flow $F$ in $H^*$ into a terminal flow $\hat F$ in $G$ via the natural structural correspondence between $H^*$ and $G$.

\begin{observation}
\label{obs: initial}
\Cref{clm: routable} holds for this initial flow $\hat F$.
\end{observation}
\begin{proof}
We also transform the terminal flow $F'$ in $H^*$ into a terminal flow $\hat F'$ in $G$ via the natural structural correspondence between $H^*$ and $G$.
\zihan{to complete}
\end{proof}

We think of the ultimate flow $\hat F^*$ as being obtained iteratively, starting from $\hat F$.
In each iteration, we take any pair $t,t'$ of terminals such that in the current terminal flow $\hat F$, the flow sent between $t,t'$ is using a path $P$ that is not the shortest $t$-$t'$ path in $G$. We remove all the flow sent through this path $P$ \tianyi{Maybe explain that this path $P$ is canonical in $H^*$ at the beginning?}, and instead send $\hat F(P)$ units of flow via the shortest $t$-$t'$ path in $G$.
It is clear that at the end of this process, we obtain the flow $\hat F^*$.
We will show that, after each iteration,
\Cref{clm: routable} continues to hold for the updated flow $\hat F$.

We now analyze an iteration where the flow between a pair $t,t'\in T$ through path $P$ is replaced with the flow through the $t$-$t'$ shortest path in $G$, which we denote by $Q$.
Instead of replacing $P$ with $Q$ directly,
we will think of $P$ as a curve $\gamma_P$, that is being morphed into a curve $\gamma_Q$ representing $Q$, in different stages, where in each stage it ``morphs over some line segments or some intersection'' within the current region surrounded by $\gamma_P$ and $\gamma_Q$. See \Cref{fig: morphing} for an illustration. \tianyi{Is this morphing process only for the cause of analysis? If we just want to construct $\hat{F}^*$, do we simply send $\hat{F}^(P)$ through the original shortest path in $G$?}


\paragraph{The structure of terminal flows.}
Flow $\hat F$ consists of a collection of paths $P$ connecting pairs of terminals and their assigned flow value $\hat F(P)$. We can think of each flow path $P$ as a curve $\gamma_P$ connecting its terminal endpoints. These curves intersect with each other, and the intersections between a curve $\gamma_P$ and other curves appear on $\gamma_P$ in some order. If replace each intersection with a new vertex, we get a graph $\hat G$, that we call the \emph{structure} of flow $\hat F$.

What the morphing process does to the flow structure $\hat G$ is that: it does not create any new intersections or remove any old intersections, but only modifies the order of the intersections appearing on these curves.
%
At the very beginning, the structure of the initial flow $\hat F$ is a subgraph of $G$. 
We will ensure that, in the end, the structure of the resulting flow $\hat F^*$ is the subgraph of $G$ where every curve $\gamma$ connecting $t,t'$ corresponds to the shortest path in $G$ connecting $t,t'$, so the resulting flow $\hat F^*$ is essentially a flow in $G$. But in an intermediate state of the process, the morphing process may produce a flow structure $\hat G$ that is not a subgraph of $G$. \tianyi{I think it is a bit misleading to morph the planar embedding of $\hat{G}$, because it sounds like you are changing $G$. Maybe you should write it in the language of rerouting flows on a fixed embedding of $G$.}

We now formally define flow structure $\hat G$. It is a graph, but instead of specifying its vertex and edge sets, we define it in a different way as follows. Graph $\hat G$ contains all terminals, and is the union of, for each pair $t,t'$ of terminals with $F_{t,t'}>0$, a path $\gamma_{t,t'}$ connecting $t$ and $t'$, such that these paths are mutually edge-disjoint. 

\tianyi{What is $\gamma_{t,t'}$ to begin with? Canonical paths in $H^*$? How can you make sure they are edge-disjoint? Canonical paths in $H^*$ are not edge-disjoint in general, and neither are shortest paths in $G$.} 

For each pair $\gamma_{t_1,t_2}, \gamma_{t_3,t_4}$ of paths, there is an intersection between them (that is, an internal vertex denoted by $(\gamma_{t_1,t_2}, \gamma_{t_3,t_4})$ that lies on both of them), iff the $t_1$-$t_2$ shortest path and the $t_3$-$t_4$ shortest path cross each other in $G$.
For each path $\gamma$, there is an order $\pi_\gamma$ specifying the order in which all intersections $(\gamma,\gamma')$ appear on $\gamma$.
For each path $\gamma_{t,t'}$, all edges on this path have  capacity $F_{t,t'}$.

\begin{figure}[h]
\centering
\subfigure
{\scalebox{0.06}{\includegraphics{figs/morphing_1.jpg}}}
\hspace{0.1cm}
\subfigure
{\scalebox{0.06}{\includegraphics{figs/morphing_2.jpg}}}
\hspace{0.1cm}
\subfigure
{\scalebox{0.06}{\includegraphics{figs/morphing_3.jpg}}}
\caption{ \label{fig: morphing}}
\end{figure}

The morphing process is essentially  modifying the flow structure $\hat G$, which we now describe in more detail.
In order to keep track, although $\hat G$ may not be a subgraph of $G$ (and in fact $\hat G$ is possibly not planar during the process), we still associate it with a drawing (which is not necessarily a planar drawing \tianyi{Why not planar? Isn't it always drawn on a plane?}), based on the drawing of $G$.
At the beginning, the drawing of the flow structure $\hat G$ of the initial flow $\hat f$ can be induced from $G$, as at this moment $\hat G$ is a subgraph of $G$.
In the iteration of processing the flow between terminals $t,t'$, we only morph the path $\gamma_{t,t'}$ and its corresponding curve in $\hat G$ and its drawing, until the image of $\gamma_{t,t'}$ coincides with the shortest $t$-$t'$ path in $G$. 
In each stage of this iteration, we will perform the following triangle flip operation to $\hat G$ and morph the curve $\gamma_{t,t'}$ accordingly.
\fi

\iffalse
A set of three paths $\set{\gamma_{1},\gamma_{2},\gamma_{3}}$ is said to form a \emph{triangle}, iff
\begin{itemize}
    \item every pair of them has an intersection; and
    \item on path $\gamma_{1}$, the intersections $(\gamma_{1},\gamma_{2})$ and $(\gamma_{1},\gamma_{3})$ appear on consecutively; \tianyi{Remove `on'. What do you mean by `consecutive'? No other intersection points in between?}
    \item on path $\gamma_{2}$, the intersections $(\gamma_{1},\gamma_{2})$ and $(\gamma_{2},\gamma_{3})$ appear on consecutively; and
    \item on path $\gamma_{3}$, the intersections $(\gamma_{1},\gamma_{3})$ and $(\gamma_{2},\gamma_{3})$ appear on consecutively.
\end{itemize}
The operation of \emph{flipping the triangle $\set{\gamma_{1},\gamma_{2},\gamma_{3}}$} includes 
\begin{itemize}
    \item flipping the order between intersections $(\gamma_{1},\gamma_{2})$ and $(\gamma_{1},\gamma_{3})$ on path $\gamma_{1}$; 
    \item flipping the order between intersections $(\gamma_{1},\gamma_{2})$ and $(\gamma_{2},\gamma_{3})$ on path $\gamma_{2}$; and
    \item flipping the order between intersections $(\gamma_{1},\gamma_{3})$ and $(\gamma_{2},\gamma_{3})$ on path $\gamma_{3}$.
\end{itemize}
We do not modify any other intersections. See \Cref{fig: triangle flip} for an illustration.

\begin{figure}[h]
\centering
\subfigure
{\scalebox{0.07}{\includegraphics{figs/triangle_1.jpg}}}
\hspace{1.0cm}
\subfigure
{\scalebox{0.07}{\includegraphics{figs/triangle_2.jpg}}}
\caption{An illustration of a triangle flip $\set{\gamma_1,\gamma_2,\gamma_3}$. We keep the images of $\gamma_2,\gamma_3$, and morph $\gamma_1$. \label{fig: triangle flip}}
\end{figure}    

\begin{claim}
In the iteration of processing the flow between terminals $t,t'$, there is a finite sequence of triangle flips that allow us to morph the image of $\gamma_{t,t'}$ to the image of $t$-$t'$ shortest path in $G$.
\end{claim}
\begin{proof}
    Let $\pi_{t, t'}$ be the curve of the drawing of the shortest path between $t, t'$. Consider the enclosed area $A_{t, t'}$ between the drawing of $\gamma_{t, t'}$ and the drawing of $\pi_{t, t'}$. We will iteratively perform triangle flips to shrink the size of $A_{t, t'}$ which is done in the following manner. In each iteration, let $C_{t, t'} = \{\gamma_{x, x'}\mid x, x'\in T\}$ be the set of all the curves that have intersections with $\gamma_{t, t'}$. \tianyi{Need some basic properties regarding $H^*$ in order to continue.}
\end{proof}

We are ready to complete the proof of \Cref{clm: routable}. Recall that we are given a terminal flow $F'$ in $H$, and the goal is to show a terminal flow $\hat F'$ in $G$, with $\hat F'_{t,t'}=F'_{t,t'}$ for all pairs $t,t'$.
%
We now view the values $\set{F'_{t,t'}}$ as a \emph{demand} on terminals. That is, we seek to route, for every pair $t,t'$ of terminals, $F'_{t,t'}$ units of flow between them. 
%
We will show that, such a demand, which we denote by $D$, is always routable in the flow structure graph $\hat G$, for any intermediate state of flow $\hat F$ during the iterative process. 

\Cref{obs: initial} actually shows that $D$ is routable in the flow structure $\hat G$ for the initial flow $\hat F$.
We prove the following claim, which states that in fact all the flow structure graphs we get during the morphing process are flow-equivalent. Its proof uses the Wye-Delta transformation.

\begin{claim}
\label{clm: Y-Delta}
Let $\hat G_1$ and $\hat G_2$ be flow structure graphs before and after any morphing step. Then any demand on terminals is routable in $\hat G_1$ iff it is routable in $\hat G_2$. \tianyi{How do you preserve dominance? Actually I cannot find any proof arguing about dominance.}
\end{claim}
\zihan{complete the proof by adding explanations to the figures}

\begin{theorem}[$Y$-$\Delta$ transformation]
For any demand $D$ on terminals $\set{1,2,3}$, $D$ is routable in the edge-capacitated $\Delta$-graph (left in \Cref{fig: YD}) iff $D$ is routable in the edge-capacitated $Y$-graph (right in \Cref{fig: YD}).
\end{theorem}
\begin{figure}[h]
\centering
\subfigure
{\scalebox{0.07}{\includegraphics{figs/YD_D.jpg}}}
\hspace{1.0cm}
\subfigure
{\scalebox{0.07}{\includegraphics{figs/YD_Y.jpg}}}
\caption{ \label{fig: YD}}
\end{figure}    


\begin{figure}[h]
\centering
\subfigure
{\scalebox{0.08}{\includegraphics{figs/YD_11.jpg}}}
\hspace{0.3cm}
\subfigure
{\scalebox{0.08}{\includegraphics{figs/YD_12.jpg}}}
\hspace{0.3cm}
\subfigure
{\scalebox{0.08}{\includegraphics{figs/YD_13.jpg}}}
\caption{ \label{fig: YD_1}}
\end{figure}    


\begin{figure}[h]
\centering
\subfigure
{\scalebox{0.08}{\includegraphics{figs/YD_21.jpg}}}
\hspace{0.3cm}
\subfigure
{\scalebox{0.08}{\includegraphics{figs/YD_22.jpg}}}
\hspace{0.3cm}
\subfigure
{\scalebox{0.08}{\includegraphics{figs/YD_23.jpg}}}
\caption{ \label{fig: YD_2}}
\end{figure}    
\end{proof}



As a corollary, $D$ is routable in the flow structure $\hat G^*$ for the ultimate flow $\hat F^*$. 
\zihan{complete the proof of \Cref{clm: routable}}

\fi

\appendix
%\newpage
\section{Missing Proofs in \Cref{sec: warmup}}

\subsection{Proof of \Cref{obs: split}}
\label{apd: Proof of obs: split}

Denote by $a_j$ the segment of the inner face boundary going clockwise from $t'_j$ to $t'_{j+1}$.
Recall that we say that $t_i$ splits at $(t'_j,t'_{j+1})$ if the region enclosed by path $P_{i,j}$, path $P_{i,j+1}$, and segment $a_j$ contains the inner face.

We first prove existence of the splitting pair. If $t_i$ splits at $(t'_j,t'_{j+1})$ for some $j=1,\ldots,\frac{k}{2}-1$, then we are done. Suppose not; then we will show it splits at the pair $(t'_{k/2},t_1)$. To see this, consider the regions $A_j$ inscribed by the compositions of paths $P_{i,j}$, $t'_{j}\to t'_{j+1}$, and $P_{i,j+1}$ for each $j=1,\ldots,\frac{k}{2}-1$. Denote their union by $A=\bigcup_{j=1,\ldots,k/2-1}A_j$, and note that $A$ doesn't contain the inner face. It can be seen by definition that $A$ is the region inscribed by the composition of paths $P_{i,1}$, $t'_1\to t'_2\to\ldots\to t'_{k/2-1}$, and $P_{i,k/2-1}$. As a result, the region inscribed by composition of $P_{i,1}$, $t'_1\to t'_{k/2-1}$, and $P_{i,k/2-1}$ must contain the inner face since the path consisting of the sequence of edges $t'_{j}\to t'_{j+1}$ for $j=1,\ldots,k/2$ inscribes the inner face. Thus, $t_i$ must split at $t_{k/2}$ and we are done.

We now prove uniqueness of the splitting point $(t'_j,t'_{j+1})$. Suppose for contradiction that $t_i$ splits at $(t'_j,t'_{j+1})$ and $(t'_\ell,t'_{\ell+1})$ for $\ell\neq j$. Observe that neither pair of paths $P_{i,j}$ and $P_{i,\ell+1}$ nor $P_{i,j+1}$ and $P_{i,\ell}$ can have a common node other than $t_i$; this would give a contradiction to the assumption that two shortest paths intersect only at a single point. But if neither pair had a common node, then both paths $P_{i,\ell}$ and $P_{i,\ell+1}$ would have to lie in the interior on the region enclosed by $(P_{i,j},P_{i,j+1},a_j)$. This implies that the region enclosed by $(P_{i,\ell},P_{i,\ell+1},a_\ell)$ cannot contain the inner face, a contradiction.


\subsection{Proof of \Cref{obs: split move}}
\label{apd: Proof of obs: split move}

Suppose that $t'_{j_2}$ didn't lie on the (closed) clockwise segment between $t'_{j_1}$ and $t'_{j_3}$; we will show that there must be two shortest paths which intersect twice. Observe that since $j_2$ isn't on the clockwise segment between $t'_{j_1}$ and $t'_{j_3}$, it must be on the counterclockwise segment between $t'_{j_1-1}$ and $t'_{j_3+1}$. Now, consider the critical path $P_{i_2,j_2+1}$. Since $t_{i_2}$ splits at $(t'_{j_2},t'_{j_2+1})$, this forces $P_{i_2,j_2+1}$ to intersect $P_{i_1,j_1}$ twice, giving a contradiction.


\subsection{Formal description of the emulator structure}
\label{apd: 2-face emulator}

Formally, the emulator consists of three parts: 
\begin{itemize}[leftmargin=*]
    \item  The first part is a planar graph $H_1$ that contains terminals $t_1,\ldots,t_{k/2}$ and has them lie on a single face $F_1$ (in the same order as they lie on the face in $G$), preserving the distances between them in $G$. In fact, we will simply use the construction of \cite{ChangO20} and \cite{goranci2020improved}, which gives a construction of $H_1$ on $O(k^2)$ vertices.
    \item The second part is a planar graph $H_2$ that contains terminals $t'_1,\ldots,t'_{k/2}$ and has them lie on a single face $F_2$ (in the same order as they lie on the face in $G$), preserving the distances between them in $G$. Similarly, we will simply use the construction of \cite{ChangO20} and \cite{goranci2020improved}, which gives a construction of $H_2$ on $O(k^2)$ vertices.
    \item the third part is a planar graph $H_3$ that contains all terminals, with terminals $t_1,\ldots,t_{k/2}$ lying on a single face $F_3$ (in the same order), and terminals $t'_1,\ldots,t'_{k/2}$ lying on another face $F_4$ (in the same order). $H_3$ is in charge of preserving the distances between inter-face pairs $t_i,t'_j$ in $G$.
\end{itemize}

\begin{figure}[h]
    \centering
    \subfigure%[Graphs $H_1,H_2,H_3$.]
    {\scalebox{0.1}{\includegraphics{figs/H_123.jpg}}}
    \hspace{1cm}
    \subfigure%[Graph $H$.]
    {\scalebox{0.1}{\includegraphics{figs/H.jpg}}}
    \caption{An illustration of graphs $H_1,H_2,H_3$ and graph $H$ (right) obtain by concatenating them.\label{fig: concatenation}}
\end{figure}

The final emulator $H$ is obtained by sticking graphs $H_1,H_2,H_3$ together. Specifically, we start from graph $H_3$, and (i) place the drawing of $H_1$ (with $F_1$ being the its outer face) into face $F_3$, and identify the two copies of each terminal $t_i$ in $H_1$ and $H_3$; (ii) place the drawing of $H_2$ (with $F_2$ being the its outer face) into face $F_4$, and identify the two copies of each terminal $t'_j$ in $H_2$ and $H_3$. See \Cref{fig: concatenation} for an illustration.

From the construction, we have that $|V(H)|\le |V(H_1)|+|V(H_2)|+|V(H_3)|\le O(k^2)+|V(H_3)|$. Therefore, if we show $|V(H_3)|=O(k^2)$, then we immediately obtain $|V(H)|=O(k^2)$.



\section{Proof of \Cref{clm: feasibility or flow}}
\label{apd: Proof of clm: feasibility or flow}

We use the following Farkas' lemma \cite{farkas1898fourier}.
\begin{lemma}[Farkas' Lemma]
For any matrix $A\in \mathbb{R}^{m\times n}$ and any vector $b\in \mathbb{R}^m$, exactly one of the two assertions holds:
\begin{itemize}
\item there exists a vector $x\in \mathbb{R}^n$ such that $Ax\le b$ and $x\ge 0$;
\item there exists a vector $y\in \mathbb{R}^m$ such that $A^Ty\ge 0$, $b^Ty<0$, and $y\ge 0$.
\end{itemize}\label{lem:farkas}
\end{lemma}
    
Note that (LP-$H^*$) can be written in the $Ax\le b$ form, and its feasibility is equivalent to the existence of a vector $x\ge 0$ that satisfies $Ax\le b$. Therefore, from Farkas' lemma, either (LP-$H^*$) is feasible, or there exists $y\ge 0$ satisfying $A^Ty\ge 0$ and $b^Ty<0$. 

To define the terminal flows $F$ and $F^\prime$, let us examine the vector $y$. If we view $A^Ty$ as a linear combination of the constraints represented by $A$, then the values in $y$ can be interpreted as the coefficients for the constraints. For each coefficient $y_Q$ of the first type of constraint on simple paths $Q$, define $F^\prime(Q)=y_Q$. For each coefficient $y_P$ of the second type of constraint on canonical paths $P$, define $F(P)=y_P$. Then define $F$ and $F^\prime$ on all other inputs to be zero. We will show that the flows $F,F^\prime$ satisfy the desired conditions.

    Condition (1) is obviously satisfied by how we define $F$. To see that Condition (2) holds, we will try to prove that $F_e-F^\prime_e\ge 0$ for each $e\in E(H)$. We claim that this is equivalent to the fact that the $A^Ty\ge 0$. Let us focus on one edge $e\in E(G)$. The coefficients of $x_e$ form a column of $A$; let us denote this column by $A_e$. Then $A^Ty\ge 0$ implies that $A^T_ey\ge 0$; we claim that this implies $F_e-F^\prime_e\ge 0$. Observe that the coordinate of $A_e$ corresponding to the first type of constraint for simple path $Q$ has value $-1$ if $e\in E(Q)$, and zero otherwise. Similarly, the coordinate of $A_e$ corresponding to the second type of constraint for canonical path $P$ has value $1$ if $e\in E(P)$, and zero otherwise. As a result,
    $$A_e^Ty=\sum_{P:e\in E(P)}y_P-\sum_{Q:e\in E(Q)}y_Q=\sum_{P:e\in E(P)}F(P)-\sum_{Q:e\in E(Q)}F(Q)=F_e-F^\prime_e,$$
    where the summations over $P$ are over canonical paths and summations over $Q$ are over all simple paths.

    Finally, we prove Condition (3) by showing $\cost(F)-\cost(F^\prime)<0$. We claim that the condition is equivalent to the condition that $b^Ty>0$ by showing $\cost(F)-\cost(F^\prime)=b^Ty$. Indeed, the coordinate of $b$ corresponding to the first type of constraint for simple path $Q$ connecting terminals $t,t^\prime$ has value $-\dist_G(t,t^\prime)$. Similarly, the coordinate of $b$ corresponding to the second type of constraint for canonical path $P$ connecting terminals $t,t^\prime$ has value $\dist_G(t,t^\prime)$. As a result, $$b^Ty=\sum_{P}y_P\cdot \dist_G(t,t^\prime)-\sum_{Q}y_Q\cdot\dist_G(t,t^\prime)=\cost(F)-\cost(F^\prime),$$
    where the first summation is over all canonical paths $P$, the second summation is over all simple paths $Q$, and $t,t^\prime$ are the terminals which paths $P$ and $Q$ connect. 




\bibliographystyle{alpha}
%\printbibliography
\bibliography{refs.bib}


\end{document}
